\chapter{Optics, Telescopes \& Detectors}

\section*{Theory problems}

\probhead{5.1}{}{2007}{Chiang Mai, Thailand}{TH 1.13}{2}
% source id: 2007_TH_1.13   tag: naked-eye diffraction limit, lunar crater
A crater on the surface of the Moon has a diameter of 80\,km. Is it possible to
resolve this crater with naked eyes, assuming the eye pupil aperture is 5\,mm?

\probhead{5.2}{}{2008}{Bandung, Indonesia}{TH S10}{20}
% source id: 2008_TH_S10   tag: CCD frame fit and position angle
The coordinates of the components of Visual Binary Star $\mu$ Sco on August 22,
2008 are given in the table below.

\begin{center}
\begin{tabular}{|l|c|c|}
\hline
 & $\alpha$ (RA) & $\delta$ (Dec) \\
\hline
$\mu$ Sco 1 (primary) & $20^{\mathrm{h}}\,17^{\mathrm{m}}\,38^{\mathrm{s}}\!.90$
  & $-12^\circ\,30'\,30''$ \\
\hline
$\mu$ Sco 2 (secondary) &
  $20^{\mathrm{h}}\,18^{\mathrm{m}}\,03^{\mathrm{s}}\!.30$
  & $-12^\circ\,32'\,41''$ \\
\hline
\end{tabular}
\end{center}

The stars are observed using Zeiss refractor telescope at the Bosscha
Observatory with aperture and focal length are 600\,mm and 10\,780\,mm,
respectively. The telescope is equipped with $765 \times 510$ pixels CCD
camera. The pixel size of the chip is
$9\,\mu\mathrm{m} \times 9\,\mu\mathrm{m}$.

\begin{description}
\item[a.] Can both components of the binary be inside the frame? (``YES'' or
  ``NO'', show it in your computation!)
\item[b.] What is the position angle of the secondary star, with respect to the
  North?
\end{description}

\probhead{5.3}{}{2009}{Tehran, Iran}{TH P7}{10}
% source id: 2009_TH_P7   tag: eyepiece FOV from star drift time
A student tries to measure field of view (FOV) of the eyepiece of his/her
telescope, using rotation of the Earth. To do this job, the observer points the
telescope towards Vega (alpha Lyr., RA: $18.5^{\mathrm{h}}$, Dec: $+39^\circ$),
turns off its ``clock drive'' and measures trace out time, $t = 5.3$ minutes,
that Vega crosses the full diameter of the FOV. What is the FOV of this
telescope in arc-minutes?

\probhead{5.4}{}{2010}{Beijing, China}{TH S8}{10}
% source id: 2010_TH_S8   tag: VLBI wavelength to resolve Sgr A*
The Galactic Center is believed to contain a super-massive black hole with a
mass $M = 4 \times 10^6\,M_\odot$. The astronomy community is trying to resolve
its event horizon, which is a challenging task. For a non-rotating black hole,
this is the Schwarzschild radius, $R_s = 3(M/M_\odot)$\,km. Assume that we have
an Earth-sized telescope (using Very Long Baseline Interferometry). What
wavelengths should we adopt in order to resolve the event horizon of the black
hole? The Sun is located at 8.5\,kpc from the Galactic Center.

\probhead{5.5}{}{2010}{Beijing, China}{TH S9}{10}
% source id: 2010_TH_S9   tag: photon count rate, Gemini 8-m
A star has a measured I-band magnitude of 22.0. How many photons per second are
detected from this star by the Gemini Telescope (8\,m diameter)? Assume that
the overall quantum efficiency is 40\% and the filter passband is flat.

\begin{center}
\begin{tabular}{lccc}
Filter & $\lambda_0(\mathrm{nm})$ & $\Delta\lambda(\mathrm{nm})$
  & $F_{VEGA}(\mathrm{W\,m^{-2}\,nm^{-1}})$ \\
$I$ & $8.00 \times 10^2$ & 24.0 & $8.30 \times 10^{-12}$ \\
\end{tabular}
\end{center}

\probhead{5.6}{}{2011}{Katowice, Poland}{TH S5}{10}
% source id: 2011_TH_S5   tag: radio dish matching optical resolution
What would be the diameter of a radiotelescope working at a wavelength of
$\lambda = 1$\,cm with the same resolution as an optical telescope of diameter
$D = 10$\,cm?

\probhead{5.7}{}{2012}{Rio de Janeiro, Brazil}{TH S12}{}
% source id: 2012_TH_S12   tag: telescope to resolve Earth's lunar wobble
What is the minimum diameter of a telescope, observing in the visible and near
ultraviolet bands, located in one of the Lagrangian points L4 or L5 of the
Sun--Earth system, in order to be able to detect the Earth's wobbling relative
to the ecliptic plane caused by the gravitational action of the Moon?

\probhead{5.8}{}{2013}{Volos, Greece}{TH S2}{10}
% source id: 2013_TH_S2   tag: telescope diameter to resolve hot Jupiter
Let us assume that we observe a hot Jupiter planet orbiting around a star at an
average distance $d = 5$\,AU. It has been found that the distance of this
system from us is $r = 250$\,pc. What is the minimum diameter, $D$, that a
telescope should have to be able to resolve the two objects (star and planet)?
We assume that the observation is done in the optical part of the
electromagnetic spectrum ($\lambda \sim 500$\,nm), outside the Earth's
atmosphere and that the telescope optics are perfect.

\probhead{5.9}{The Vega star in the mirror}{2014}{Suceava, Romania}{TH P12}{}
% source id: 2014_TH_P12   tag: double image photometry with mirror
Inside a photo camera a plane mirror is placed along the optical axis of the
lens of the objective (as seen in figure below). The length of the mirror is
half of the focal distance of the lens of the objective. The photo camera is
oriented as on the photographic plate situated in the focal plane of the photo
camera are captured two images with different illuminations of the Vega star.
Find out the difference between the apparent photographical magnitudes of the
two images of the Vega stars.

\ioaafig{2014_TH_P12-f1.pdf}
% Vega starlight entering a lens objective (aperture r), a plane mirror of
% length f/2 lying on the optical axis in front of the photographic plate at
% focal distance f, with the two image positions Sigma_1 and Sigma_2 separated
% by r/2 on the plate. No auto-extracted figure exists for 2014_TH_P12.

\probhead{5.10}{}{2015}{Magelang, Indonesia}{TH S4}{}
% source id: 2015_TH_S4   tag: resolving supermassive black hole
Let us assume that an observer using a hypothetical, far-infrared, Earth-sized
telescope (wavelength range 20 to 640\,$\mu$m) found a static and neutral
supermassive black hole with a mass of $2.1 \times 10^{10}\,\mathrm{M}_\odot$.
Determine the maximum distance at which this black hole can be resolved by the
observer.

\probhead{5.11}{AstroSat}{2016}{Bhubaneswar, India}{TH T12}{50}
% source id: 2016_TH_T12   tag: X-ray instrument count rates, SNR
India astronomy satellite, AstroSat, launched in September 2015, has five
different instruments.

\ioaafig{2016_TH_T12-f1.pdf}

In this question, we will discuss three of these instruments (SXT, LAXPC,
CZTI), which point in the same direction and observe in X-ray wavelengths. The
details of these instruments are given in the table below.

\begin{center}
\small
\begin{tabular}{|l|c|c|c|p{3.2cm}|c|}
\hline
Instrument & Band & Collecting & Effective Photon & Saturation level & No.\ of
\\
 & [keV] & Area [m$^2$] & Detection Efficiency & [counts] & Pixels \\
\hline
SXT & 0.3--80 & 0.067 & 60\% & 15000 (total) & $512 \times 512$ \\
\hline
LAXPC & 3--80 & 1.5 & 40\% & 50000 (in any one counter) or 200000 (total)
  & --- \\
\hline
CZTI & 10--150 & 0.09 & 50\% & --- & $4 \times 4096$ \\
\hline
\end{tabular}
\end{center}

You should note that LAXPC energy range is divided into 8 different energy band
counters of equal bandwidth with no overlap.

\begin{description}
\item[(T12.1)] Some X-ray sources like Cas A have a prominent emission line at
  0.01825\,nm corresponding to a radioactive transition of $^{44}$Ti. Suppose
  there exists a source which emits only one bright emission line corresponding
  to this transition. What should be the minimum relative velocity ($v$) of the
  source, which will make the observed peak of this line to get registered in a
  different energy band counter of LAXPC as compared to a source at rest?
\end{description}

These instruments were used to observe an X-ray source (assumed to be a point
source), whose energy spectrum followed the power law,
\[ F(E) = K E^{-2/3} \qquad
   [\text{in units of counts/keV/m}^2\text{/s}] \]
where $E$ is the energy in keV, $K$ is a constant and $F(E)$ is photon flux
density at that energy. Photon flux density, by definition, is given for per
unit collecting area (m$^2$) per unit bandwidth (keV) and per unit time
(seconds). From prior observations, we know that the source has a flux density
of 10\,counts/keV/m$^2$/s at 1\,keV, when measured by a detector with 100\%
photon detection efficiency. The ``counts'' here mean the number of photons
reported by the detector.

As the source flux follows the power law given above, we know that for a given
energy range from $E_1$ (lower energy) to $E_2$ (higher energy) the total
photon flux ($F_{\mathrm{T}}$) will be given by
\[ F_{\mathrm{T}} = 3K\left(E_2^{1/3} - E_1^{1/3}\right) \qquad
   [\text{in units of counts/m}^2\text{/s}] \]

\begin{description}
\item[(T12.2)] Estimate the incident flux density from the source at 1\,keV,
  5\,keV, 40\,keV and 100\,keV. Also estimate what will be the total count per
  unit bandwidth recorded by each of the instruments at these energies for an
  exposure time of 200 seconds.
\item[(T12.3)] For this source, calculate the maximum exposure time
  ($t_{\mathrm{S}}$), without suffering from saturation, for the CCD of SXT.
\item[(T12.4)] If the source became 3500 times brighter, calculate the expected
  counts per second in LAXPC counter 1, counter 8 as well as total counts
  across the entire energy range. If we observe for longer period, will the
  counter saturate due to any individual counter or due to the total count?
  Tick the appropriate box in the Summary Answersheet.
\item[(T12.5)] Assume that the counts reported by CZTI due to random
  fluctuations in electronics are about 0.00014 counts per pixel per keV per
  second at all energy levels. Any source is considered as ``detected'' when
  the SNR (signal to noise ratio) is at least 3. What is minimum exposure time,
  $t_c$, needed for the source above to be detected in CZTI?\\
  Note that the ``noise'' in a detector is equal to the square root of the
  counts due to random fluctuations.
\item[(T12.6)] Let us consider the situation where the source shows variability
  in number flux, so that the factor $K$ increases by 20\%. AstroSat observed
  this source for 1 second before the change and 1 second after this change in
  brightness. Calculate the counts measured by SXT, LAXPC and CZTI in both the
  observations. Which instrument is best suited to detect this change? Tick the
  appropriate box in the Summary Answersheet.
\end{description}

\probhead{5.12}{GMRT beam transit}{2016}{Bhubaneswar, India}{TH T5}{10}
% source id: 2016_TH_T5   tag: source transit through radio dish FWHM
Giant Metrewave Radio Telescope (GMRT), one of the world's largest radio
telescopes at metre wavelengths, is located in western India (latitude:
$19^\circ 6'$\,N, longitude: $74^\circ 3'$\,E). GMRT consists of 30 dish
antennas, each with a diameter of 45.0\,m. A single dish of GMRT was held fixed
with its axis pointing at a zenith angle of $39^\circ 42'$ along the northern
meridian such that a radio point source would pass along a diameter of the
beam, when it is transiting the meridian.

What is the duration $T_{\mathrm{transit}}$ for which this source would be
within the FWHM (full width at half maximum) of the beam of a single GMRT dish
observing at 200\,MHz?

\textbf{Hint:} The FWHM size of the beam of a radio dish operating at a given
frequency corresponds to the angular resolution of the dish. Assume uniform
illumination.

\probhead{5.13}{Telescope optics}{2016}{Bhubaneswar, India}{TH T7}{20}
% source id: 2016_TH_T7   tag: magnification, Barlow lens, CCD pixel scale
In a particular ideal refracting telescope of focal ratio $f/5$, the focal
length of the objective lens is 100\,cm and that of the eyepiece is 1\,cm.

\begin{description}
\item[(T7.1)] What is the angular magnification, $m_0$, of the telescope? What
  is the length of the telescope, $L_0$, i.e.\ the distance between its
  objective and eyepiece?
\end{description}

An introduction of a concave lens (Barlow lens) between the objective lens and
the prime focus is a common way to increase the magnification without a large
increase in the length of the telescope. A Barlow lens of focal length 1\,cm is
now introduced between the objective and the eyepiece to double the
magnification.

\begin{description}
\item[(T7.2)] At what distance, $d_{\mathrm{B}}$, from the prime focus must the
  Barlow lens be kept in order to obtain this desired double magnification?
\item[(T7.3)] What is the increase, $\Delta L$, in the length of the telescope?
\end{description}

A telescope is now constructed with the same objective lens and a CCD detector
placed at the prime focus (without any Barlow lens or eyepiece). The size of
each pixel of the CCD detector is 10\,$\mu$m.

\begin{description}
\item[(T7.4)] What will be the distance in pixels between the centroids of the
  images of the two stars, $n_{\mathrm{p}}$, on the CCD, if they are $20''$
  apart on the sky?
\end{description}

\probhead{5.14}{U-Band photometry}{2016}{Bhubaneswar, India}{TH T8}{20}
% source id: 2016_TH_T8   tag: photon counting, sky background
A star has an apparent magnitude $m_{\mathrm{U}} = 15.0$ in the $U$-band. The
$U$-band filter is ideal, i.e., it has perfect (100\%) transmission within the
band and is completely opaque (0\% transmission) outside the band. The filter
is centered at 360\,nm, and has a width of 80\,nm. It is assumed that the star
also has a flat energy spectrum with respect to frequency. The conversion
between magnitude, $m$, in any band and flux density, $f$, of a star in Jansky
($1\,\mathrm{Jy} = 1 \times 10^{-26}\,\mathrm{W\,Hz^{-1}\,m^{-2}}$) is given by
\[ f = 3631 \times 10^{-0.4 m}\,\mathrm{Jy} \]

\begin{description}
\item[(T8.1)] Approximately how many $U$-band photons, $N_0$, from this star
  will be incident normally on a $1\,\mathrm{m}^2$ area at the top of the
  Earth's atmosphere every second?
\end{description}

This star is being observed in the $U$-band using a ground based telescope,
whose primary mirror has a diameter of 2.0\,m. Atmospheric extinction in
$U$-band during the observation is 50\%. You may assume that the seeing is
diffraction limited. Average surface brightness of night sky in $U$-band was
measured to be 22.0\,mag/arcsec$^2$.

\begin{description}
\item[(T8.2)] What is the ratio, $R$, of number of photons received per second
  from the star to that received from the sky, when measured over a circular
  aperture of diameter $2''$?
\item[(T8.3)] In practice, only 20\% of $U$-band photons falling on the primary
  mirror are detected. How many photons, $N_{\mathrm{t}}$, from the star are
  detected per second?
\end{description}

\probhead{5.15}{GOTO}{2017}{Phuket, Thailand}{TH T10}{25}
% source id: 2017_TH_T10   tag: plate scale, CCD noise, SNR, exposure
The Gravitational-Wave Optical Transient Observer (GOTO) aims to carry out
searches of optical counterparts of any Gravitational Wave (GW) sources within
an hour of their detection by the LIGO and VIRGO experiments. The survey needs
to cover a big area on the sky in a short time to search all possible regions
constrained by the GW experiments before the optical burst signal, if any,
fades away. The GOTO telescope array is composed of 4 identical reflective
telescopes, each with 40-cm diameter aperture and f-ratio of 2.5, working
together to image large regions of the sky. For simplicity, we assume that the
telescopes' fields-of-view (FoV) do not overlap with one another.

\begin{description}
\item[a)] Calculate the projected angular size per mm at the focal plane, i.e.\
  plate scale, of each telescope. \hfill[6]
\item[b)] If the zero-point magnitude (i.e.\ the magnitude at which the count
  rate detected by the detector is 1 count per second) of the telescope system
  is 18.5\,mag, calculate the minimum time needed to reach 21\,mag at
  Signal-to-Noise Ratio (SNR) $= 5$ for a point source. We first assume that
  the noise is dominated by both the Read-Out Noise (RON) at 10 counts/pixel
  and the CCD dark (thermal) noise (DN) rate of 1 count/pix/minute. The CCDs
  used with the GOTO have a 6-micron pixel size and gain (conversion factor
  between photo-electron and data count) of 1. The typical seeing at the
  observatory site is around 1.0\,arcsec. \hfill[8]\\
  The Signal-to-Noise Ratio is defined as
  \[ \mathrm{SNR} \equiv
     \frac{\text{Total Source Count}}{\sqrt{\sum_i \mathrm{Noise}_i^2}}
     = \frac{\text{Total Source Count}}
            {\sqrt{\sigma_{\mathrm{RON}}^2 + \sigma_{\mathrm{DN}}^2 + \dots}},
  \]
  \[ \sigma_{\mathrm{RON}} = \sqrt{N_{\mathrm{pix}} \cdot \mathrm{RON}^2},
     \qquad
     \sigma_{\mathrm{DN}} = \sqrt{N_{\mathrm{pix}} \cdot \mathrm{DN} \cdot t},
  \]
  where $t$ is the exposure time.
\item[c)] Normally when the exposure time is long and the source count is high
  then Poisson noise from the source is also significant. Determine the
  relation between SNR and exposure time in the case that the noise is
  dominated by Poisson noise of the source. Recalculate the minimum exposure
  time required to reach 21\,mag with SNR $= 5$ from part b) if Poisson noise
  is also taken into consideration. The Poisson noise (standard deviation) of
  the source is given by
  $\sigma_{\mathrm{source}} = \sqrt{\text{Source Count}}$. In reality, there is
  also the sky background which can be important source of Poisson noise. For
  our purpose here, please ignore any sky background in the calculation.
  \hfill[6]
\item[d)] The typical localisation uncertainty of the GW detector is about 100
  square-degrees and we would like to cover the entire possible location of any
  candidate within an hour after the GW is detected. Estimate the minimum side
  length of the square CCD needed for each telescope in terms of the number of
  pixels. You may assume that the time taken for the CCD read-out and the
  pointing change are negligible. \hfill[5]
\end{description}

\probhead{5.16}{Radio Galaxy}{2018}{Beijing, China}{TH T9}{25}
% source id: 2018_TH_T9   tag: radiometer equation, integration time
An observer wants to use the Five-hundred-meter Aperture Spherical radio
Telescope (FAST) in China to observe a radio galaxy at redshift of $z = 0.06$.
We assume that the radio source is compact compared to the beam size of the
telescope at the observing frequencies, i.e., the source is point-like as seen
through the telescope. To detect a point source with FAST, it must be
sufficiently strong (bright) relative to the noise level (for single
polarization observations), $\sigma$, which depends on the bandwidth,
$\Delta\nu$, and the integration time (the radio astronomy equivalent of
exposure time), $t_i$, as follows:
\[ \sigma = \frac{2 k_B T_{sys}}{A_e \sqrt{t_i\,\Delta\nu}} \]
where $T_{sys}$ is the system temperature (about 150\,K in the frequency range
of 0.28\,GHz -- 0.56\,GHz and 25\,K in the frequency range of 1.05\,GHz --
1.45\,GHz), and $A_e = 4.6 \times 10^4\,\mathrm{m}^2$ is the effective area of
the telescope taking into account the total efficiency of the instrument.

This radio galaxy has an observed continuum flux density of
$f_\nu = 2.5 \times 10^{-3}$\,Jy at an observing frequency of 0.4\,GHz. The
bandwidth $\Delta\nu$ for the continuum observation centered at 0.4\,GHz is
$2.8 \times 10^8$\,Hz.

\begin{description}
\item[(a)] In order to detect the continuum flux density at 0.4\,GHz with a
  signal-to-noise ratio of 30 (a so-called 30$\sigma$ detection), what is the
  required integration time, $t_i$?
\item[(b)] We want to search for the neutral Hydrogen (HI) in the galaxy using
  21\,cm absorption line. The HI 21\,cm line, with rest frame frequency of
  1.4204\,GHz. Calculate the observed frequency ($\nu_{obs}$) of the HI line
  for this galaxy.
\item[(c)] The radio continuum emission from this galaxy can be described by a
  power law $f_\nu \sim \nu^\alpha$, with a spectral index of $\alpha = -0.2$.
  Calculate the continuum flux density at $\nu_{obs}$ for this galaxy.
\item[(d)] The line width of the HI 21\,cm absorption line is 90\,km/s.
  Calculate the line width in Hz at the observing frequency of $\nu_{obs}$.
  According to Figure 1, the HI 21\,cm line absorbs 4\% of the continuum flux
  density (on average) over the line width of 90\,km\,s$^{-1}$. In order to
  detect the absorption line at $\geq 3\sigma$ in three consecutive
  30\,km\,s$^{-1}$ channels, what is the required integration time?
\end{description}

% Figure 1: spectrum of the HI 21cm absorption relative to the continuum
% emission in the radio galaxy.
\ioaafig{2018_TH_T9-f1.pdf}

\probhead{5.17}{Photographing a nanosatellite}{2019}{Keszthely, Hungary}{TH 14}{60}
% source id: 2019_TH_14   tag: cubesat magnitude, CCD detection, trailing
The very first, entirely Hungarian-made, satellite was ``MASAT--1'', a
nanosatellite ``cubesat''. It was made mostly of aluminium with total mass
1\,kg, sides length $l = 10$\,cm and a longer communication aerial. It was
designed and prepared at the Technical University of Budapest (BME) in 2009 by
students. The launch occurred on February 13, 2012, using a Vega rocket from
the Kourou launch site --- together with several other cubesats from other
countries. It operated successfully right up to the last minutes of its
lifetime (the transmitted last data packages were captured by radio receivers
on the evening of January 9, 2015, a few hours later the nanosatellite
re-entered the atmosphere and disintegrated).

\ioaafig{2019_TH_14-f1.pdf}

The orbital altitude of MASAT--1 changed between $h_{\mathrm{min}} = 350$\,km
and $h_{\mathrm{max}} = 1450$\,km (due to its highly eccentric orbit), but in
this question assume it to be in a circular orbit 900\,km above the sea level.

The MASAT team wished to photograph their nanosatellite from the ground. Thus,
they called the staff of Baja Astronomical Observatory (South Hungary,
$\lambda_{\mathrm{B}} = 19.010843^\circ$,
$\varphi_{\mathrm{B}} = 46.180329^\circ$, $h_{\mathrm{B}} = 100$\,m) to
photograph the orbiting MASAT--1 with their telescope.

The observatory has a Ritchey-Chr\'etien type reflecting telescope with a
diameter of 50\,cm, and focal ratio of $f/8.4$. The CCD camera installed on the
telescope had a $4096 \times 4096$ chip with 9\,$\mu$m sized square pixels. The
quantum efficiency of the CCD was about 70\,\%. The practical detection limit
was about $19.5^{\mathrm{m}}$ visual band applying a
$\tau_{\mathrm{exp}} = 2$\,min exposure time. Disregard any effects of seeing.
The orbital inclination of MASAT--1 was about $i = 70^\circ$ and the direction
of the orbital motion was the same as the Earth rotation. In addition, let us
assume that the reflection area always is 100\,cm$^2$. Throughout this event,
the telescope was pointed to the local zenith, and its RA motor was tracking
the stars. Consider only light from the Sun, you can ignore light reflected
from Earth and Moon. Ignore the effects of atmospheric extinction.

\begin{description}
\item[a)] Calculate the apparent visual magnitude of this cubesat under ideal
  observational circumstances, i.e.\ when it was at the local zenith of the
  observing site (Baja, Hungary) at midnight. Omit all atmospheric effects, and
  consider the Earth to be a sphere. The albedo of a specular aluminium plate
  is $a \approx 0.70$. \hfill(10\,p)\\
  \textit{Hint:} Use a comparison of MASAT--1 with the Full Moon.
\item[b)] What was the Observatory's answer to the MASAT team; would it ever be
  possible to photograph their nanosatellite with the existing equipment at the
  observatory? \textbf{Mark ``YES'' or ``NO'' on the answer sheet.} Support
  your answer with detailed calculations. \hfill(42\,p)
\item[c)] What would have been the answer if they had taken into account the
  blurring effect of the atmosphere, the so-called seeing? In Hungary the
  typical FWHM (Full Width at Half Maximum) value of the blurred stellar images
  --- which can be approximated by a symmetric 2D Gaussian profile close to the
  zenith --- is about $3.5''$. \textbf{Mark ``YES'' or ``NO'' or ``Close to the
  limit'' on the answer sheet.} Support your answer with a short calculation.
  \hfill(8\,p)\\
  \textit{Hint:} Although the illumination of the seeing spot in the focal
  plane can be approximated by a symmetric 2D Gaussian profile, you can take it
  to be homogeneous in your calculation.
\end{description}

\probhead{5.18}{Improving a common reflecting telescope}{2019}{Keszthely, Hungary}{TH 4}{10}
% source id: 2019_TH_4   tag: mirror reflectivity, limiting-magnitude gain
A student has an average quality Cassegrain telescope, with primary and
secondary mirrors having $\varepsilon_1 = 91$\,\% reflectivity aluminium
layers.

\begin{description}
\item[a)] What will be the change in the limiting magnitude of this telescope
  by replacing the mirror coatings with ``premium'' quality
  $\varepsilon_2 = 98$\,\% reflectivity ones? \hfill(5\,p)
\item[b)] Assuming the student also uses a star diagonal mirror, also with
  reflectivity $\varepsilon_1$ with the original telescope --- what will be the
  improvement if he/she also replaces this piece with an
  $\varepsilon_3 = 99$\,\% reflectivity (``dielectric'' mirror) model, combined
  with the new $\varepsilon_2$ mirrors? \hfill(3\,p)\\
  (A star diagonal mirror is a flat mirror, inclined to the optical axis by
  $45^\circ$.)
\item[c)] Is this difference obviously detectable by the human eye? \textbf{
  Mark ``YES'' or ``NO'' on the answer sheet.} \hfill(2\,p)
\end{description}

Consider the whole visual band and disregard any wavelength dependence and
geometric effects.

\probhead{5.19}{Astrophotography}{2020}{GeCAA (Estonia, remote)}{TH 1}{10}
% source id: 2020_TH_1   tag: camera FOV, pixel scale, star trails
An astrophotographer, based on the Equator, uses a good digital camera on a
tripod, but with no tracking. The camera has a telescopic lens with focal
length of 150\,mm and aperture (objective diameter) of 84\,mm. The sensor has
effective light collecting diameter of 22.5\,mm. The photographic target is a
star field at the observer's Zenith.

\begin{description}
\item[(a)] Calculate the field of view (the angular width of the image) which
  can be captured on the sensor using this equipment. \hfill(2 points)
\item[(b)] The pixels in the camera's sensor are separated by a distance of
  $3.23\,\mu$m. What is the maximum possible exposure time for a single
  frame, so that no star trails appear on the exposed image?
  \hfill(5 points)
\item[(c)] For a better-quality image of the star field, the photographer
  decides to take multiple shots at the exposure time calculated in b) and
  then to stack them together. The total time for all these shots is 30\,s
  (ignore any time taken to write data to the memory card) What proportion of
  the total field of view is possible at the higher signal to noise ratio?
  \hfill(3 points)
\end{description}

\probhead{5.20}{Lucy: The First Mission to the Trojan Asteroids}{2021}{Bogot\'a (remote)}{TH Q13}{15}
% source id: 2021_TH_Q13   tag: CCD radiation noise, excited pixels
CCD cameras on space probes are very sensitive and exposed to space weather
conditions. Intense radiation passing through the CCD produces electron-hole
pairs in the silicon of the CCD chip. The rate at which these pairs are
produced is an important parameter when operating cameras on board spacecraft
and can be calculated for radiation of any given energy.

A high energy particle or photon of radiation passing through the CCD will
deposit some energy in the chip with each electron-hole pair it creates. The
`stopping power' of silicon for a given type of particle can be measured as
the energy per areal density ($\mathrm{areal\ density} = \mathrm{mass\ per\
unit\ area}$) that the silicon `takes away' from the travelling particle.

\ioaafig{2021_TH_Q13-f1.pdf}

NASA's Lucy mission will be the first to study the Trojan asteroids and will
revolutionize our understanding of the formation of the Solar System. One of
the instruments on board is L'LORRI (Lucy LOng Range Reconnaissance Imager),
which contains a sensitive CCD in order to produce detailed images of the
Trojan asteroids. Unfortunately, the radiation around Jupiter is very intense
and it can generate a lot of `noise' in the pixels of the CCD.

Let us assume that an average charged particle trapped in Jupiter's magnetic
field has an energy of 15\,MeV and that the flux of such particles in this
region is equivalent to about 600 electrons $\mathrm{s}^{-1}\,
\mathrm{cm}^{-2}$. Also assume that for each electron-hole pair which a
particle passing through a pixel creates, it deposits exactly the excitation
energy of the pair in that pixel. After the pixel crosses a threshold number
of electron-hole pairs it is `excited' and no more pairs can be produced in
that pixel. Any remaining energy in the particle is passed to the next pixel
(and so on).

Using the data given below for the CCD chip in the L'LORRI camera, answer the
following questions:

\begin{description}
\item[(13.1)] How many pixels will be excited by one such particle of
  radiation passing through the CCD when the spacecraft is near Jupiter's
  orbit? \hfill[10.0\,pt]
\item[(13.2)] Given the radiation flux near Jupiter, what percentage of the
  total number of pixels in an image will be excited? \hfill[5.0\,pt]
\end{description}

CCD Data:
\begin{itemize}
\item Exposure time of an image $= 30$\,ms
\item Pixels on the CCD $= 1024 \times 1024$
\item CCD Area $= 13\,\mathrm{mm} \times 13\,\mathrm{mm}$
\item CCD chip thickness $= 0.06$\,cm
\item Density of silicon, $\rho = 2.34\,\mathrm{g\,cm}^{-3}$
\item Excitation energy of single pair $= 2.36$\,eV
\item Excitation threshold of a single pixel $= 250$ pairs
\item `Stopping power' of silicon for a 15\,MeV electron
  $= 3.012\,\mathrm{MeV\,g}^{-1}\,\mathrm{cm}^{2}$
\end{itemize}

\probhead{5.21}{ALMA - Calculating photons}{2021}{Bogot\'a (remote)}{TH Q4}{10}
% source id: 2021_TH_Q4   tag: photon rates, interferometer resolution
ALMA is a radio observatory with a revolutionary design. It consists of 66
high-precision antennas, operating in the wavelength range from 0.32\,mm to
8.60\,mm. The principal array has fifty antennas of 12\,m diameter each that
can work together as a single telescope in the so-called interferometric
mode. There is also another array of four 12\,m antennas, and twelve smaller
antennas of 7\,m diameter each.

Imagine that a single 12\,m antenna is being calibrated, pointing to a source
with a known incident flux of $1 \times 10^{-20}\,\mathrm{W/m}^2$.

\begin{description}
\item[(4.1)] Assuming that all the flux arrives at the shortest wavelength of
  ALMA sensitivity, determine the average number of photons that would reach
  the detector every second. \hfill[2.0\,pt]
\item[(4.2)] Compare it to the average number of photons that would have
  reached the detector, if all the flux arrived at the longest wavelength of
  operation. \hfill[2.0\,pt]
\item[(4.3)] What is the angular resolution (in arcsec) of a single 12\,m
  antenna, operating at 74.9\,GHz? \hfill[2.0\,pt]
\item[(4.4)] Imagine the principal array operating at 74.9\,GHz in the
  interferometric mode. Assuming for simplicity that the spatial resolution
  is solely given by the longest baseline (largest distance between any pair
  of antennas), which turns to be $D_{\max} = 16$\,km, what would be the
  angular resolution (in arcsec) in this case? Treat this case as a single
  slit aperture instead of a circular one. \hfill[2.0\,pt]
\item[(4.5)] For a radio antenna, the term SEFD refers to `System Equivalent
  Flux Density', which is a characteristic energy flux density of the
  antenna, depending on its temperature and size. We also note that for
  energy estimation of radio photons, Rayleigh-Jeans approximation is valid.
  Assuming a system temperature of 691\,K, what would be the SEFD of the full
  ALMA observatory in Jansky if all the 66 antennas could work together?
  \hfill[2.0\,pt]
\end{description}

\probhead{5.22}{Microlensing}{2023}{Katowice, Poland}{TH Q3}{5}
% source id: 2023_TH_Q3   tag: telescope diameter for equal S/N
A faint subdwarf star ($I = 20.4$\,mag) in the Galactic bulge was observed to
brighten to $I' = 15.2$\,mag as a result of gravitational microlensing,
allowing a high-resolution spectrum to be obtained with the UVES spectrograph
on the Very Large Telescope (mirror diameter $D = 8.2$\,m).

Estimate the diameter of the telescope needed to obtain a spectrum of the
same quality with the same instrument and exposure time for this star at its
normal apparent brightness. The fiber aperture is small enough so that the
sky background is negligible.

\probhead{5.23}{Bolometer}{2023}{Katowice, Poland}{TH Q6}{13}
% source id: 2023_TH_Q6   tag: conical cavity absorption, effective temperature
The entrance cavity of a particular bolometer is a cone with an opening angle
of $30^\circ$, the surface of which has an energy absorption coefficient of
$a = 0.99$. Assume that there is no scattering of the incident radiation on
the walls of the cavity, only multiple mirror-like (specular) reflections.
The bolometer is connected to a cooler which keeps the bolometer cavity
surface at practically 0\,K temperature. The instrument is orbiting at 2\,au
from the Sun and is pointed directly at the centre of the Solar disk.

Calculate the temperature of a black body which would radiate the same
amount of energy as the entire bolometer opening does per unit surface area.

Note: the opening angle is defined as twice the angle between the axis of
the cone and its generatrix.

\probhead{5.24}{Cluster Photography}{2024}{Rio de Janeiro, Brazil}{TH T6}{20}
% source id: 2024_TH_T6   tag: plate scale, CCD signal-to-noise
An astronomer takes pictures, in the $V$-band, of a faint celestial target,
from a place with no light pollution. The selected target is the globular
cluster Palomar 4, which has an angular diameter of $\theta = 72.0''$ and a
uniform surface brightness in the $V$-band of $m_V = 20.6$\,mag/arcsec$^2$.
The observation equipment consists of one reflector telescope, with diameter
$D = 305$\,mm and f-ratio $f/5$, and a prime focus CCD with quantum
efficiency $\eta = 80\%$ and square pixels with size $\ell = 3.80\,\mu$m.

Given data:
\begin{itemize}
\item $V$-band central wavelength: $\lambda_V = 550$\,nm
\item $V$-band bandwidth: $\Delta\lambda_V = 88.0$\,nm
\item Photon flux for a 0-magnitude object in the $V$-band:
  $10000\,\mathrm{counts/nm/cm^2/s}$
\end{itemize}

\begin{description}
\item[(a)] Calculate the plate scale (the angle of sky projected per unit
  length of the sensor) of the observation equipment in arcmin/mm.
  \hfill(3 points)
\item[(b)] Estimate the number of pixels, $n_p$, covered by the cluster
  image on the CCD. \hfill(4 points)
\item[(c)] With an exposure time of $t = 15$ seconds, the astronomer obtains
  a signal-to-noise ratio of $S/N = 225$. Compute the brightness of the sky
  at the observation site, knowing that the CCD has a readout noise
  (standard deviation) of 5 counts per pixel and dark noise of 6 counts per
  pixel per minute. Give your answer in mag/arcsec$^2$. You may find useful:
  $\sigma^2_{RON} = n_p \cdot 1 \cdot RON^2$ and
  $\sigma^2_{DN} = n_p \cdot DN \cdot t$. \hfill(13 points)
\end{description}

\probhead{5.25}{Daksha Mission}{2025}{Mumbai, India}{TH T01}{10}
% source id: 2025_TH_T01   tag: detector geometry, timing localization
``Daksha'' is a proposed Indian mission consisting of two satellites S1 and
S2 orbiting the Earth in the same circular orbit of radius $r = 7000$\,km but
with $180^\circ$ phase difference. These satellites observe the universe in
high energy domain (X-rays and $\gamma$-rays). Each of the satellites of
Daksha uses several flat, rectangular detectors.

To understand how to localize a source in the sky, we shall use a simplified
model of the Daksha mission. Assume that S1 has only two identical detectors
D1 and D2, each of area $A = 0.50\,\mathrm{m}^2$, attached to an opaque mount
M as shown in the figure below. The detectors lie symmetrically around the
$y$ axis in planes perpendicular to the $x$-$y$ plane and make an angle
$\alpha = 120^\circ$ with each other.

\ioaafig{2025_TH_T01-f1.pdf}

\begin{description}
\item[(T01.1)] When observing a distant source located in the $x$-$y$ plane,
  detector D1 records a power $P_1 = 2.70 \times 10^{-10}\,\mathrm{J\,s}^{-1}$
  and detector D2 records a power
  $P_2 = 4.70 \times 10^{-10}\,\mathrm{J\,s}^{-1}$. Estimate the angle $\eta$
  made by the position vector of the source with the positive $y$-axis, with
  counter-clockwise angle from the positive $y$ axis being considered
  positive. \hfill[5]
\end{description}

Consider a single pulse from a distant source (not necessarily in the
$x$-$y$ plane) recorded by both satellites (S1 and S2) of Daksha. The times
of the peaks of the pulses recorded by S1 and S2 are $t_1$ and $t_2$,
respectively.

\begin{description}
\item[(T01.2)] If $t_1 - t_2$ was measured to be $10.0 \pm 0.1$\,ms then
  determine the fraction, $f$, of the celestial sphere where the source
  might lie. \hfill[5]
\end{description}

\probhead{5.26}{Atmospheric Seeing}{2025}{Mumbai, India}{TH T09}{35}
% source id: 2025_TH_T09   tag: refraction layers, speckles
A telescope with an achromatic convex objective lens of diameter
$D = 15$\,cm and focal length $f = 200$\,cm is pointed to a star at the
zenith.

\begin{description}
\item[(T09.1)] Find the diameter (in m), $d_{\mathrm{image}}$, of the image
  of a point source as produced by the objective lens at its focal plane for
  green light ($\lambda = 550$\,nm), considering only the effects of
  diffraction. \hfill[1]
\end{description}

The image of an astronomical source is also affected by the so-called
``atmospheric seeing''. The boundaries between the layers in the atmosphere
as well as the refractive indices of the layers change continuously due to
turbulence, temperature variation and other factors. This leads to tiny
changes in the position of the image in the focal plane of the telescope,
known as the ``twinkling effect''. For rest of the problem, apart from the
diffraction limited finite size of the image of the star discussed above, no
interference effects will be considered.

The left panel of the figure below shows a vertical cross-section of the
atmosphere with multiple layers of different refractive indices
($n_1, n_2, n_3, \ldots$). The right panel shows the zoomed in view of a
thin vertical segment of the atmosphere and the boundary between the two
lowest atmospheric layers of refractive indices $n_1$ and $n_2$
($n_1 > n_2$). We consider only these two layers and their boundary for this
problem. The diagrams are not to scale.

\ioaafig{2025_TH_T09-f1.pdf}

\begin{description}
\item[(T09.2)] Let the boundary between the two layers be at a height
  $H = 1$\,km directly above the telescope objective, with a tilt of
  $\theta = 30^\circ$ with respect to the horizontal plane. In all parts of
  this problem $\theta$ is taken to be positive in the anti-clockwise
  direction. For a monochromatic light source, $n_1 = 1.00027$ and
  $n_2 = 1.00026$. Let the angular shift in the position of the image at the
  focal plane of the telescope for a star at the zenith be $\alpha$.
  \begin{description}
  \item[(T09.2a)] Draw an appropriately labelled ray-diagram at the boundary
    showing $n_1$, $n_2$, $\theta$ and $\alpha$. \hfill[2]
  \item[(T09.2b)] Find the expression for $\alpha$ in terms of $\theta$,
    $n_1$ and $n_2$. Use the small angle approximations:
    $\sin\alpha \approx \alpha$ and $\cos\alpha \approx 1$. \hfill[2]
  \item[(T09.2c)] Calculate the displacement, $\Delta x_\theta$ (in m), in
    the position of the image if $\theta$ increases by 1\% (keeping $n_1$
    and $n_2$ fixed). \hfill[3]
  \item[(T09.2d)] Calculate the displacement, $\Delta x_n$ (in m), in the
    position of the image if $n_2$ increases by 0.0001\% (keeping $n_1$ and
    $\theta$ fixed). \hfill[3]
  \end{description}
\item[(T09.3)] For white light coming from a star at the zenith, choose
  which of the following most closely describes the shape and colour of the
  image by ticking the appropriate box (only one) in the Summary
  Answersheet. Note $x$ increases from left to right in the figure.
  \hfill[2]

  \begin{tabular}{clccc}
  & Image colour & Image shape & Left edge & Right edge \\
  A & White & Circular & & \\
  B & White & Elliptical & & \\
  C & Coloured & Circular & Blue & Red \\
  D & Coloured & Circular & Red & Blue \\
  E & Coloured & Elliptical & Blue & Red \\
  F & Coloured & Elliptical & Red & Blue \\
  \end{tabular}
\end{description}

For all remaining parts of this question we consider monochromatic green
light with $\lambda = 550$\,nm. We model the boundary between the layers as
a set of infinite zigzag planes (running perpendicular to the plane of the
page) separated by $d = 10$\,cm along $x$-axis, with either
$\theta = 10^\circ$ or $\theta = -10^\circ$. The figure below (not to scale)
shows a cross-section of this model of the atmosphere of width $W$
($W \ll H$). For telescopes with large aperture this zigzag nature of the
boundary results in formation of speckles in the focal plane.

\ioaafig{2025_TH_T09-f2.pdf}

\begin{description}
\item[(T09.4)] Consider an atmosphere modelled as above.
  \begin{description}
  \item[(T09.4a)] A section of the atmosphere with consecutive zigzag
    planes, with same parameters as stated above, is shown in the diagram
    below (not to scale).

    \ioaafig{2025_TH_T09-f3.pdf}

    In this diagram, reproduced in the Summary Answersheet, draw the paths
    of the incident light rays up to the plane where the telescope objective
    is placed, shown by the gray dotted line. Mark the region(s), if any, by
    ``X'' in the diagram where no light rays will reach. \hfill[4]
  \item[(T09.4b)] Calculate the width $W_X$ of such region(s). \hfill[3]
  \item[(T09.4c)] Find the largest diameter, $D_{\max}$, of the telescope
    objective with which it will be possible to obtain a single image of a
    star, by appropriately choosing the location of the telescope relative
    to the structure of the boundary. \hfill[4]
  \end{description}
\item[(T09.5)] Consider the case when the zigzag shape of the boundary is
  allowed in both $x$ and $y$ directions, (like a field of pyramids), and
  $D = 100$\,cm (with $f = 200$\,cm). Draw the qualitative pattern of the
  resulting speckles in the box given in the Summary Answersheet. \hfill[6]
\item[(T09.6)] For a turbulent atmosphere again consider the same parallelly
  running zigzag shape of the boundary layer only along $x$-direction, but
  now the angle between two planes are changing at a uniform rate from
  $10^\circ$ to $-10^\circ$ in 1.0\,s. Assume that this leads to a uniform
  rate of shift of the position of the image.

  Consider a telescope with $D = 8$\,cm and $f = 1$\,m. Estimate the longest
  exposure time $t_{\max}$ allowed for its CCD camera so that one gets only
  a single image, and any possible deviation in its position remains less
  than 1\% of the diffraction limited diameter of the image. \hfill[5]
\end{description}

\section*{Data-analysis problems}

\probhead{5.27}{CCD Image Processing}{2009}{Tehran, Iran}{DA D1}{60}
% source id: 2009_DA_D1   tag: aperture photometry, zero point, seeing
As an exercise of image processing, this problem involves use of a simple
calculator and tabular data (table 1.1) which contains the pixel values of an
image during the given exposure time. This image, which is a part of a larger
CCD image, was taken by a small CCD camera, installed on an amateur telescope
and using a $V$ band filter. Figure 1.1 shows this $50 \times 50$ pixels
image that contains 5 stars.

\ioaafig{2009_DA_D1-f1.pdf}

% TABLE ABRIDGED — table 1.1 (50x50 grid of pixel values) in crop edition

In table 1.1 the first row and column indicates the pixels' x and y
coordinates. Table 1.2 gives the telescope and the image specifications.

\begin{center}
Table 1.2\\[2pt]
\begin{tabular}{|l|c|}
\hline
Telescope focal length & $1.20$\,m \\
\hline
CCD pixel size & $25 \times 25\,\mu$m \\
\hline
Exposure time & $450$\,s \\
\hline
Telescope zenith angle & $25^\circ$ \\
\hline
Average extinction coefficient in $V$ band & 0.3\,mag/airmass \\
\hline
\end{tabular}
\end{center}

The observer identified stars 1, 3 and 4 by comparing this image with
standard star catalogues. Table 1.3 shows stars true magnitudes ($m_t$) as
given in the catalogue.

\begin{center}
Table 1.3\\[2pt]
\begin{tabular}{|c|c|}
\hline
Star & $m_t$ \\
\hline
1 & 9.03 \\
\hline
3 & 6.22 \\
\hline
4 & 8.02 \\
\hline
\end{tabular}
\end{center}

\begin{description}
\item[(a)] Using the available data, determine the instrumental magnitudes
  of the stars in the image. Assume the dark current is negligible and the
  image is flat fielded. For simplicity you can use a square aperture.

  Hint: The instrumental magnitude is calculated using the difference
  between the measured flux from the star in the aperture and the flux from
  an equivalent area of dark sky.
\item[(b)] The instrumental magnitude of a star in a CCD image is related to
  true magnitude as
  \[ m_I = m_t + KX - Zmag \]
  where $K$ is extinction, $X$ is airmass, $m_I$ and $m_t$ are respectively
  instrumental and true magnitude of star and $Zmag$ is zero point constant.
  Calculate the zero point constant ($Zmag$) for identified stars. Calculate
  average zero point constant ($\overline{Zmag}$).

  Hint: Zero point constant is the constant reducing extinction-free
  magnitudes to the true magnitude.
\item[(c)] Calculate true magnitudes of stars 2 and 5.
\item[(d)] Calculate CCD pixel scale for the CCD camera in units of arcsec.
\item[(e)] Calculate average brightness of dark sky in magnitude per square
  arcsec ($m_{sky}$).
\item[(f)] Use a suitable plot to estimate astronomical seeing in arcsec.
\end{description}

\probhead{5.28}{CCD image}{2010}{Beijing, China}{DA D1}{35}
% source id: 2010_DA_D1   tag: plate scale, chip size, moving objects
Information:

Picture 1 presents a negative image of sky taken by a CCD camera attached to
a telescope whose parameters are presented in the accompanying table (which
is part of the FITS datafile header). Picture 2 consists of two images: one
is an enlarged view of part of Picture 1 and the second is an enlarged image
of the same part of the sky taken some time earlier. Picture 3 presents a
sky map which includes the region shown in the CCD images.

The stars in the images are far away and should ideally be seen as point
sources. However, diffraction on the telescope aperture and the effects of
atmospheric turbulence (known as `seeing') blur the light from the stars.
The brighter the star, the more of the spread-out light is visible above the
level of the background sky.

Questions:

\begin{enumerate}
\item Identify any 5 bright stars (mark them by Roman numerals) from the
  image and mark them on both the image and map.
\item Mark the field of view of the camera on the map.
\item Use this information to obtain the physical dimensions of the CCD chip
  in mm.
\item Estimate the size of the blurring effect in arcseconds by examining
  the image of the star in Picture 2. (Note that due to changes in contrast
  necessary for printing, the diameter of the image appears to be about 3.5
  times the full width at half maximum (FWHM) of the profile of the star.)
\item Compare the result with theoretical size of the diffraction disc of
  the telescope.
\item Seeing of 1 arcsecond is often considered to indicate good conditions.
  Calculate the size of the star image in pixels if the atmospheric seeing
  was 1 arcsecond and compare it with the result from question 4.
\item Two objects observed moving relative to the background stars have been
  marked on Picture 1. The motion of one (``Object 1'') was fast enough that
  it left a clear trail on the image. The motion of the other (``Object 2'')
  is more easily seen on the enlarged image (Picture 2A) and another image
  taken some time later (Picture 2B).

  Using the results of the first section, determine the angular velocity on
  the sky of both objects. Choose which of the statements in the list below
  are correct, assuming that the objects are moving on circular orbits.
  (Points will be given for each correct answer marked and deducted for each
  incorrect answer marked.) The probable causes of the different angular
  velocities are:
  \begin{description}
  \item[a)] different masses of the objects,
  \item[b)] different distances of the objects from Earth,
  \item[c)] different orbital velocities of the objects,
  \item[d)] different projections of the objects' velocities,
  \item[e)] Object 1 orbits the Earth while Object 2 orbits the Sun.
  \end{description}
\end{enumerate}

Data: for Picture 1, the data are,

\begin{center}
\begin{tabular}{lll}
BITPIX & 16 & Number of bits per pixel \\
NAXIS & 2 & Number of axes \\
NAXIS1 & 1024 & Width of image (in pixels) \\
NAXIS2 & 1024 & Height of image (in pixels) \\
DATE-OBS & 2010-09-07 05:00:40.4 & Middle of exposure \\
TIMESYS & UT & Time Scale \\
EXPTIME & 300.00 & Exposure time (seconds) \\
OBJCTRA & 22 29 20.031 & RA of center of the image \\
OBJCTDEC & +07 20 00.793 & DEC of center of the image \\
FOCALLEN & 3.180m & Focal length of the telescope \\
TELESCOP & 0.61m & Telescope aperture \\
\end{tabular}
\end{center}

\ioaafig{2010_DA_D1-f1.pdf}

For Picture 2A (the same area observed some time earlier) the data are:
DATE-OBS = 2010-09-07 04:42:33.3 (middle of exposure). Picture 2B is an
enlargement of Picture 1 around Object 2.

\ioaafig{2010_DA_D1-f2.pdf}

\ioaafig{2010_DA_D1-f3.pdf}

\probhead{5.29}{}{2013}{Volos, Greece}{DA D1}{}
% source id: 2013_DA_D1   tag: camera focal length from plate scale
In Figure 1, part of the constellation of Ursa Major is shown. It was taken
with a digital camera with a large CCD chip
($17\,\mathrm{mm} \times 22\,\mathrm{mm}$). Find the focal length, $f$, of
the optical system and give the error of your results.
\ioaafig{2013_DA_D1-f1.pdf}

\probhead{5.30}{Photometric comparison of surveys}{2024}{Rio de Janeiro, Brazil}{DA D1}{75}
% source id: 2024_DA_D1   tag: SDSS/DES magnitude errors, calibration
You are an astronomer working with large photometric surveys, such as the
Sloan Digital Sky Survey (SDSS) and the Dark Energy Survey (DES), both of
which have your host, Observat\'orio Nacional, as a participant. SDSS used a
2.5\,m telescope in Apache Point, USA, during the 2000s, and DES used a 4\,m
telescope in Cerro Tololo, Chile, from 2013 to 2019. Even though they mostly
covered different hemispheres of the sky, they had an equatorial region in
common known as Stripe 82 that you can use to compare and calibrate the
photometry of different data sets, like SDSS and DES.

The following tables containing object positions and magnitudes from Stripe
82 were downloaded for analysis. However, due to a file system corruption on
the computer, the file names were scrambled, and now you cannot tell which
table belongs to which survey.

Tables 1 and 2 appear next to each other below, with an identification
number for each source, its equatorial coordinates, and its magnitude in the
$g$-band ($m_g$) with its error (err $m_g$).

\begin{description}
\item[(a)] From these tables, which survey (SDSS or DES) is Table 1 and
  which is Table 2? Assume that both surveys are equivalent regarding
  detector response, exposure times, and site characteristics.
  \hfill(5 points)
\item[(b)] Using the data in the table, plot the magnitude ($m_g$) on the
  $x$-axis (linear scale) and the error in magnitude (err $m_g$) on the
  $y$-axis (logarithmic scale) using the semi-log paper marked as Graph 1.
  Estimate the angular coefficient A (slope) and linear coefficient B
  ($y$-axis intercept) for each dataset. There is no need to calculate the
  associated errors. \hfill(35 points)
\item[(c)] The Signal to Noise ratio ($S/N$) is approximately the inverse of
  the error in the magnitude, $S/N \approx 1/(\mathrm{err}\ m_g)$. Using the
  linear fit calculated in the previous part, what is the $S/N$ reached for
  each survey at a magnitude of $m_g = 21.5$\,mag? \hfill(5 points)
\item[(d)] An object in Table 1 that is within 1 arcsecond of an object in
  Table 2 can be considered to be the same object. By looking at the RA and
  Dec of the objects in both tables, identify the objects in common and
  write down a new table with the matching IDs, $ID_1$ and $ID_2$.
  \hfill(15 points)
\item[(e)] Using the matched table from part (d), plot the $g$-band
  magnitude of each survey against the other, Table 1 on $x$-axis, and
  Table 2 on $y$-axis using the millimetre (linear) paper marked as Graph 2.
  Draw on error bars for each point in both horizontal and vertical
  directions, using values \textbf{double} err $m_g$ (known as a $2\sigma$
  uncertainty). From your graph, identify the stars that would be suitable
  for photometric calibration between the two surveys and write down their
  correspondings IDs from Table 1. \hfill(15 points)
  % sic: "correspondings IDs" as printed in the source paper.
\end{description}

\begin{center}
Table 1\\[2pt]
\begin{tabular}{ccccc}
$ID_1$ & RA & Dec & $m_g$ & err $m_g$ \\
 & (deg) & (deg) & (mag) & (mag) \\
\hline
1 & 0.047255 & 0.000406 & 21.7649 & 0.0120 \\
2 & 0.064741 & 0.021568 & 21.1111 & 0.0067 \\
3 & 0.064911 & 0.026395 & 21.3931 & 0.0084 \\
4 & 0.098343 & 0.054871 & 21.3934 & 0.0088 \\
5 & 0.022256 & 0.039129 & 21.9933 & 0.0157 \\
6 & 0.006188 & 0.066928 & 21.5490 & 0.0088 \\
7 & 0.083945 & 0.074259 & 21.9395 & 0.0126 \\
8 & 0.076715 & 0.079496 & 21.4808 & 0.0089 \\
9 & 0.057422 & 0.076006 & 21.8897 & 0.0127 \\
10 & 0.024412 & 0.087688 & 21.8341 & 0.0126 \\
11 & 0.044723 & 0.091782 & 21.8868 & 0.0172 \\
12 & 0.071089 & 0.091053 & 21.4390 & 0.0098 \\
\end{tabular}
\end{center}

\begin{center}
Table 2\\[2pt]
\begin{tabular}{ccccc}
$ID_2$ & RA & Dec & $m_g$ & err $m_g$ \\
 & (deg) & (deg) & (mag) & (mag) \\
\hline
1 & 0.006167 & 0.066874 & 21.9020 & 0.0576 \\
2 & 0.018660 & 0.007450 & 21.8039 & 0.0529 \\
3 & 0.047853 & 0.061487 & 21.3007 & 0.0418 \\
4 & 0.050870 & 0.015659 & 21.1678 & 0.0388 \\
5 & 0.051270 & 0.020812 & 21.2524 & 0.0401 \\
6 & 0.057414 & 0.075999 & 21.8884 & 0.0578 \\
7 & 0.064745 & 0.021583 & 21.3634 & 0.0422 \\
8 & 0.064910 & 0.026419 & 21.6428 & 0.0488 \\
9 & 0.071102 & 0.091058 & 21.9259 & 0.0751 \\
10 & 0.074946 & 0.002792 & 21.3258 & 0.0410 \\
11 & 0.076709 & 0.079474 & 21.5303 & 0.0476 \\
12 & 0.092635 & 0.077395 & 21.6995 & 0.0513 \\
13 & 0.098343 & 0.054854 & 21.6542 & 0.0499 \\
14 & 0.099332 & 0.093711 & 21.8802 & 0.0577 \\
\end{tabular}
\end{center}

\section*{Observation --- written and sky map}

\probhead{5.31}{Binoculars: Hyades limiting magnitude, M31 shape}{2007}{Chiang Mai, Thailand}{OBS-SKY 2}{10}
% source id: 2007_OBS-SKY_2   tag: binocular observation, limiting magnitude
Use the provided binoculars to observe the objects and then answer the
questions by writing or drawing directly in the question sheets.

\begin{description}
\item[(2.1)] The open star cluster ``Hyades'' in constellation Taurus is one
  of the nearest clusters to us, being only 151 light years away. From the
  provided chart with brightness of some stars indicated by the apparent
  magnitude in parentheses, please estimate the apparent magnitude of the
  star Gamma-Tauri ($\gamma$-Tauri) to the nearest first decimal digit.
  \hfill(5 points)

\ioaafig{2007_OBS-SKY_2-s1.pdf}

\item[(2.2)] Observe the Andromeda Galaxy (M31) then draw the approximate
  shape and size of the galaxy that you see through the binoculars in the
  frame below with correct orientation (in equatorial coordinates). The
  field of view of the binoculars is 6.8 degrees. \hfill(5 points)

\ioaafig{2007_OBS-SKY_2-s2.pdf}
\end{description}

\section*{Observation --- telescope}

\probhead{5.32}{CCD imaging and identification of a cluster field}{2008}{Bandung, Indonesia}{OBS-TEL 1}{5}
% source id: 2008_OBS-TEL_1   tag: telescope + CCD photography, object ID
Participants have to identify as many stars as possible in the field of a
celestial photograph using the telescope and CCD provided by the committee.
Participants have to choose one out of the five recommended regions in the
sky listed below. Then, point the telescope to the direction of the selected
sky region. Take three photographs with different exposure times and record
the images of the sky by the CCD camera. Save the observational data.
Transfer the data to the printing facilities to print out the result. Ask
the technical assistant for a help. Choose the best prints-out and use the
image to identify the stars in the field of observations. The following
procedures are,

\begin{enumerate}
\item Choose only one out of the following five recommended regions to the
  directions of (marked by the following bright clusters):
  \begin{description}
  \item[i.] M7 ($\alpha = 17^{\mathrm{h}}53.^{\mathrm{m}}3$,
    $\delta = -34^\circ 46.'2$)
  \item[ii.] M8 ($\alpha = 18^{\mathrm{h}}04.^{\mathrm{m}}2$,
    $\delta = -24^\circ 22.'0$)
  \item[iii.] M20 ($\alpha = 18^{\mathrm{h}}02.^{\mathrm{m}}4$,
    $\delta = -22^\circ 58.'7$)
  \item[iv.] M21 ($\alpha = 18^{\mathrm{h}}04.^{\mathrm{m}}2$,
    $\delta = -22^\circ 29.'5$)
  \item[v.] M23 ($\alpha = 17^{\mathrm{h}}57.^{\mathrm{m}}0$,
    $\delta = -18^\circ 59.'4$)
  \end{description}
  You may not change your choice.
\item Type in your country code, your student code, and your choice of
  region into the answer page provided in the computer.
\item Point the telescope to the chosen cluster by using telescope
  controller. If necessary you may move the telescope slightly to get the
  best position in the frame of the CCD by checking the display of CCDops
  software.
\item Display the region in The Sky map software provided in the computer,
  to confirm that the telescope is pointed to the selected object in the
  sky. You may change field of view of the sky chart.
\item You may invert the background images into white color, as in the chart
  mode. Copy and paste the sky chart from The Sky into the answer page. Use
  ``Ctrl c'' and ``Ctrl v'' buttons from the keyboard, respectively.
\item Type the equatorial coordinates of the centre of that object in the
  answer page as indicated in The Sky.
\item Take three photographs of the chosen object by using the attached CCD
  camera and CCDops software, with various exposure times. Choose exposure
  time in the range between 1 and 120 seconds. Image is automatically
  subtracted with dark frame with the same exposure time.
\item You must invert the background images into white color. Copy and paste
  images from CCDops into the answer page. Use ``Ctrl c'' and ``Ctrl v''
  buttons from the keyboard, respectively.
\item Save your answer page into hard disk in a Microsoft Word file format.
\item Print your answer page which consists of the photographs and the
  corresponding sky chart.
\item Go to the identification room and bring with you the prints-out of the
  sky chart and the photographs. Ask some help from technical assistant, if
  necessary.
\item Choose the best out of the three printed images and identify as many
  objects as possible on it.
\item Use the assigned computer and The Sky software to identify objects.
  Type your identification to your answer page.
\item Type on the answer page the names (or catalogue number), RA, Dec, and
  magnitude of each identified star and put the sequential number on the
  photograph. Make sure that you list the stars on the answer page following
  the same order and number as on the photograph.
\item Estimate the limiting magnitude of the photograph you choose,
  empirically.
\end{enumerate}

\probhead{5.33}{Telescope pointing and field identification}{2009}{Tehran, Iran}{OBS-TEL 3}{80}
% source id: 2009_OBS-TEL_3   tag: telescope pointing, field ID
Note: You have only 13 Minutes to answer all parts of this Question.

Before starting this part, please note: the telescope is pointed by the
examiner towards Caph ($\beta$ Cas). Please note the readings on the grade
circles before moving the telescope (to be used in 3.2).

\begin{description}
\item[(3.1)] Choose one of the 4 recommended stars listed below; write down
  the name of the selected star in table 3-1 and point the telescope to that
  star. Then, notify the examiner to check it. \hfill(40 Points)
  \begin{description}
  \item[1-] Deneb (Alpha Cygni)
  \item[2-] Alfirk (Beta Cephei)
  \item[3-] Algol (Beta Persei)
  \item[4-] Capella (Alpha Aurigae)
  \end{description}
  (You have 6 Minutes to answer 3.1)
\item[(3.2)] The Telescope was parked to Caph in Cassiopeia constellation
  (RA: $0^{\mathrm{h}}9.7^{\mathrm{m}}$; Dec: $59^\circ 12'$). Using the
  clock beside the telescope write down the local time (with the format of
  HH:MM:SS) in the appropriate field in Table B. Then, by using the graded
  circle on the telescope mount, estimate the ``declination'' and the ``hour
  angle'' of the target measured from South, which you chose in part one of
  this question. Then, complete Table 3-2. \hfill(40 Points)

  (You have 7 Minutes to answer 3.2)
\end{description}

\probhead{5.34}{Field of view of the telescope (NGC 6231)}{2012}{Rio de Janeiro, Brazil}{OBS-TEL 1}{}
% source id: 2012_OBS-TEL_1   tag: estimating eyepiece field of view
Estimate the field of view of this telescope, using 10\,mm Pl\"oss eyepiece
and star chart-1, showing nearby region of open cluster NGC 6231. Star chart
1 shows two angular distances. Use them as reference. Express your answer in
arc minutes and tenths of it.

\ioaafig{2012_OBS-TEL_1-s1.pdf}

\probhead{5.35}{Night observation round}{2014}{Suceava, Romania}{OBS-TEL 1}{}
% source id: 2014_OBS-TEL_1   tag: telescope observation tasks
For the observational test outdoor, a Newtonian telescope on equatorial
mount EQ5 ($D = 200$\,mm, $F = 1000$\,mm) is used. Note: the telescope is
already aligned, but not necessarily calibrated --- do not change the
position of the tripod! The observational round in the field should take a
maximum of 30 minutes.

\begin{enumerate}
\item Name any five constellations which will be at the meridian 2 hours
  from the start of your observation test?
\item Point the telescope to M39. When you finish, ask your assistant to
  verify. Write in the box the number which corresponds to the object type
  (1 - globular cluster, 2 - double cluster, 3 - open cluster, 4 - galaxy,
  5 - nebula).
\item The right ascension and the declination for $\beta$ Aql (Alshain) are
  $\alpha = 19^{\mathrm{h}}55^{\mathrm{m}}$ and $\delta = 6^\circ 26'$. By
  using the telescope find out the right ascension and the declination for
  $\delta$ Cep. Write down the values in the appropriate boxes.
\item Point the telescope to the coordinate
  $\alpha = 2^{\mathrm{h}}22^{\mathrm{m}}$ and $\delta = 57^\circ 10'$. When
  you finish, ask your assistant to verify. Write in the box the number
  which corresponds to the object type (1 - globular cluster, 2 - double
  cluster, 3 - open cluster, 4 - galaxy, 5 - nebula).
\item Estimate UT when the meridian, the ecliptic and the equator are
  intersecting at the same point in this night. You may use the telescope or
  any other method.
\item Estimate the galactic latitude of $\xi$ Dra (Grumium).
\item Estimate the ecliptic latitude of $\varepsilon$ Cyg (Gienah).
\end{enumerate}

\probhead{5.36}{Main telescope: point to M7 and observe}{2015}{Magelang, Indonesia}{OBS-TEL 1}{}
% source id: 2015_OBS-TEL_1   tag: telescope pointing and observation
Main Telescope (Total Duration: 20 minutes)

\begin{enumerate}
\item Point your telescope to NGC 6475 (M7)\\
  RA (J2000): $17^{\mathrm{h}}\,53^{\mathrm{m}}\,54^{\mathrm{s}}$\\
  Dec (J2000): $-34^\circ 49'$\\
  Show the object to your telescope assistant. If you fail to find the
  object in the first 10 minutes, ask the assistant to help, but you will
  lose the mark. \hfill(23 points)
\item Write down the cardinal directions (North, South, East and West) on
  Chart 1 in the answer sheet. \hfill(10 points)
\item Three stars of M7 are missing from the Chart 1. Mark the position of
  the missing stars with crosses (x). \hfill(21 points)
\item The apparent magnitudes of comparison stars A, B and C are 7.6, 7.2,
  and 5.6 respectively. Estimate and write down the apparent magnitudes of
  the missing stars beside respective cross marks. \hfill(21 points)
\item Estimate the FOV of the instrument, using the given stopwatch. Show
  your calculation in the answer sheet. \hfill(25 points)
\end{enumerate}

\ioaafig{2015_OBS-TEL_1-s1.pdf}

\probhead{5.37}{Telescope already set: identify the field}{2016}{Bhubaneswar, India}{OBS-TEL OT1}{10}
% source id: 2016_OBS-TEL_OT1   tag: telescope field identification
When you arrive at your observing station, DO NOT disturb the telescope
before attempting this question.

The telescope is already set to a deep sky object. Identify the object and
tick the correct box in the Summary Answersheet.

\emph{Note:} You can use any technique to identify the object. However, if
you disturb the telescope, you will NOT be helped to bring it back to the
original position.

\probhead{5.38}{Telescope observation task}{2016}{Bhubaneswar, India}{OBS-TEL OT2}{20}
% source id: 2016_OBS-TEL_OT2   tag: telescope observation
\begin{description}
\item[(OT2.1)] Point the telescope to M45. Show the object to the
  examiner. \hfill(5 points)

  \emph{Note:} 1. After 5 minutes, 1 mark will be deducted for a delay of
  every minute (or part thereof) in pointing the telescope.
  2. You have a single chance to be evaluated. If your pointing is incorrect
  the examiner will change the pointing to M45 for the next part of the
  question.
\item[(OT2.2)] Your Summary Answersheet shows the telescopic field of M45.
  In the image, the seven (7) brightest stars of the cluster are replaced by
  a `+' sign. Compare the image with the field you see in the telescope and
  number the `+' marks from 1 to 7 in the order of decreasing brightness
  (brightest is 1 and faintest is 7) of the corresponding stars. \hfill(15 points)
\end{description}

\textit{[M45 answersheet image not in available sources]}

\probhead{5.39}{Telescope observation task}{2016}{Bhubaneswar, India}{OBS-TEL OT3}{20}
% source id: 2016_OBS-TEL_OT3   tag: telescope observation
The examiner will give you a moon filter, an eyepiece with a cross-wire and
a stopwatch. Point the telescope towards the Moon. Attach the filter to the
telescope. On the surface of the Moon, you will see several ``seas'' (maria)
which are nearly circular in shape. Estimate the diameter of Mare
Serenitatis, $D_{\mathrm{MSr}}$, labelled as ``1'' in the figure below, as a
fraction of the lunar diameter, $D_{\mathrm{Moon}}$, by measuring the
telescope drift times, $t_{\mathrm{Moon}}$ and $t_{\mathrm{MSr}}$, for the
Moon and the mare, respectively.

\ioaafig{2016_OBS-TEL_OT3-s1.pdf}

\probhead{5.40}{Night observation with the real sky}{2017}{Phuket, Thailand}{OBS-TEL 1}{}
% source id: 2017_OBS-TEL_1   tag: telescope observation tasks
This part of the exam involves observation of the real sky. You must
complete two tasks using the equipment provided. You have 6 minutes to
complete the first task, and 4 minutes to complete the second task. Once you
complete a task, call out to the proctor to have it graded; once graded,
that answer is considered final and you may not return to it again. Once the
timer has expired any ungraded task will be graded as is.

\textbf{N1: Observation with Equatorial Mount Telescope}

Use the Equatorial Mount Telescope to observe the target given in the star
chart. Write down the Name and make sure the target is focused.

Included:
\begin{itemize}
\item Equatorial Mount Telescope
\item Mount is already polar aligned
\item Eyepiece: 25 mm
\item Finder scope already aligned with the main scope
\item Telescope starts already pointing at ``starting star'' (see map)
\item Star chart with starting and target position marked.
\end{itemize}

Target name: \underline{\hspace{5cm}}

\ioaafig{2017_OBS-TEL_1-s1.pdf}

\textbf{N2: Observation with Dobsonian Telescope}

Use the Dobsonian Telescope provided to observe only one of the following
objects and focus on the selected object.

\begin{center}
\begin{tabular}{cll}
 & Name & Bayer Designation \\
\hline
$\square$ & Menkar & $\alpha$ Cet \\
$\square$ & Markab & $\alpha$ Peg \\
$\square$ & Hamal & $\alpha$ Ari \\
$\square$ & Fomalhaut & $\alpha$ PsA \\
$\square$ & Alpheratz & $\alpha$ And \\
\end{tabular}
\end{center}

Included:
\begin{itemize}
\item Dobsonian Telescope
\item Eyepieces: 12 and 25 mm
\item Finder scope is already aligned with the main scope.
\end{itemize}

\probhead{5.41}{Eyepiece field of view}{2018}{Beijing, China}{OBS-TEL O5}{10}
% source id: 2018_OBS-TEL_O5   tag: field of view measurement
You will use a telescope, equipped with one eyepiece but no finderscope, to
observe five red LED screens at a distance.

The eyepiece's field of view is $45^\circ$. Please estimate the field of
view (with an accuracy of $0.1^\circ$) of the telescope when observing.

\probhead{5.42}{Text observed through the telescope}{2018}{Beijing, China}{OBS-TEL O6}{40}
% source id: 2018_OBS-TEL_O6   tag: telescope resolution/reading
You will use a telescope, equipped with one eyepiece but no finderscope, to
observe five red LED screens at a distance. Some of these screens display
the names of planets, some show Messier numbers, and some show equatorial
coordinates.

Write down the text you observed on the screen.

\probhead{5.43}{Telescope task 1}{2019}{Keszthely, Hungary}{OBS-TEL 1}{5}
% source id: 2019_OBS-TEL_1   tag: telescope observation
\emph{Telescope: 150/750 Newton. Eyepieces: 25 mm, 10 mm, Barlow lens:
2$\times$. The telescope is already polar aligned. Use the 25 mm eyepiece
for this task.}

\textbf{Finderscope alignment} \hfill available time: 5 minutes

The finderscope is NOT aligned at the beginning. Point the telescope to
Saturn and align the finderscope parallel to the main tube.

If the alignment of Saturn is not within the crosshair of the finderscope,
the telescope assistant will correct it --- and you receive only partial or
no points.

\medskip
\emph{Alternative task (used in case of bad sky conditions at low
altitudes):} The finderscope is NOT aligned at the beginning. Point the
telescope to Altair ($\alpha$ Aql) and align the finderscope parallel to the
main tube. If the alignment is not satisfactory, the telescope assistant
will correct it --- and you receive only partial or no points.

\probhead{5.44}{Telescope task 2}{2019}{Keszthely, Hungary}{OBS-TEL 2}{15}
% source id: 2019_OBS-TEL_2   tag: telescope observation
\emph{Telescope: 150/750 Newton. For this task, we recommend using the 10 mm
eyepiece and the 2$\times$ Barlow lens.}

\textbf{Observation of Saturn} \hfill available time: 10 minutes

\begin{itemize}
\item In the upper box, the circle represents the disk of Saturn and the
  horizontal line is the E--W direction on the sky. Pay attention to the
  direction of North (see top right corner).

  Mark the position of Titan by a cross.
\item The smaller box on the bottom right corner of the first box is for
  drawing the rings of Saturn. Again the circle represents the disk of
  Saturn.

  Draw the rings of Saturn in this box with the correct size and
  orientation. Both the outer and inner edges of the ring are necessary, no
  faint ring details or gaps are needed. Keep the orientation of the image
  the same as the orientation in the upper box.
\item Estimate the angular distance (in arcsec) and position angle (in
  degrees) of Titan relative to the center of Saturn. You may do your
  calculations besides the answer.
\end{itemize}

Apparent major axis of the ring: $43''$


\medskip
\emph{Alternative task (used in case of bad sky conditions at low
altitudes): \textbf{Epsilon Lyrae}}

\begin{itemize}
\item Find $\epsilon$ Lyr, and make a drawing of the field of view (with the
  object and other stars) with the 10 mm eyepiece. Label the directions
  North and East by two arrows and mark them as `N' and `E'.
\item Estimate the angular distance between the wide pair
  ($\epsilon_1$--$\epsilon_2$), and estimate the position angle of the same
  pair.
\item Increase the magnification with the 2$\times$ Barlow lens to be able
  to resolve and separate the two close pairs. Estimate the angle (in
  degrees to the nearest integer) subtended by the two close pairs relative
  to each other (the enclosed angle of the two lines going through the two
  narrow pairs). Do not give any PA, only the relative angle of the two
  close pairs. No drawing is needed.
\end{itemize}

If you cannot find $\epsilon$ Lyr, the assistant can point to it for you,
but only after 5 minutes. In this case you will lose the marks for pointing
the telescope to the object. The telescope assistant will check the object at
the end of the 10 min limit. If you are ready sooner, keep the star in the
FOV, and wait for the check.

\probhead{5.45}{Telescope task 3}{2019}{Keszthely, Hungary}{OBS-TEL 3}{10}
% source id: 2019_OBS-TEL_3   tag: telescope observation
\emph{Telescope: 150/750 Newton. Use the 25 mm eyepiece for this task.}

\textbf{M57 --- field of view} \hfill available time: 10 minutes

\begin{itemize}
\item Find the planetary nebula M57 (in constellation Lyra), and put it in
  the centre of the field of view in the main scope.
\item The star chart in the answer sheet shows a part of constellation Lyra.
  In this chart, draw the FOV circle around M57 as accurately as possible.
\end{itemize}

If you cannot find M57, the assistant will help you, but only after 5
minutes. In this case you will lose the marks for pointing to the object.

\ioaafig{2019_OBS-TEL_3-s1.pdf}

\probhead{5.46}{Telescope task 4}{2019}{Keszthely, Hungary}{OBS-TEL 4}{15}
% source id: 2019_OBS-TEL_4   tag: telescope observation
\emph{Telescope: 150/750 Newton. Use the 25 mm eyepiece for this task.}

\textbf{Variable star --- AF Cygni} \hfill available time: 15 minutes

\begin{itemize}
\item Use the given charts of the constellation Cygnus to find the variable
  star AF Cyg.

  The large scale finder chart has normal orientation (N is up, E is to the
  left). The smaller scale chart has `telescope' orientation (S is up, W is
  to the left). Brightnesses of reference stars are given without decimal
  points, e.g.\ `97' means 9.7 magnitude.

  If you do not find AF Cyg, the telescope assistant cannot help you to
  point to it in this task.
\item Estimate the magnitude of AF Cyg by comparing it with the reference
  stars and write it down, with decimal point, at one decimal accuracy
  (i.e.\ 9.7).

  Write the time of your observation in UTC. You may ask the telescope
  assistant for the time in the local time zone (CEST).
\end{itemize}

\textit{[AF Cyg finder charts not in source PDFs]}

\probhead{5.47}{Telescope task 5}{2019}{Keszthely, Hungary}{OBS-TEL 5}{5}
% source id: 2019_OBS-TEL_5   tag: telescope observation
\emph{You are not allowed to use the telescope for this task.}

\textbf{Naked eye brightness estimation} \hfill available time: 5 minutes

\begin{itemize}
\item Estimate the visual magnitude of the two naked-eye stars marked on the
  stellar chart of constellation Ursa Minor:
  \begin{description}
  \item[(a)] $\zeta$ UMi (zeta UMi = Alifa) --- STAR 2
  \item[(b)] $\gamma$ UMi (gamma UMi = Pherkad) --- STAR 1
  \item[(c)] Write your estimate with one decimal accuracy (e.g.\ 8.6).
  \end{description}
\item Estimate the angular distance between $\gamma$ UMi (STAR 1) and
  Polaris in degrees.
\end{itemize}

\textit{[UMi chart not in source PDFs]}

\probhead{5.48}{Night observation O1}{2022}{Kutaisi, Georgia}{OBS-TEL O1}{15}
% source id: 2022_OBS-TEL_O1   tag: telescope observation
\emph{Time available: 5 min}

Your task is to identify the circled objects in the images that you will see
through the telescopes (numbered 1--5). The images are posted in the windows
of the academic building and the telescopes are already pointed to each
image. A single object will be marked on each board. For each object in the
list below, if it was marked on one of the boards, write the corresponding
telescope number in the list. You will get 30 seconds per telescope.

\begin{center}
\begin{tabular}{lc|lc}
Object & Number & Object & Number \\
\hline
Jupiter & & $\beta$ Her / Antilicus & \\
Saturn & & $\alpha$ Oph / Rasalhague & \\
Mars & & $\epsilon$ Peg / Enif & \\
$\alpha$ And / Alpheratz & & $\alpha$ Per / Mirfak & \\
$\alpha$ Aql / Altair & & $\alpha$ Sco / Antares & \\
$\alpha$ Boo / Arcturus & & $\alpha$ Ser / Unukalhai & \\
$\alpha$ CrB / Alphecca & & $\epsilon$ Sgr / Kaus Australis & \\
$\beta$ Dra / Rastaban & & $\beta$ UMi / Kochab & \\
\end{tabular}
\end{center}

\probhead{5.49}{Night observation O2}{2022}{Kutaisi, Georgia}{OBS-TEL O2}{15}
% source id: 2022_OBS-TEL_O2   tag: telescope observation
\emph{Time available: 4 min}

Determine the field of view FOV of the given telescope using the 25 mm
eyepiece. Show your calculations for the method used.

\probhead{5.50}{Night observation O3}{2022}{Kutaisi, Georgia}{OBS-TEL O3}{6}
% source id: 2022_OBS-TEL_O3   tag: telescope observation
\emph{Time available: 2 min}

Using the result from problem O2, determine the focal length of the
telescope. Assume the apparent field of view of the 25 mm eyepiece is 45
degrees.

Show and explain your calculations.

\probhead{5.51}{Night observation O4}{2022}{Kutaisi, Georgia}{OBS-TEL O4}{14}
% source id: 2022_OBS-TEL_O4   tag: telescope observation
\emph{Time available: 4 min}

$\beta$ Cyg (Albireo) is a visual binary star. You will be shown a picture
of this binary system from one of the windows of the building nearby. Aim
the telescope at this picture, and estimate the acute angle between the
horizon and the line going through the system's stars.

\probhead{5.52}{Observation round: the alien saucer screens}{2023}{Katowice, Poland}{OBS-TEL 1}{50}
% source id: 2023_OBS-TEL_1   tag: telescope observation tasks
\textbf{General instructions.} Scientists have discovered a crashed alien
flying saucer. High up inside the hold, they found several screens
transmitting views of the sky and telescopes have been set up to let you see
them clearly from the level of the deck. Use your telescope to observe the
(simulated) targets on the screens and record your results.

There are 5 screens on the opposite side: the central one will display video
for tasks 1 and 2, the other four will display static images for tasks 3 and
4. The two screens closer to the centre will display the (same) image for
task 3, and the two outer screens will display the (same) image for task 4.
Point your telescope at the screens furthest away from you.

You will have a total of 30 minutes to complete the observing tasks, however
tasks 1 and 2 will only be displayed once: just as with real observations
you will only have one opportunity to collect the data. There will be two
clocks visible showing the time remaining in the round.

At the start of the round a clock on the central screen will show the
simulated time at the observer's location. The clock will have the correct
orientation when seen through the telescope. The time will be shown for 3
minutes after which it will disappear; use this to set a start time for your
observations.

Caution: the scale of the field of view is different between the video and
still images.

\medskip
\textbf{Observation 1: `Asteroid occultation'} (15 points)

Calculations based on the orbital elements predict that an asteroid will
occult the star HD 163390 for 21 s, with the maximum occultation (mid-time)
occurring at 23:03:32 UT. However, the ephemeris is not perfect and the
prediction may be wrong by up to 20 s for the time and by 10 s for the
duration.

Based on your observations, find the true mid-time and duration of the
occultation. To identify the star use Map 1 and the following coordinates:
\[ \mathrm{HD\ 163390:}\quad \mathrm{RA} = 17^h58^m05^s, \quad
   \mathrm{DEC} = -18^\circ 50' 46.14'' \]
The map and the sky are in the same epoch.

[Answer table: mid-time of occultation $\pm$ error; duration of occultation
$\pm$ error]

\ioaafig{2023_OBS-TEL_1-s1.pdf}

\medskip
\textbf{Observation 2: `Starlink'} (15 points)

In the same star field as for Question 1, a `train' of Starlink satellites
will appear near the meridian of $17^h59^m$ at around 23:05 UT. Their
passage will last for around three minutes.

You may assume that the centre of the star field is at an altitude of
$20^\circ$ and that the satellites are 400 km above the Earth's surface
moving on circular orbits with equal distances between them. You may also
assume that satellites will move vertically (perpendicular to the horizon).

\begin{description}
\item[(a)] Measure the angular velocity of the satellites as seen by an
  observer on the simulated sky.
\item[(b)] Measure the time interval between the passes of successive
  satellites and mark their path on the sky chart (Map 1).
\item[(c)] Calculate the theoretical angular velocity of the satellites as
  seen by the observer, using the information given in the question.
\item[(d)] Estimate the distance in km between two consecutive satellites.
\end{description}

Constants: $G = 6.674 \times 10^{-11}\,\mathrm{N\,m^2\,kg^{-2}}$;
$M_\oplus = 5.972 \times 10^{24}$\,kg; $R_\oplus = 6378$\,km.

\medskip
\textbf{Observation 3: `Planetary Moons'} (10 points)

The screen will display an image of one of the planets of the Solar System
as seen on August 15, 2023, at 00:00 UT. Identify any five moons and mark
them on the answer sheet (you may use the moon position chart attached below
and the table showing their brightness).

\ioaafig{2023_OBS-TEL_1-s2.pdf}

\begin{center}
\begin{tabular}{cll}
Number & Name & Magnitude \\
\hline
I & Mimas & 13.0 \\
II & Enceladus & 11.8 \\
III & Tethys & 10.4 \\
IV & Dione & 10.6 \\
V & Rhea & 9.9 \\
VI & Titan & 8.5 \\
VII & Hyperion & 14.4 \\
VIII & Japetus & 11.0 \\
\end{tabular}
\end{center}

[Answer sheet: image of the planet on which the positions of any 5 moons are
to be marked with a dot and labelled with their numbers (I, II, \ldots)]

\medskip
\textbf{Observation 4: `Supernova'} (10 points)

The other screen presents the view of a galaxy and a bright (mag $< 11$)
object which was not visible previously. Estimate the right ascension (RA)
and declination (DEC) coordinates of this star and estimate its magnitude.
You may use Map 2, with stellar coordinates and a list of magnitudes.

[Answer table: right ascension, declination, estimated magnitude]

\ioaafig{2023_OBS-TEL_1-s3.pdf}

\probhead{5.53}{Discovery of 'Z'}{2025}{Mumbai, India}{OBS-TEL OT01}{25}
% source id: 2025_OBS-TEL_OT01   tag: sky map + telescope object discovery
You are given a sky map ``Map-OT01'' along with the question paper. This map
shows only stars but no diffuse objects.

\ioaafig{2025_OBS-TEL_OT01-s1.pdf}

Four known stars (S1, S2, S3 and S4), which are present in the map above,
are given in the table below with their Common names, Bayer designations and
Equatorial coordinates.

\begin{center}
\begin{tabular}{cllll}
Sr. No. & Common Name & Bayer Name & RA & Dec \\
\hline
S1 & Alpheratz & $\alpha$ Andromedae & $00^h08^m24^s$ & $29^\circ 5'16''$ \\
S2 & Markab & $\alpha$ Pegasi & $23^h04^m46^s$ & $15^\circ 12'17''$ \\
S3 & Scheat & $\beta$ Pegasi & $23^h03^m47^s$ & $28^\circ 4'58''$ \\
S4 & Algenib & $\gamma$ Pegasi & $00^h13^m14^s$ & $15^\circ 10'59''$ \\
\end{tabular}
\end{center}

Complete the tasks (OT01.1) and (OT01.2) while you are planning the
observations.

\begin{description}
\item[(OT01.1)] Your first task is to mark these 4 stars (with a circle
  around each star) and label them as S1, S2, S3 and S4 on the sky map
  ``Map-OT01'' provided to you. \hfill[6]
\item[(OT01.2)] An astronomer has discovered a new diffuse object `Z' at the
  following coordinates ---
  RA: $21^h36^m10.6^s$, Dec: $-26^\circ 10'24.4''$.

  Mark the position of this diffuse object on the same sky map ``Map-OT01''
  with a $\oplus$ sign and label it as `Z'. Assume that a linear rectangular
  grid for equatorial coordinates is valid in the relevant region of the
  map. \hfill[7]
\end{description}

The following task is to be performed once you reach the telescope station.

On the screen diagonally opposite to your station, initially a welcome
message, followed by a sample sky (sky unrelated to the question) will be
displayed along with a countdown timer. You may use this time to orient the
telescope towards the screen and familiarize yourself with other equipment
provided at the station. At the end of this time a part of the sky given in
the map ``Map-OT01'' will be projected on the screen for the next 6 minutes.
Note that the scale of the projection shown on the screen is different from
the actual scale seen in the sky.

\begin{description}
\item[(OT01.3)] Find the new object `Z' with the telescope using any
  appropriate eyepiece. Then centre the object properly in the field of view
  of the eyepiece with the crosshair, and show it to the examiner at your
  station.

  At the end of 6 minutes, the projection will be blurred for 20 seconds. At
  this point you must step away from the telescope. The projection will be
  restored for the examiner to check the view through the telescope. This
  marks the end of the first task. \hfill[12]
\end{description}

\probhead{5.54}{Lensing Time Delay}{2025}{Mumbai, India}{OBS-TEL OT02}{25}
% source id: 2025_OBS-TEL_OT02   tag: telescope observation
Gravitational lensing can result in multiple images of a background source
if the source, the lensing object and the observer are nearly aligned. These
multiple images take different times to reach the observer, and if the
background source is variable, each image shows the same feature in its
variability after specific time delays. These time delay measurements are
extremely useful to estimate the current expansion rate of the Universe, the
Hubble constant.

We will consider the gravitational lens system shown in the above figure.
The left hand panel shows a galaxy cluster (lens) together with 4 images of
a background quasar formed due to gravitational lensing. The 4 images,
labelled A, B, C and D, have different fluxes as each image is magnified by
a different amount. For any given image the magnification does not change
with time. Light takes the largest time to travel for the image labelled D.

\ioaafig{2025_OBS-TEL_OT02-s1.pdf}

The light coming from this quasar is variable, and astronomers have
monitored this system for more than a decade now. The right hand panel of
the figure shows the light curve for the image D.

On the screen opposite to your station you will see a movie of the
gravitational lens system. This movie is 28 s long and will loop 6 times
with breaks of 1 minute or 2 minutes between runs. Every second on the watch
corresponds to 250.0 days in the actual lens system.

\begin{description}
\item[(OT02.1)] Let the time delays of image D with respect to images A, B
  and C be given as $t_{DA} = t_D - t_A$, $t_{DB} = t_D - t_B$, and
  $t_{DC} = t_D - t_C$, respectively. Find these time delays, taking any
  necessary steps to reduce the uncertainty in your results. \hfill[25]
\end{description}

\clearpage
\section*{Solutions to Chapter 5}

\solhead{5.1}{}
% source id: 2007_TH_1.13
\ioaafig{2007_TH_1.13-s1.pdf}

A and B are two distant objects. $A'$ and $B'$ are the central maxima of
their diffracted images. The angular position $\phi_1$ of the first minimum
relative to central maximum of each diffraction pattern is given by
$\phi_1 = 1.22\,\dfrac{\lambda}{D}$.

According to Lord Rayleigh the minimum angle of resolution is $\phi_1$.
Hence the images of A and B will be resolved if $\theta > \phi_1$,
$D > 1.22\,\dfrac{\lambda}{\theta}$.

We may take distance AB to be the diameter of the crater. Hence the
diameter $D$ of the eye's aperture must be, at least,
$D_{\min} = 1.22\,\dfrac{\lambda}{\theta}$.
\begin{align*}
\theta &= \frac{80\ \mathrm{km}}{\text{distance from Earth to Moon}}\\
&= \frac{80 \times 10^{3}\ \mathrm{m}}{3.844 \times 10^{8}\ \mathrm{m}}
 = 2.081 \times 10^{-4}\ \text{radian}
\end{align*}
\[ \lambda = 500 \times 10^{-9}\ \mathrm{m} \]
\begin{align*}
\therefore\ D_{\min}
&= \frac{1.22 \times 500 \times 10^{-9}}{2.081 \times 10^{-4}}\ \mathrm{m}
 = 2.93 \times 10^{-3}\ \mathrm{m}\\
&= 2.9\ \mathrm{mm}
\end{align*}
Hence, it is possible to resolve the 80\,km-diameter crater with naked eye.

\solhead{5.2}{}
% source id: 2008_TH_S10
\begin{description}
\item[a.] Chip size:
\begin{align*}
(9 \times 765)\ \mu\mathrm{m} \times (9 \times 510)\ \mu\mathrm{m}
&= 6885\ \mu\mathrm{m} \times 4590\ \mu\mathrm{m}\\
&\approx 6.9\ \mathrm{mm} \times 4.6\ \mathrm{mm}
\end{align*}
\hfill(10\%)

For Zeiss telescope: $FL = 10780$\,mm
\[
\text{Pixel size (arc sec)}
 = 206265 \times \text{pixel size}\ (\mu\mathrm{m})
   /(1000 \times \text{focal length (mm)})''
\]
\hfill(10\%)
\[
 = 206265 \times 9/(1000 \times 10780) = 0.1722''
\]
\hfill(10\%)
\begin{align*}
\text{FOV CCD chip} &= (765 \times 510) \times 0.1722
  = 131.7330'' \times 87.8220''\\
&= 2.1956' \times 1.4637'
\end{align*}
\hfill(10\%)

For $\mu$ Sco 1: $\alpha_1 = 20^{\mathrm{h}}\,17^{\mathrm{m}}\,38^{\mathrm{s}}.90
  = 20.2941^{\mathrm{h}} = 304.4121^\circ$,\quad
$\delta_1 = -12^\circ\,30'\,30'' = -12.5083^\circ$

For $\mu$ Sco 2: $\alpha_2 = 20^{\mathrm{h}}\,18^{\mathrm{m}}\,03^{\mathrm{s}}.30
  = 20.3009^{\mathrm{h}} = 304.5138^\circ$,\quad
$\delta_2 = -12^\circ\,32'\,41'' = -12.5447^\circ$

Let the angular separation between the stars be $\gamma$.
Use spherical trigonometry:
\[
\cos\gamma = \cos(90^\circ - \delta_1)\cos(90^\circ - \delta_2)
 + \sin(90^\circ - \delta_1)\sin(90^\circ - \delta_2)\cos(\alpha_2 - \alpha_1)
\]
\hfill(10\%)
\begin{align*}
\cos\gamma &= \cos(90^\circ - (-12^\circ.5083))\cos(90^\circ - (-12^\circ.5447))\\
&\quad + \sin(90^\circ - (-12^\circ.5083))\sin(90^\circ - (-12^\circ.5447))
   \cos(304.5138^\circ - 304.4121^\circ)\\
\cos\gamma &= \cos(102.5083^\circ)\cos(102.5447^\circ)\\
&\quad + \sin(102.5083^\circ)\sin(102.5447^\circ)\cos(0.1017^\circ)
 = 0.999998298
\end{align*}
\hfill(10\%)
\[
\gamma = 0.1057^\circ = 6.3424'
\]
\ioaafig{2008_TH_S10-s1.pdf}

NO, because $\gamma > \sqrt{2.1956^2 + 1.4637^2}$ (diagonal of FOV Chip).
\hfill(10\%)

\item[b.] Position angle ($S_1$):
\[
\frac{\sin(\alpha_2 - \alpha_1)}{\sin\gamma}
 = \frac{\sin S_1}{\sin(90^0 - \delta_2)}
\]
\[
\sin S_1 = \frac{\sin(0^\circ 6'\,6'')\,\sin(102^\circ 32'41'')}{\sin(6'.1020)}
 = 0.9388
% sic: the printed denominator is sin(6'.1020); numerically the quoted result
% 0.9388 corresponds to dividing by sin(gamma) = sin(6'.3424).
\]
\[
S_1 = 69.8535^\circ \text{ or } 110.1466^\circ
\]
\hfill(20\%)

From the figure above, we found that the correct position angle is
$110.15^\circ$ ($\pm 6^\circ$ is full mark). \hfill(10\%)
\end{description}

\solhead{5.3}{}
% source id: 2009_TH_P7
\ioaafig{2009_TH_P7-s1.pdf}

Figure shows that if FOV of telescope is $\beta$ then we have:
\[ \beta = \phi\cos\delta \]
\hfill(4 points)

As the earth rotate, Vega moves through the FOV with constant angular
velocity of the earth
\[
\omega = \frac{2\pi}{86164} = 7.29 \times 10^{-5}\ (\mathrm{rad/s})
\]
\hfill(2 points)
\[
\phi = \omega t = 7.29 \times 10^{-5} \times 5.3 \times 60 = 0.023\ \mathrm{(rad)}
\]
\[
FOV = \beta = \phi\cos\delta = 0.023\cos 39^\circ
 = 0.018\ \mathrm{(rad)} \simeq 62\ \mathrm{min}
\]
\hfill(4 points)

\solhead{5.4}{}
% source id: 2010_TH_S8
Observationally, the diameter of the Galactic black hole at the distance of
$L = 8.5\,\mathrm{kpc}$ has the angular size,
\[
\theta_{BH} = 2R_s/L
\]
\hfill(2)

On the other hand, an Earth-sized telescope ($D = 2R_e$) has the resolution,
\[
\theta_{tel} = 1.22\lambda/(2R_e)
\]
\hfill(2)

In order to resolve the black hole at Galactic center, we need to have
$\theta_{BH} \geq \theta_{tel}$, which marginally we consider
$\theta_{BH} = \theta_{tel}$.

This leads to,
\[
\lambda = 4R_e R_s/(1.22L)
\]
\hfill(4)

Taking the values, we have
\[
\lambda \approx 0.9\,\mathrm{mm}
\]
\hfill(2)

This means that we need to observe at least at near sub-mm frequencies, which
is in radio or far-infrared band.

\solhead{5.5}{}
% source id: 2010_TH_S9
The definition of the magnitude is:
\[
m_I = -2.5\lg F_I + const
\]
where $F_I$ is the flux received from the source. Using the data above, we can
obtain the constant:
\begin{align*}
0.0 &= -2.5\lg(0.83 \times 10^{-11}) + const\\
% sic: the source prints 0.83 x 10^{11} (exponent sign lost in print); the
% value meant is 0.83 x 10^{-11} = 8.3 x 10^{-12}, as const = -27.7 confirms.
const &= -27.7
\end{align*}
Thus,
\begin{align*}
m_I &= -2.5\lg F_I - 27.7\\
F_I &= 10^{\frac{m_I + 27.7}{-2.5}}
  = 1.3 \times 10^{-20}\,\mathrm{W\,m^{-2}\,nm^{-1}}
\end{align*}
\hfill(4)

For our star, at an effective wavelength $\lambda_0 = 800$\,nm, using this
flux, the number of photons detected per unit wavelength per unit area is the
flux divided by the energy of a photon with the effective wavelength:
\[
N_I = \frac{1.3 \times 10^{-20}}{hc/\lambda_0}
 = 5.3 \times 10^{-2}\ \mathrm{photons\,s^{-1}\,m^{-2}\,nm^{-1}}
\]
\hfill(3)

Thus the total number of photons detected from the star per second by the 8\,m
Gemini telescope over the $I$ band is
\begin{align*}
N_I(total) &= (\text{tel.\ collecting area}) \times QE
  \times Bandwidth \times N_I\\
&= (\pi \times 4^2) \times 0.4 \times 24 \times N_I\\
&= 26\ \mathrm{photons/s} \approx 30\ \mathrm{photons/s}
\end{align*}
\hfill(3)

\solhead{5.6}{}
% source id: 2011_TH_S5
Angular resolution: $\sin(\lambda) = 1.22\,\lambda/D$
% sic: the source writes sin(lambda) for the sine of the resolution angle
\hfill[2]
\[
D_{\mathrm{radio}}\,[\mathrm{m}]
 = (\lambda_{\mathrm{radio}} \cdot D_{\mathrm{optic}})/\lambda_{\mathrm{optic}}
 = (1\,\mathrm{cm} \cdot 10\,\mathrm{cm})/550\,\mathrm{nm}
 = (0.01\ \mathrm{m} \cdot 0.1\ \mathrm{m})/0.00000055\ \mathrm{m}
 = 1820\ \mathrm{m}
\]

equation of angular resolution as a function of $\lambda$ and $D$ \hfill[2]

proportion of
$\lambda_{\mathrm{radio}}/D_{\mathrm{radio}}
 = \lambda_{\mathrm{optic}}/D_{\mathrm{optic}}$ \hfill[2]

proper calculation \hfill[2]

and result assuming $\lambda_{\mathrm{optic}}$ in proper range:
500--600\,nm; 400--900\,nm; out of range \hfill[2]

\solhead{5.7}{}
% source id: 2012_TH_S12
The phenomenon is as follow:
\ioaafig[1]{2012_TH_S12-s1.pdf}

The Center of Mass of the Earth--Moon system (CM) lies on the plane of the
ecliptic, so as the Lagrangian points of its orbit around the Sun. Let be $D$
the distance between Earth and Moon and $x$ the distance between Earth and
the CM. We then have:
\[
M_\oplus \cdot x = M_{\mathrm{Moon}} \cdot (D - x)
\ \therefore\
x = \left(\frac{M_{\mathrm{Moon}}}{M_\oplus + M_{\mathrm{Moon}}}\right)D
\]
\hfill(2 pt)

The Earth's displacement perpendicular to the ecliptic is $A = 2x\sin\theta$,
where $\theta = 5^\circ$ is the inclination of the plane of Moon's orbit
relative to the plane of ecliptic. Calculating:
\[
A = \frac{2D\,M_{\mathrm{Moon}}}{M_\oplus + M_{\mathrm{Moon}}}\sin 5^\circ
 = 810\ \mathrm{km}
\]
\hfill(4 pt)

The telescope resolution is given by
\[
r'' = 1.22\frac{\lambda}{\delta} = \frac{A}{d}
\]
\hfill(3 pt)

where $\delta$ is the diameter of the objective lens and $d$ is the distance
to the telescope to Earth. But, at L4(L5), Earth and Sun forms an equilateral
triangle, so $d = 1$\,UA. Using $\lambda = 400$\,nm (that is, the visible
wavelength that allows the smallest telescope), we find
$\delta = 9.01\ \mathrm{cm}$.
\hfill(1 pt)

\solhead{5.8}{}
% source id: 2013_TH_S2
\textit{Answer:} From the figure we have:
\ioaafig{2013_TH_S2-s1.pdf}
\[
\omega(\mathrm{rad}) \simeq \frac{d}{r} = \frac{5\,\mathrm{AU}}{250\,\mathrm{pc}}
 = \frac{5 \times 1.5 \times 10^{11}\,\mathrm{m}}
        {250 \times 3.09 \times 10^{16}\,\mathrm{m}}
 = 9.70 \times 10^{-8}\,\mathrm{rad}
\]
\hfill(5 Points)

Let us assume that $D$ is the minimum diameter of our space telescope.
Its angular resolution is
\[
\Delta\theta = 1.22 \times \frac{\lambda}{D}
\quad\rightarrow\quad
9.70 \times 10^{-8}\,\mathrm{rad}
 = 1.22 \times \frac{500 \times 10^{-9}\,\mathrm{m}}{D}
\quad\rightarrow\quad
D = 6\ \mathrm{m}
\]
\hfill(5 Points)

\solhead{5.9}{The Vega star in the mirror}
% source id: 2014_TH_P12
The light beam arriving from \textbf{Vega Star} can be considered paraxial,
due to the distance from it to the observer on Earth. The explanation for the
existence of two distinct images of the star is that the optical axis of the
objective is not parallel with the light beam from the star.

The images on the camera plate are symmetrical placed relative to the
principal optical axis.

\ioaafig{2014_TH_P12-s1.pdf}

Each of the point images of the Vega Star $\Sigma_1$ and $\Sigma_2$ didn't
concentrate the same light fluxes. In the down below figure it can be seen
the sections of the lens which correspond to each image. The sector APBC is
passed by the light which concentrates in the image $\Sigma_2$ and the light
passing the sector ACBQ concentrates into the point image $\Sigma_1$. See the
picture in figure 13.

\ioaafig{2014_TH_P12-s2.pdf}

The ratio between the light fluxes concentrated into the two image points
will directly depend on the ratio of the two sectors areas.

From the geometry of the figure 2 results:
\begin{align*}
&MN = OM;\quad N\Sigma_1 = OC = \frac{r}{2};\\
&\angle(CBO) = 30^\circ;\quad \angle(BOC) = 60^\circ;\quad
 \angle(AOB) = 120^\circ;
\end{align*}
\[
\frac{S_1}{S_2} = \frac{8\pi + 3\sqrt{3}}{4\pi - 3\sqrt{3}} \approx 4.
\]

Using the Pogson formula:
\[
\log\frac{E_1}{E_2}
 = \log\frac{\dfrac{\sigma T_{\mathrm{V}}^4 \cdot 4\pi R_{\mathrm{V}}^2}
                   {4\pi d_{\mathrm{PV}}^2} \cdot S_1}
            {\dfrac{\sigma T_{\mathrm{V}}^4 \cdot 4\pi R_{\mathrm{V}}^2}
                   {4\pi d_{\mathrm{PV}}^2} \cdot S_2}
 = -0.4(m_1 - m_2);
\]
\[
\log\frac{S_1}{S_2} = -0.4(m_1 - m_2);
\]
\[
m_2 - m_1 = 1.5^{\mathrm{m}}.
\]

\solhead{5.10}{}
% source id: 2015_TH_S4
This problem can be solved by using relationship between the telescope's
resolution and the angular size of black holes which can be seen in this
equation:
\[
\phi_{\mathrm{telescope}} \leq \phi_{\mathrm{blackhole}}
\quad\Rightarrow\quad
1.2\frac{\lambda}{D_{\mathrm{telescope}}}
 \leq \frac{\text{Diameter of event horizon}}
           {\text{Observer distance to the Galactic center}}
\]
\[
1.2\frac{\lambda}{D_{\mathrm{telescope}}}
 \leq \frac{2R_{\mathrm{blackhole}}}{d}
\]
\hfill(25)

Then, it is assumed that the radius of black hole is the Schwarzschild
radius; and the telescope is also Earth-sized telescope. Therefore:
\[
1.2\frac{\lambda}{2R_{\mathrm{Earth}}}
 \leq \frac{2R_{\mathrm{Schwarzschild,BH}}}{d}
\]
\begin{align*}
d &\leq \frac{4 \times R_{\mathrm{Schwarzschild,BH}} \times R_{\mathrm{Earth}}}
             {1.2 \times \lambda}
\quad\Rightarrow\quad
d \leq \frac{4 \times M_{\mathrm{Black\,Hole}}
             \times R_{\mathrm{Schwarzschild,Sun}} \times R_{\mathrm{Earth}}}
            {1.2 \times M_{\mathrm{Sun}} \times \lambda}\\
d &\leq \frac{4 \times 2.1 \times 10^{10}\,M_{\mathrm{Sun}}
             \times R_{\mathrm{Schwarzschild,Sun}} \times R_{\mathrm{Earth}}}
            {1.2 \times M_{\mathrm{Sun}} \times \lambda}
\quad\Rightarrow\quad
d \leq \frac{8.4 \times 10^{10}
             \times R_{\mathrm{Schwarzschild,Sun}} \times R_{\mathrm{Earth}}}
            {1.2 \times \lambda}
\end{align*}
\hfill(50)

To find the maximum distance, then we should use the minimum $\lambda$.
Therefore:
\[
d_{\max} = \frac{8.4 \times 10^{10}
            \times R_{\mathrm{Schwarzschild,Sun}} \times R_{\mathrm{Earth}}}
           {1.2 \times \lambda_{\min}}
 = \frac{8.4 \times 10^{10} \times 2.95 \times 10^{3}
         \times 6.3708 \times 10^{6}}
        {1.2 \times 20 \times 10^{-6}}\ \mathrm{m}
 = 6.6 \times 10^{25}\ \mathrm{m}
\]
\hfill(25)

\solhead{5.11}{AstroSat}
% source id: 2016_TH_T12
\textit{No official solution for the AstroSat problem was ever published: both the archived and the official 2016 solutions files contain the solution to a different (replaced) problem --- ``Exoplanets'' --- in place of T12. The AstroSat problem appears only in the exam paper.}

\solhead{5.12}{GMRT beam transit}
% source id: 2016_TH_T5
As the dish is pointed towards northern meridian at zenith angle of
$39.7^\circ$, altitude of the centre of the beam is
\[
a = 90.00^\circ - z = 90.00^\circ - 39.70^\circ = 50.30^\circ \tag*{(1.0)}
\]
Thus, declination of the source should be,
\[
\delta = 90.00 - a + \phi
 = 90.00^\circ - 50.30^\circ + 19.10^\circ = 58.80^\circ \tag*{(2.0)}
\]
FWHM beam size (for uniform illumination) will be given by
\begin{align*}
\theta &= \frac{1.22\lambda}{D} \tag*{(1.5)}\\
&= \frac{1.22c}{D\nu}
 = \frac{1.22 \times 2.998 \times 10^{8}}{45.0 \times 2 \times 10^{8}}\\
&= 0.0406\ \mathrm{rad}\\
\theta &= 2.33^\circ \tag*{(1.5)}
\end{align*}
\begin{align*}
T_{\mathrm{transit}} &= \frac{\theta \times 3.99\ \mathrm{min}}{\cos\delta}
  \tag*{(3.0)}\\
&= \frac{2.33 \times 3.99\ \mathrm{min}}{\cos 58.8}
\end{align*}
\[
\boxed{T_{\mathrm{transit}} = 17.9\ \mathrm{min}} \tag*{(1.0)}
\]

\solhead{5.13}{Telescope optics}
% source id: 2016_TH_T7
\begin{description}
\item[(T7.1)] The magnification will be given by,
\begin{align*}
m_0 &= \frac{f_{\mathrm{o}}}{f_{\mathrm{e}}} \tag*{(1.0)}\\
&= \frac{100}{1} = 100 \tag*{(1.0)}
\end{align*}
The magnification is $m_0 = 100$.

Length of the telescope will be
\begin{align*}
L_0 &= f_{\mathrm{o}} + f_{\mathrm{e}} \tag*{(1.0)}\\
&= 100 + 1 = 101\ \mathrm{cm} \tag*{(1.0)}
\end{align*}
The telescope length will be $L_0 = 101$\,cm.

\item[(T7.2)] We use the following sign convention. Lens is the origin.
Direction along the direction of light is taken as positive. The lens formula
is $\dfrac{1}{f} = \dfrac{1}{v} - \dfrac{1}{u}$ ($f$ is positive for convex
lens and negative for concave lens). Magnification is $m = \dfrac{v}{u}$.
Solutions using other sign conventions are acceptable.

Let $v$ be image distance from the Barlow lens.
\[
\frac{1}{f_{\mathrm{B}}} = \frac{1}{v} - \frac{1}{u} \tag*{(1.0)}
\]
Distance of Barlow lens, $d_{\mathrm{B}}$, before the prime focus is same as
the object distance, $u$, in this case. \hfill(1.0)
\[
\frac{1}{f_{\mathrm{B}}} = \frac{1}{v} - \frac{1}{d_B}
\]
Also, $m_{\mathrm{B}} = 2 = \dfrac{v}{u} = \dfrac{v}{d_{\mathrm{B}}}$
\hfill(0.5)
\begin{align*}
\therefore\ \frac{1}{d_{\mathrm{B}}} &= \frac{2}{v} \tag*{(1.0)}\\
\therefore\ \frac{1}{-1} &= \frac{1}{v} - \frac{2}{v}\\
-1 &= \frac{-1}{v}\\
v &= 1\ \mathrm{cm}\\
d_{\mathrm{B}} &= \frac{v}{2} = \frac{1\ \mathrm{cm}}{2} = 0.5\ \mathrm{cm}
  \tag*{(2.5)}
\end{align*}
The positive sign for $d_{\mathrm{B}}$ indicates that the Barlow lens was
introduced $\boxed{0.5\ \mathrm{cm}}$ before the prime focus.

\item[(T7.3)] The increase in the length will be,
\begin{align*}
\Delta L &= v - d_{\mathrm{B}} \tag*{(2.0)}\\
&= 1.0 - 0.5 = 0.5\ \mathrm{cm} \tag*{(2.0)}
\end{align*}
Thus, the length will be increased by $\boxed{\Delta L = 0.5\ \mathrm{cm}}$.

\item[(T7.4)] Plate scale at prime focus is given by,
\[
s = \frac{1}{f_o} = \frac{1\ \mathrm{rad}}{1\ \mathrm{m}}
 = 0.206\,265\ \mathrm{arcsec}/\mu\mathrm{m} \tag*{(2.0)}
\]
Since each pixel is $10\ \mu$m in size,
\[
s_{\mathrm{p}} = 10 \times 0.206\ \mathrm{arcsec}/\mu\mathrm{m}
 = 2.06\ \mathrm{arcsec/pixel} \tag*{(2.0)}
\]
Two stars will be separated by,
\[
n_{\mathrm{p}} = \frac{20''}{2.06''}\ \mathrm{pixels}
 \simeq 10\ \mathrm{pixels} \tag*{(2.0)}
\]
\end{description}

\solhead{5.14}{U-Band photometry}
% source id: 2016_TH_T8
\begin{description}
\item[(T8.1)]
The $U$-band is defined as $(360 \pm 40)$\,nm. Thus, the maximum, minimum and
average frequencies of the band will be,
\begin{align*}
\nu_{\mathrm{max}} &= \frac{c}{\lambda_{\mathrm{max}}}
  = 9.369 \times 10^{14}\,\mathrm{Hz} \\
\nu_{\mathrm{min}} &= 7.495 \times 10^{14}\,\mathrm{Hz} \\
\nu_{\mathrm{avg}} &= 8.432 \times 10^{14}\,\mathrm{Hz} \\
\Delta\nu &= \nu_{\mathrm{max}} - \nu_{\mathrm{min}} \\
  &= 1.874 \times 10^{14}\,\mathrm{Hz}
\end{align*}
\hfill(2.0 points)
\begin{align*}
f_{\mathrm{st1}} &= 3631 \times 10^{-0.4 \times 15} \\
  &= 3.631\,\mathrm{mJy}
   = 3.631 \times 10^{-29}\,\mathrm{W\,Hz^{-1}\,m^{-2}}
\end{align*}
\hfill(2.0 points)

Now, $N_0 \times h\nu_{\mathrm{avg}} = \Delta\nu \times f_{\mathrm{st1}}
\times A \times \Delta t$, \hfill(2.0 points)\\
where $A = 1\,\mathrm{m}^2$ \& $\Delta t = 1$\,s
\begin{align*}
\therefore N_0 &= \frac{1.874 \times 10^{14} \times 3.631 \times 10^{-29}}
  {6.626 \times 10^{-34} \times 8.432 \times 10^{14}} \\
  &\simeq 12180
\end{align*}
\hfill(2.0 points)

Exact calculation including integration is accepted with full credit (exact
answer: 12190).

\item[(T8.2)]
Let us call sky flux per square arcsec as $\Phi$ and total sky flux for the
given aperture as $\phi_{\mathrm{sky}}$. Let total star flux be
$\phi_{st}$.
\[ \phi_{\mathrm{sky}} = A\Phi = \pi \times (1\,\mathrm{arcsec})^2 \times \Phi
   = \pi\Phi \]
\hfill(3.0 points)
\begin{align*}
\therefore m_{\mathrm{sky}}
  &= 22.0 + 2.5\log_{10}\left(\frac{\Phi}{\phi_{\mathrm{sky}}}\right)
  && \text{(1.0 points)} \\
  &= 22.0 + 2.5\log_{10}\left(\frac{\Phi}{\pi\Phi}\right) \\
  &= 22.0 - 2.5\log_{10}(\pi) \\
m_{\mathrm{sky}} &= 20.76~\mathrm{mag}
  && \text{(1.0 points)}
\end{align*}
As extinction is 50\% \hfill(1.0 points)
\begin{align*}
R &= \frac{\phi_{\mathrm{st2}}}{\phi_{\mathrm{sky}}}
   = \frac{0.5\,\phi_{\mathrm{st1}}}{\phi_{\mathrm{sky}}}
   = 0.5 \times 10^{(20.76 - 15)/2.5} \\
  &\simeq 100
\end{align*}
\hfill(2.0 points)

A student may also calculate number of photons incident per second per metre
square in this case and then compare it with the answer in the first case to
get the correct ratio.

\item[(T8.3)]
\begin{align*}
N_{\mathrm{t}} \times 1\,\mathrm{m}^2
  &= N_0 \times 0.5 \times 0.2 \times A_{\mathrm{t}}
  && \text{(2.0 points)} \\
N_{\mathrm{t}} &= 12180 \times 0.5 \times 0.2 \times
  \pi\left(\frac{2.0}{2}\right)^2 = 1233\pi
  && \text{(1.0 points)} \\
N_{\mathrm{t}} &\simeq 3813
  && \text{(1.0 points)}
\end{align*}
\end{description}

\solhead{5.15}{GOTO}
% source id: 2017_TH_T10
\begin{description}
\item[a)]
\ioaafig{2017_TH_T10-s1.pdf}
For any two points on the sky separated by small angular distance $\theta$,
at the focal plane of the telescope the distance is $S$
\[ \tan\theta = \frac{S}{f}, \]
\hfill(0.5 marks)\\
where $f$ is the focal length.

For small angle $\theta$
\[ \tan\theta \approx \theta \]
Then plate scale is
\[ \frac{\theta}{S} = \frac{1}{f} \]
\hfill(0.5 marks)
\[ f = \text{F-ratio} \times \text{Diameter} \]
\hfill(1 mark)
\[ \text{Plate scale} = \frac{1}{\text{F-ratio} \times \text{Diameter}} \]
\hfill(1 mark)
\begin{align*}
\text{Plate Scale}
  &= \frac{1}{2.5 \times 40 \times 10^{-2}\,(\mathrm{m})
     \times 10^{3}\,(\mathrm{mm\,m^{-1}})}
     \times \frac{180 \times 3600\,(\mathrm{arcsec})}{\pi}
  && \text{(1 mark)} \\
  &= 206~\mathrm{arcsec/mm} \\
\text{Plate Scale} &= 3.44~\mathrm{arcmin/mm}
  && \text{(1 mark)}
\end{align*}

\item[b)]
From plate scale calculated in a), size of a star for typical seeing 1\,arcsec
\begin{align*}
  &= \frac{1/60~\mathrm{arcmin}}
     {\left(3.44~\mathrm{arcmin/mm} \times 6 \times 10^{-3}~\mathrm{mm}\right)}
     ~\text{pixels}
  && \text{(1 mark)} \\
  &= 0.8~\text{pixel}
\end{align*}
$\Rightarrow$ light from star is mostly contained within
$N_{\mathrm{pix}} = 1$ pixel. \hfill(1 mark)

\emph{Alternative solution:} If a student gets the correct answer for
1\,arcsec $= 0.8$ pixel but clearly explains that the chosen aperture size is
$N \times$ ``seeing'' or a circular aperture with radius $R_{N\mathrm{pix}}$
(where $R_{N\mathrm{pix}}$ is not larger than $2\times$ seeing), calculates
the total number of pixels as the closest round-up integer, and correctly uses
\[ \mathrm{SNR} = \frac{\mathrm{CR} \times t}
   {\sqrt{N_{\mathrm{pix}}\,\mathrm{RON}^2
    + (N_{\mathrm{pix}} \times \mathrm{DN} \times t)}} \]
then the student will be awarded full marks for the relevant parts (this
also applies to part c).

From the definition of the zero-point magnitude given in the question
\[ m = -2.5\log_{10}(\text{count rate}) + \text{Zero-Point} \]
\hfill(2 marks)
\[ m = 21,~\mathrm{ZP} = 18.5 \]
$\Rightarrow$ source count rate
$\mathrm{CR} = 10^{-(21-18.5)/2.5} = 0.1$ per second \hfill(1 mark)

Signal-to-noise ratio
\begin{gather*}
\mathrm{SNR} = \frac{\mathrm{CR} \times t}
  {\sqrt{\mathrm{RON}^2 + (\mathrm{DN} \times t)}} \\
(\mathrm{CR} \times t)^2 = 25(\mathrm{RON}^2 + (\mathrm{DN} \times t))
\end{gather*}
\hfill(1 mark)

Solve quadratic equation
\[ t = \frac{25\,\mathrm{DN}
   + \sqrt{(25\,\mathrm{DN})^2 + 100\,\mathrm{CR}^2\,\mathrm{RON}^2}}
   {2\,\mathrm{CR}^2} \]
\hfill(1 mark)
\[ t = \frac{\dfrac{25}{60}
   + \sqrt{\left(\dfrac{25}{60}\right)^2 + 100\,(0.1)^2\,10^2}}
   {2 \times 0.1^2}
   = 521~\mathrm{s} = 8.68~\text{minutes} \]
\hfill(1 mark)

\item[c)]
If we include the source Poisson noise,
\[ \mathrm{SNR} = \frac{\mathrm{CR} \times t}
   {\sqrt{\mathrm{RON}^2 + (\mathrm{DN} \times t) + \mathrm{CR} \times t}} \]
\hfill(1 mark)

For the Poisson noise dominated case
\[ \mathrm{SNR} \approx \frac{\mathrm{CR} \times t}
   {\sqrt{\mathrm{CR} \times t}} \]
\hfill(1 mark)

Hence
\[ \mathrm{SNR} \propto t^{\frac{1}{2}} \]
\hfill(1 mark)

Re-calculate exposure time
\begin{gather*}
\mathrm{SNR} = \frac{\mathrm{CR} \times t}
  {\sqrt{\mathrm{RON}^2 + (\mathrm{DN} \times t) + \mathrm{CR} \times t}}, \\
(\mathrm{CR} \times t)^2
  = 25(\mathrm{RON}^2 + (\mathrm{DN} + \mathrm{CR})t)
\end{gather*}
\hfill(1 mark)

Solve quadratic equation
\[ t = \frac{25(\mathrm{DN} + \mathrm{CR})
   + \sqrt{25^2(\mathrm{DN} + \mathrm{CR})^2
   + 100\,\mathrm{CR}^2\,\mathrm{RON}^2}}{2\,\mathrm{CR}^2} \]
\hfill(1 mark)
\[ t = 666~\mathrm{s} = 11.1~\text{minutes} \]
\hfill(1 mark)

\item[d)]
In 1\,hr, one can observe $\dfrac{60}{11.1} = 5.4$ \hfill(1 mark)\\
$\Rightarrow$ 5 pointings (round down to the nearest integer). \hfill(1 mark)

Therefore, we need to cover $100~\mathrm{deg}^2$ $\Rightarrow$ one pointing
needs to cover
\[ \frac{100}{(5 \times 4~\text{telescope})} = 5~\mathrm{deg}^2 \]
\hfill(1 mark)\\
or $2.24 \times 2.24$ deg per pointing per telescope. \hfill(1 mark)

Calculate size of each CCD
\begin{align*}
  &= \frac{2.24~\mathrm{deg} \times 60}
     {(3.4~\mathrm{arcmin/mm})(6 \times 10^{-3}~\mathrm{mm})}~\text{pixels}
  && \text{(1 mark)} \\
  &= 6588.2~\text{pixels}
   \;\Rightarrow\; 6589~\text{pixels (pixel number needs to be integer)}
\end{align*}
Therefore, we need minimum $\mathrm{CCD~size} = 6589 \times 6589$ pixels.
\hfill(1 mark)
\end{description}

\solhead{5.16}{Radio Galaxy}
% source id: 2018_TH_T9
\begin{description}
\item[(a)] For signal to noise ratio of 30 of flux density of
$2.5 \times 10^{-3}$\,Jy,
\[ \sigma = 2.5 \times 10^{-3}/30~\mathrm{Jy}
   \approx 8.33 \times 10^{-5}~\mathrm{Jy}
   = 8.33 \times 10^{-31}~\mathrm{W\,m^{-2}\,Hz^{-1}} \]
\hfill(4 points)
\[ T_{\mathrm{sys}} = 150~\mathrm{K}, \quad
   \Delta\nu = 280~\mathrm{MHz}, \]
Using the equation above, $t_{\mathrm{int}} \sim 42$\,s (more exactly
41.8\,s). \hfill(3 points)

\item[(b)] $\nu_{\mathrm{obs}} = 1.4204~\mathrm{GHz}/(1+z)
= 1.4204/1.06 = 1.34~\mathrm{GHz}$ \hfill(3 points)

\item[(c)]
\[ f_{\nu_{\mathrm{obs}}=1.34\,\mathrm{GHz}}
   = f_{0.4\,\mathrm{GHz}}\left(\frac{1.34}{0.4}\right)^{-0.2}
   = 2.5 \times 10^{-3}\,\mathrm{Jy}
     \times \left(\frac{1.34}{0.4}\right)^{-0.2}
   = 1.96 \times 10^{-3}\,\mathrm{Jy} \]
\hfill(3 points)

\item[(d)] Line width:
$90~\mathrm{km\,s^{-1}}/c \times \nu_{\mathrm{obs}}
 = 90/(2.9979 \times 10^{5}) \times 1.34~\mathrm{GHz}
 = 0.402~\mathrm{MHz}$ \hfill(2 points)

4\% of $1.96 \times 10^{-3}$\,Jy continuum
$= 7.84 \times 10^{-5}$\,Jy \hfill(1 point)

$\geq 3\sigma$ means
$\sigma = 7.84 \times 10^{-5}~\mathrm{Jy}/3
 = 2.61 \times 10^{-5}~\mathrm{Jy}
 = 2.61 \times 10^{-31}~\mathrm{W\,m^{-2}\,Hz^{-1}}$
(values around $2.6 \times 10^{-5}$\,Jy are ok)

in three consecutive 30\,km\,s$^{-1}$ channels:
$\Delta\nu = 0.402~\mathrm{MHz}/3 = 0.134~\mathrm{MHz}$ \hfill(6 points)
\[ T_{\mathrm{sys}} = 25~\mathrm{K} \]
Using equation (1), $t_{\mathrm{int}} = 24700\,\mathrm{s} \approx 6.9$ hours.
\hfill(3 points)
\end{description}

\solhead{5.17}{Photographing a nanosatellite}
% source id: 2019_TH_14
\begin{description}
\item[a)] As it was suggested in the hint the key is a comparison of the
sunlight scattered from a $10\,\mathrm{cm} \times 10\,\mathrm{cm}$ aluminium
plate and the Moon to the observer (both are nearly at the local zenith).

Let the distance of the Moon and the satellite from the observing site be
$d_{\mathrm{Moon}}$, and $d_{\mathrm{MASAT}}$, respectively. Under the
mentioned "ideal" circumstances, when both the Moon and the satellite is near
the local zenith: $d_{\mathrm{MASAT}} = 900$\,km, while
\begin{equation*}
d_{\mathrm{Moon}} = s_{\mathrm{Moon}} - R_{\mathrm{Earth}}
  - R_{\mathrm{Moon}}, \tag{14.1}
\end{equation*}
\hfill(2 p)\\
where the average distance between the centre of Earth and Moon
$s_{\mathrm{Moon}}$ is included in the \emph{Table of constants}.

Let the intensity coming from the Sun and scattered back to the direction of
the observer be $I$, this is considered practically the same at the solar
distance of the Moon and the satellite, but necessarily multiplied by the
albedo of the given surface quality (the lunar albedo $a_{\mathrm{Moon}}$ is
given in the Constant table, while $a_{\mathrm{MASAT}}$ of the satellite is
in the text).

Finally, for using the magnitude formula, to find the observable magnitude of
the satellite, we have to consider the resulted flux of the scattered light
at the observer, which is directly proportional to the intensity, the albedo
and the scattering area, but is inversely proportional to the distance on the
square:
\begin{equation*}
I_{\mathrm{scat}} \approx
  \frac{a_{\mathrm{obj}}\,I\,A_{\mathrm{obj}}}{d_{\mathrm{obj}}^2}
\tag{14.2}
\end{equation*}
\hfill(3 p)

Thus, the magnitude formula gives:
\begin{equation*}
m_{\mathrm{MASAT}} - m_{\mathrm{Moon}}
  = -2.5\log\frac{I_{\mathrm{MASAT}}}{I_{\mathrm{Moon}}}
  = -2.5\log\frac{\dfrac{a_{\mathrm{MASAT}}\,I\,A_{\mathrm{MASAT}}}
      {d_{\mathrm{MASAT}}^2}}
    {\dfrac{a_{\mathrm{Moon}}\,I\,A_{\mathrm{Moon}}}{d_{\mathrm{Moon}}^2}}
\tag{14.3}
\end{equation*}
\hfill(3 p)

With numerical values:
\[ m_{\mathrm{MASAT}} = 9.72^{\mathrm{m}} \]
\hfill(2 p)

\item[b)] A moving point-like object will generally cause a line on the CCD
image, resulting less photoelectrons in a given pixel, than a similarly
bright unmoving object. Since the MASAT--1 satellite is so small, it is not
resolved by the telescope -- it will cause a seeing-sized spot at the focus
of a telescope, like any other stars on the sky, but since the text of the
problem described: we should omit the seeing. Thus, the image of any point
source will be formed by the diffraction -- so we shall use the approximate
formula of the Airy disk.

By the diffraction-limited imaging, the point-like sources illuminate a spot
in the focal plane having the angular diameter of the so-called Airy disc
determined by the formula:
\begin{equation*}
\varphi_{\mathrm{Airy}} = 2.44\,\frac{\lambda}{D}, \tag{14.4}
\end{equation*}
\hfill(2 p)\\
while its linear size at the focus of $f$ focal length telescope is:
\begin{equation*}
d_{\mathrm{Airy}} = f\varphi_{\mathrm{Airy}} \tag{14.5}
\end{equation*}
\hfill(2 p)

For the telescope of Baja Observatory, described in the text, at a typical
visual wavelength of $0.55\,\mu\mathrm{m}$:
\[ d_{\mathrm{Airy}} = 2.44\,\frac{f\lambda}{D}
   = 2.44 \times \frac{4200\,\mathrm{mm} \times 0.55\,\mu\mathrm{m}}
     {500\,\mathrm{mm}}
   = 11.27\,\mu\mathrm{m} \]
\hfill(2 p)

This is very close to the pixel size of the applied CCD. However, this
number is not necessary for our calculation, since both the stars and the
MASAT--1 are point-like sources, illuminating the same number of pixels, thus
the number of pixels will drop out from all comparative calculations. So
does the quantum efficiency value of the applied CCD.

For deriving the effect of the moving object, let us first calculate the
scale of the image at the focus of the 50\,cm RC telescope of Baja ($N$ is
the f-number of the telescope, $N = f/D$):
\[ S = \frac{206\,265''}{DN}
     = \frac{206\,265''}{500\,\mathrm{mm} \times 8.4}
     = 49\,''\,\mathrm{mm^{-1}} = 0.049\,''\,\mu\mathrm{m}^{-1} \]
\hfill(3 p)

This means, that one pixel of the attached CCD describes
$P = 0.049\,''\,\mu\mathrm{m}^{-1} \times 9\,\mu\mathrm{m} = 0.44''$
side length area on the sky. \hfill(3 p)

Now let us estimate the relative flux, resulted by a star having the limiting
magnitude of the telescope-CCD system ($19.5^{\mathrm{m}}$), by the magnitude
formula:
\begin{equation*}
19.5 = -2.5\log\Phi_{19.5} \rightarrow
  \Phi_{19.5} = 10^{-7.8}\,\mathrm{s^{-1}} \tag{14.6}
\end{equation*}
\hfill(2 p)

And similarly, for the MASAT--1 brightness -- observed with the same system,
under the same conditions, if it were an unmoving object:
\begin{equation*}
9.72 = -2.5\log\Phi_{\mathrm{MASAT}} \rightarrow
  \Phi_{\mathrm{MASAT}} = 10^{-3.9}\,\mathrm{s^{-1}} \tag{14.7}
\end{equation*}
\hfill(2 p)

Next, we have to calculate the angular speed of MASAT--1 at 900\,km altitude
above the observing site, in order to derive the fluxes falling onto one
pixel of the CCD on the telescope.

The velocity of the motion on a circular orbit with orbital radius $r$:
\begin{equation*}
v_r = \sqrt{\frac{GM_{\mathrm{Earth}}}{r}}, \tag{14.8}
\end{equation*}
\hfill(2 p)\\
where $r = h + R_{\mathrm{Earth}}$. \hfill(2 p)

With numerical values:
\[ v_r = 7404\,\mathrm{m\,s^{-1}} \]
\hfill(1 p)

Hence, the visible angular speeds (aside from the Earth rotation):
\begin{equation*}
\omega_r = \frac{v_r}{r} \tag{14.9}
\end{equation*}
\hfill(2 p)
% sic: the source prints omega_r = v_r / r; for the angular speed seen from
% the observing site directly below, the distance should be the altitude h.
With numerical values:
\[ \omega_r = 0.47\,\mathrm{deg\,s^{-1}} = 1696.8\,''\,\mathrm{s^{-1}}, \]
\hfill(1 p)\\
when the satellite is near the local zenith.

For deriving the relative speed, which will be really exposed onto the CCD,
we have to calculate the angular speed of the observing site caused by the
terrestrial rotation. At the latitude of $\varphi$ it is given as:
\begin{equation*}
\omega_{\mathrm{rot},\varphi} = \omega_{\mathrm{rot,eq}}\cos\varphi
\tag{14.10}
\end{equation*}
\hfill(2 p)

For Baja its numerical value:
\[ \omega_{\mathrm{rot,B}} = 10.4\,''\,\mathrm{s^{-1}} \]
\hfill(1 p)

Its component parallel to the direction of the satellite's orbital motion:
\begin{equation*}
\omega_{\mathrm{rot,B},i} = \omega_{\mathrm{rot,B}}\cos i_{\mathrm{MASAT}}
\tag{14.11}
\end{equation*}
\hfill(2 p)

For Baja its numerical value:
\[ \omega_{\mathrm{rot,B},i} = 3.6\,''\,\mathrm{s^{-1}} \]
\hfill(1 p)

Finally, the relative angular speed of the nanosatellite above the telescope:
\begin{equation*}
\Delta\omega = \omega_a - \omega_{\mathrm{rot,B},i} \tag{14.12}
\end{equation*}
\hfill(2 p)

With numerical values:
\[ \Delta\omega = 1693.3\,''\,\mathrm{s^{-1}} \]
\hfill(1 p)

\ioaafig{2019_TH_14-s1.pdf}

Since the satellite image is moving with a given $\Delta\omega$ angular
speed, this means that the flux coming from the satellite is falling onto a
given pixel only for a very limited time (independently of the exposure time
of the image, but acting similarly, than that). Considering our very simple
approximation we get an approximate "illumination time" for one pixel falling
onto the line of the orbit:
\begin{equation*}
\tau = \frac{P}{\Delta\omega} \tag{14.13}
\end{equation*}
\hfill(2 p)

With numerical values:
\[ \tau = \frac{0.44''}{1693.3\,''\,\mathrm{s^{-1}}}
        = 2.61 \times 10^{-4}\,\mathrm{s} \]
\hfill(1 p)

For finishing our estimation -- to get the magnitude level of any pixel on
the observable trace of this nanosatellite, let us use the magnitude formula
for them and one non-moving light source at the detection limit. For doing
this we have to multiply the relative fluxes $\Phi_{\mathrm{MASAT}}$ and
$\Phi_{19.5}$ with the respective illumination times. (For a
$19.5^{\mathrm{m}}$ stars we should use the 2\,min-long exposure time
mentioned in the text, which is needed to collect enough photoelectrons for
the safe detection of such a faint star.)
\begin{equation*}
m_{\mathrm{MASAT,trace}} - 19.5^{\mathrm{m}}
  = -2.5\log\frac{\tau\Phi_{\mathrm{MASAT}}}
    {\tau_{\mathrm{exp}}\Phi_{19.5}} \tag{14.14}
\end{equation*}
\hfill(2 p)

With numerical values:
\[ m_{\mathrm{MASAT,trace}} = 23.9^{\mathrm{m}} \]
\hfill(2 p)

So, the trace of the moving $23.9^{\mathrm{m}}$ magnitude nanosatellite looks
so faint on a (2\,min-long exposure) CCD image, as a $\sim 24^{\mathrm{m}}$
stars would be seen, but they fall surely below the detection limit for the
described telescope. Thus, the answer is definitely

\textbf{NO}, WE COULD NOT RECORD THE MASAT--1 \hfill(2 p)\\
using that telescope.

\item[c)] If we consider a real atmosphere, due to the effect of the
atmospheric turbulences (the so-called seeing) during the 2 minutes-long
exposure the images of the stars will never be described by the ideal Airy
disk, but expanded to a much larger "seeing spot", which can be approximated
by a Gaussian profile. As mentioned in the text: in Hungary, its diameter at
half-maxima is typically $3.5''$. This means, that the flux coming from a
$19.5^{\mathrm{m}}$ star will be distributed over about a
\[ \frac{\mathrm{FWHM}}{P} \times \frac{\mathrm{FWHM}}{P}
   = \frac{3.5''}{0.44''} \times \frac{3.5''}{0.44''} \approx 8 \times 8 \]
\hfill(1 p)\\
pixels area. Thus, now the $F_{19.5}$ flux will be distributed amongst
$4 \cdot 4\pi \approx 50$ pixels (for the sake of simplicity:
homogeneously), i.e.\ one pixel gets:
\begin{equation*}
\Phi_{19.5} = \frac{10^{-7.8}}{50}\,\mathrm{s^{-1}}
  \approx 10^{-9.5}\,\mathrm{s^{-1}} \tag{14.15}
\end{equation*}
\hfill(1 p)

As we have seen above, the satellite image is moving so fast along the focal
plane (0.26\,ms per a pixel), that during this short time, the atmosphere can
be considered to be still. Thus, the image of MASAT--1 can further be
approximated with the Airy-disc size also in a real atmosphere.

Let us correct our last magnitude formula for this fact:
\begin{equation*}
m_{\mathrm{MASAT,trace}} - 19.5^{\mathrm{m}}
  = -2.5\log\frac{\tau\Phi_{\mathrm{MASAT}}}
    {\tau_{\mathrm{exp}}\dfrac{\Phi_{19.5}}{50}} \tag{14.16}
\end{equation*}
\hfill(2 p)

With numerical values:
\[ m_{\mathrm{MASAT,trace}} = 19.6^{\mathrm{m}} \]
\hfill(2 p)

Under real conditions the situation is much better: the trace of the small
nanosatellite is near the detectability! Thus, owing to further secondary
effects (pixel sensitivity inhomogenities, temporarily changing transparency
of the atmosphere, etc.), the trace of the cubesat can easily be seen on the
CCD of the 50\,cm reflector. So in real cases, the answer can be YES.

\textbf{YES}, WE COULD HAVE RECORDED THE MASAT--1. \hfill(2 p)
\end{description}

\solhead{5.18}{Improving a common reflecting telescope}
% source id: 2019_TH_4
\begin{description}
\item[a)] Let us designate the (unknown) original limiting magnitude of this
telescope as $m_1$, and refer to the intensity (at the focal plane) as $I_1$.
The new limiting magnitude is $m_2$. Due to the additional intensity, which
is due to the better reflectivity of the mirror surfaces, the stellar
intensity which will result in the same limiting illumination at the focal
plane, will be proportionally smaller, with the ratio of the reflectivities
of the new and old mirror coatings:
\begin{equation*}
I_2 = \frac{\varepsilon_1\varepsilon_1}{\varepsilon_2\varepsilon_2}\,I_1
\tag{4.1}
\end{equation*}
\hfill(2 p)

The reflectivities should be multiplied twice, since the light coming from
the star, loses some energy twice, suffering reflection on each of the two
mirrors.

Now applying the well-known magnitude formula:
\begin{equation*}
m_1 - m_2 = -2.5\log\frac{I_1}{I_2}
          = -2.5\log\frac{\varepsilon_2^2}{\varepsilon_1^2}
\tag{4.2}
\end{equation*}
\hfill(2 p)

With numerical values:
\[ m_1 - m_2 = -0.16^{\mathrm{m}} \]
And thus:
\[ m_2 = m_1 + 0.16^{\mathrm{m}} \]
\hfill(1 p)

\item[b)] % sic: the official solution swaps parts b) and c) relative to
% the question paper -- this item answers question c) (eye detectability),
% and the next item answers question b) (the star diagonal calculation).
Yes, it is well appreciable by most human eyes. Therefore,
especially the deep-sky-object hunters are using higher reflectivity mirrors
in their telescopes. \hfill(2 p)

\item[c)] The similar way of thinking gives:
\begin{equation*}
m_2 = m_1 + 2.5\log\frac{I_1}{I_2}
    = m_1 + 2.5\log\frac{\varepsilon_2^2\varepsilon_3}{\varepsilon_1^3}
\tag{4.3}
\end{equation*}
\hfill(2 p)

With numerical values:
\[ m_1 - m_2 = -0.25^{\mathrm{m}} \]
And thus:
\[ m_2 = m_1 + 0.25^{\mathrm{m}} \]
\hfill(1 p)
\end{description}

\solhead{5.19}{Astrophotography}
% source id: 2020_TH_1
\begin{description}
\item[(a)] By simple right angle triangle,
\begin{align*}
\mathrm{FOV} &= 2 \times \tan^{-1}
  \left(\frac{\text{sensor width}}{2 \times \text{Focal Length}}\right)
  = 2 \times \tan^{-1}\left(\frac{22.5}{2 \times 150}\right) \\
\alpha &\approx 8.58^\circ
\end{align*}
\hfill(2.0 points)

\item[(b)] Number of pixels on the sensor will be given by,
\begin{align*}
N &= \frac{\text{Sensor width}}{\text{pixel width}}
   = \frac{22.5}{3.23 \times 10^{-3}} \\
  &= 6965.94
\end{align*}
\hfill(1.0 points)

As the number of pixels have to be integer, we take $N = 6966$.
\hfill(1.0 points)\\
Angular coverage of each pixel will be,
\begin{align*}
\text{Pixel view} &= \frac{8.578^\circ}{6966} = 0.00123^\circ/\text{pixel} \\
  &\approx 4.43''/\text{pixel}
\end{align*}
\hfill(1.0 points)

As the stars complete one full circle in 23.9344 hours, \hfill(1.0 points)
\begin{align*}
t_1 &= 1.23 \times 10^{-3}\,\frac{23.9344^{h}}{360} = 0.294\,\mathrm{s} \\
  &\approx 0.3\,\mathrm{s}
\end{align*}
\hfill(1.0 points)

(in reality this is probably a factor of 10 smaller than the eye would
detect)

\item[(c)] In 30 seconds, the sky will move by,
\[ \Delta\alpha = \frac{360^\circ}{23.9344 \times 120} = 0.125\,342^\circ \]
\hfill(1.0 points)
\begin{align*}
\therefore \frac{\alpha - \Delta\alpha}{\alpha}
  &= \frac{8.578 - 0.125342}{8.578} \\
  &= 98.5\%
\end{align*}
\hfill(2.0 points)

(or $8.453^\circ$ is total field of view in high resolution images from the
stack)
\end{description}

\solhead{5.20}{Lucy: The First Mission to the Trojan Asteroids}
% source id: 2021_TH_Q13
\begin{description}
\item[(13.1)] A single electron deposits energy in the CCD, as follows:
\[ E_{\mathrm{deposited}}
   = \text{stopping power} \times \rho_{\mathrm{si}} \times \text{thickness} \]
\[ E_{deposited}
   = 3.012\,\frac{MeV\,cm^2}{g} \times \frac{2.34\,g}{cm^3}
     \times 0.06\,cm = 422.9~keV \]
\hfill(4 points)

Since an electron with 15\,MeV energy deposits 422.9\,keV, we must calculate
the number of pairs electron/hole
\[ 422.9\,keV \times \frac{1}{2.36\,eV}
   = 1.79 \times 10^{5}\,\frac{e}{h} \]
\hfill(4 points)

How many pixels will be able to excite $1.79 \times 10^{5}$ pairs of e/h?
\[ 1.79 \times 10^{5} \times \frac{1}{250} = 716~pixels \]
\hfill(2 points)

\item[(13.2)] The number of electrons entering the detector area is:
\[ flux \,\times\, area \,\times\, t \]
\hfill(1 point)
\[ \frac{600\,e}{cm^2\,s} \times [1.3 \times 1.3]\,cm^2 \times 0.03\,s
   = 30\,e^{-} \]
\hfill(1 point)

$\sim$30 electrons enter the detector area, and we know that one electron
excites 716 pixels, the total number of excited pixels is:
\[ 716\,\frac{pixels}{electron} \times 30.42~electrons
   = 2.18 \times 10^{4}\,pixels \]
\hfill(1 point)

If the CCD has a total of $1024 \times 1024$ pixels
$= 1.048 \times 10^{6}\,pixels$ then the total percentage of pixels that
will be driven in a single image is $\sim 2.08\%$. \hfill(2 points)
% sic: "driven" (presumably "excited") as printed in the source.
\end{description}

\solhead{5.21}{ALMA - Calculating photons}
% source id: 2021_TH_Q4
\begin{description}
\item[(4.1)] To get the number of photons per second we have to multiply the
Flux by the area of the dish, to know the incoming energy per second, and
divide this by the energy of a photon.

For $\lambda_1 = 0.32\,mm = 3.2 \times 10^{-4}\,m$:

Energy of a photon:
\[ E = \frac{hc}{\lambda} = 6.2076 \times 10^{-22}\,Jules \]
% sic: the source spells "Jules" for Joules.
\hfill(1 point)

Area of the disk:
\[ A = \pi R^2 = 113\,m^2 \]

Number of photons per second:
\[ n_1 = \frac{flux \times A}{E} \approx 1820~photon/s \]
\hfill(1 point)

\item[(4.2)] Same calculation, just changing the energy of the individual
photons:

For $\lambda_2 = 8.6\,mm = 8.6 \times 10^{-3}\,m$

Energy of a photon:
\[ E = \frac{hc}{\lambda} = 2.31 \times 10^{-23}\,Jules \]
\hfill(1 point)

Number of photons per second:
\[ n_2 = \frac{flux \times A}{E} \approx 48900~photon/s \]
\hfill(1 point)

\item[(4.3)] Spatial resolution of a single telescope is given by:
\[ \theta = 1.22\,\frac{\lambda}{D} \]
where $D$ represents the diameter of the dish. This value will be given in
radians, so it must be converted to arcsec afterwards.

For a frequency of 74.9\,GHz, the corresponding wavelength is:
\[ \lambda = \frac{c}{f} = 4\,mm \]
\hfill(1 point)

And the spatial resolution:
\[ \theta = 1.22\,\frac{4\,mm}{12\,m} \approx 83.9~arcsec \]
\hfill(1 point)

\item[(4.4)] For an array the correct expression is:
\[ \theta = \frac{\lambda}{B} \quad\text{or}\quad
   \theta = \frac{\lambda}{2B} \]
\hfill(1 point)\\
being $B$ the longest baseline in the array.

So in this case:
\[ \theta = \frac{4\,mm}{16\,km} \approx 0.0516~arcsec
   \quad\text{or}\quad 0.0258~arcsec \]
\hfill(1 point)

\item[(4.5)] The SEFD is a characteristic flux of a system, found by
dividing the characteristic energy associated to the so-called temperature of
the antenna by its effective area:
\[ SEFD = \frac{2kT_{sys}}{A_e} \]
\hfill(1 point)

As no additional information is given about the effective area, the actual
physical area of the array should be used:
\[ A = 54\left(\pi \times 6^2\right) + 12\left(\pi \times 3.5^2\right)
   \approx 6569\,m^2 \]
\hfill(0.5 points)

Substituting the Boltzmann constant, the given temperatura of the antenna,
and converting the answer to Jansky we get:
% sic: "temperatura" appears in the source.
\[ SEFD \approx 290.5~Jy \]
\hfill(0.5 points)
\end{description}

\solhead{5.22}{Microlensing}
% source id: 2023_TH_Q3
Let $F$ be the unmagnified flux of the star. During gravitational
microlensing, the apparent flux is magnified by a factor of $A$, thus using
the formula relating magnitude to flux:
\[ m_1 - m_2 = -2.5\log_{10}\left(\frac{F_1}{F_2}\right), \]
\hfill(1 point)
\[ I - I' = -2.5\log_{10}\left(\frac{F}{AF}\right) = 2.5\log A \]
\hfill(1 point)\\
and so
\[ A = 10^{2.08} \approx 120. \]
\hfill(1 point)

Let $D$ be the effective mirror diameter of a telescope which would collect
the same number of photons in unit time from the unbrightened star as VLT
from the brightened star. We have:
\[ FD^2 = AFD_{\mathrm{VLT}}^2, \]
\hfill(1 point)\\
therefore:
\[ D = \sqrt{A}\,D_{\mathrm{VLT}} \approx 90\,\mathrm{m}. \]
\hfill(1 point)

\solhead{5.23}{Bolometer}
% source id: 2023_TH_Q6
From the table of constants, the Solar luminosity
$L_\odot = 3.826 \times 10^{26}$\,W and
$1\,\mathrm{au} = 1.496 \times 10^{11}$\,m, therefore the distance to the
bolometer is:
\[ r = 2\,\mathrm{au} = 2.992 \times 10^{11}\,\mathrm{m} \]
and the surface area $A$ of a sphere of that radius is:
\[ A_{\mathrm{sphere}} = 4\pi r^2 = 1.125 \times 10^{24}\,\mathrm{m}^2 \]
and the incident flux is:
\[ F(r) = L_\odot/A_{\mathrm{sphere}} = 340.1\,\mathrm{W\,m^{-2}} \]
\hfill(3 points)

The incoming radiation can be assumed to be initially parallel to the axis of
the cone. As the rays hit the surface they are reflected, after the first
reflection the rays are travelling at $30^\circ$ to their original path
($2 \times 15^\circ$). They next meet the surface at
$30^\circ + 15^\circ = 45^\circ$ and are reflected by $90^\circ$, and so on
with each reflection, such that after $N = 6$ reflections the ray will be
sent back out of the aperture as shown in the figure below. These exiting
rays are what we are interested in; all the energy entering the bolometer
itself is removed by the 100\% efficient cooler.

\ioaafig{2023_TH_Q6-s1.pdf}

\hfill(\emph{conceptually difficult:} 4 points)

Thus the fraction of energy leaving the bolometer will be given by:
\[ S = (1-a)^N = 0.01^6 = 10^{-12}. \]
\hfill(2 points)

Treating the opening as if it were a black body radiator, we can apply the
Stefan-Boltzmann law:
\begin{gather*}
\Phi = \sigma T^4 \\
\implies T = (\Phi/\sigma)^{0.25},
\end{gather*}
\hfill(1 point)\\
and
\[ \Phi = F(r) \cdot S = 340.1 \times 10^{-12}\,\mathrm{W\,m^{-2}} \]
\hfill(2 points)

From the table of constants:
$\sigma = 5.670 \times 10^{-8}\,\mathrm{W\,m^{-2}\,K^{-4}}$, therefore the
effective temperature is:
\[ T = \left(\frac{340.1 \times 10^{-12}}{5.670 \times 10^{-8}}\right)^{0.25}
     = 0.28\,\mathrm{K} \]
\hfill(1 point)

\solhead{5.24}{Cluster Photography}
% source id: 2024_TH_T6
\begin{description}
\item[(a)] The plate scale is simply the angular size of the image per unit
length as projected at the focal plane. Hence, for small angles,
\[ \tan\theta_i = \frac{\ell_i}{f} \approx \theta_i \]
\hfill(1.0 points)\\
where $\theta_i$ is the angular size of the image, $\ell_i$ is the unit
length, and $f$ is the focal length of the telescope, which is given by
($D \times$ F-ratio). \hfill(0.5 points)
\begin{gather*}
\theta_i = \frac{1\,\mathrm{mm}}{5 \times 305\,\mathrm{mm}}
  = 6.56 \times 10^{-4}\,\mathrm{rad} = 2.254' \\
PS \approx 2.25\,\mathrm{arcmin/mm}
\end{gather*}
\hfill(1.5 points)

\item[(b)] Given the plate scale, the diameter of the image on the focal
plane is
\[ d_{GC} = \frac{72.0''}{2.25' \times 60} = 0.532\,\mathrm{mm} \]
\hfill(1.0 points)

This implies that the area covered by the image is
\[ S_{GC} = \frac{\pi}{4} \cdot d_{GC}^2 = 0.223\,\mathrm{mm}^2 \]
which can be divided by the area of a single pixel to estimate the pixel
coverage
\[ n_p = \frac{0.223}{(3.8 \times 10^{-3})^2}
   \approx 15\,400~\text{pixels} \]
\hfill(2.0 points)

Since there are other methods to estimate such quantity, the accepted range
is
\[ 15300 \leqslant n_p \leqslant 15600 \]
\hfill(1.0 points)

\item[(c)] The signal-to-noise ratio is given by the expression
\begin{align*}
S/N &= \frac{N_{GC}}
  {\sqrt{\sigma_{GC}^2 + \sigma_{sky}^2 + \sigma_{RON}^2 + \sigma_{DN}^2}} \\
  &= \frac{N_{GC}}
  {\sqrt{N_{GC} + N_{sky} + n_p \cdot RON^2 + n_p \cdot DN \cdot t}}
\end{align*}
\hfill(3.0 points)\\
where,
\[
\begin{cases}
N_{GC} \rightarrow~\text{Total Source Count} \\
\sigma_{GC} = \sqrt{N_{GC}} \rightarrow~\text{Poisson Noise} \\
N_{sky} \rightarrow~\text{Total Sky Count} \\
\sigma_{sky} = \sqrt{N_{sky}} \rightarrow~\text{Sky Noise} \\
\sigma_{RON} = \sqrt{n_p \cdot 1 \cdot RON^2}
  \rightarrow~\text{Readout Noise} \\
\sigma_{DN} = \sqrt{n_p \cdot DN \cdot t} \rightarrow~\text{Dark Noise}
\end{cases}
\]

The noise values associated with the CCD operation can be easily calculated:
\[
\begin{cases}
\sigma_{RON}^2 = 15400 \times 5^2 = 385000~\text{counts} \\
\sigma_{DN}^2 = 15400 \times 6 \times \dfrac{1}{60} \times 15
  = 23100~\text{counts}
\end{cases}
\]
\hfill(2.0 points)

Now, the apparent visual magnitude of the globular cluster must be
calculated,
\[ V_{GC} - m_V = -2.5\log\left(\frac{\Omega_{GC}}{1}\right) \]
with $\Omega_{GC}$ being the solid angle subtended by the globular cluster,
\[ \Omega_{GC} = \pi \cdot \left(\frac{\theta}{2}\right)^2
   = \pi \times \left(\frac{72.0''}{2}\right)^2
   = 4071.5\,\mathrm{arcsec}^2 \]
\hfill(1.0 points)

Hence,
\[ V_{GC} = 20.6 - 2.5\log(4071.5) = 11.6\,\mathrm{mag} \]
\hfill(1.0 points)\\
and the associated flux of photons can now be evaluated using the
0-magnitude as a reference:
\begin{align*}
V_{GC} - V_0 &= -2.5\log\left(\frac{\phi_{GC}}{\phi_0}\right) \\
\phi_{GC} &= \phi_0 \cdot 10^{-V_{GC}/2.5} \\
  &\approx 0.234\,\mathrm{counts/nm/cm^2/s}
\end{align*}
\hfill(1.0 points)

With the calculated parameters, the total source count is:
\begin{align*}
N_{GC} &= \eta \cdot \phi_{GC} \cdot \frac{\pi}{4} \cdot D^2
  \cdot \Delta\lambda \cdot t \\
N_{GC} &= 0.80 \times 0.234 \times \frac{\pi}{4} \times (30.5)^2
  \times 88.0 \times 15 \\
  &= 180\,540~\text{counts}
\end{align*}
\hfill(1.0 points)

By referring to the signal-to-noise ratio expression and substituting the
values,
\begin{align*}
225 &= \frac{180540}{\sqrt{180540 + N_{sky} + 385000 + 23100}} \\
N_{sky} &= \left(\frac{180540}{225}\right)^2 - 385000 - 180540 - 23100 \\
N_{sky} &= 55250~\text{counts}
\end{align*}
\hfill(1.0 points)

The inverse procedure to determine the Total Source Count is taken:
\begin{align*}
\phi_{sky} &= \frac{N_{sky}}
  {\eta \cdot \frac{\pi}{4} \cdot D^2 \cdot \Delta\lambda \cdot t} \\
  &= \frac{55250}
  {0.80 \times \frac{\pi}{4} \times 30.5^2 \times 88.0 \times 15} \\
\phi_{sky} &= 7.16 \times 10^{-2}\,\mathrm{counts/nm/cm^2/s}
\end{align*}
\hfill(1.0 points)\\
which implies
\[ V_{sky} = V_0 - 2.5\log\left(\frac{\phi_{sky}}{\phi_0}\right)
   = 12.9\,\mathrm{mag} \]
\hfill(1.0 points)\\
and finally, since $\Omega_{sky} = \Omega_{GC}$
\begin{align*}
m_{sky} &= V_{sky} + 2.5\log\left(\frac{\Omega_{sky}}{1}\right) \\
  &= 12.8 + 2.5\log(4071.5) \\
% sic: the source prints 12.8 here although V_sky = 12.9 mag above.
m_{sky} &= 21.9\,\mathrm{mag/arcsec^2}
\end{align*}
\hfill(1.0 points)

Considering the accepted range of pixels from the previous item, the
accepted answers for the brightness of the sky fall within the interval
\[ 21.7 \leqslant m_{sky} \leqslant 22.1\,\mathrm{mag/arcsec^2} \]
\end{description}

\solhead{5.25}{Daksha Mission}
% source id: 2025_TH_T01
\begin{description}
\item[(T01.1)]
Since the two detectors make an angle with each other, photons from the
source will not be incident on the detectors at the same angle. Using this
the source can localized.
% sic: "can localized" as printed in the source.

Let $\widehat{n_1}$ be the unit vector perpendicular to the plane of
D$_1$ and $\widehat{n_2}$ be the unit vector perpendicular to the plane of
D$_2$. We have,
\begin{align*}
\widehat{n_1} &= \frac{1}{2}\widehat{x} + \frac{\sqrt{3}}{2}\widehat{y}
  && \text{(0.5 marks)} \\
\widehat{n_2} &= -\frac{1}{2}\widehat{x} + \frac{\sqrt{3}}{2}\widehat{y}
  && \text{(0.5 marks)}
\end{align*}
Let $\widehat{s} = s_x\widehat{x} + s_y\widehat{y}$ be the unit vector in
the direction of the source and $F$ be the flux of the source. The
effective area of capture for a detector will be
\[ A_{\mathrm{eff}} = A(|\widehat{s} \cdot \widehat{n}|) \]
\hfill(0.5 marks)
\ioaafig{2025_TH_T01-s1.pdf}

Hence, the power received is,
\[ P = AF(|\widehat{s} \cdot \widehat{n}|) \]
\hfill(1.0 marks)\\
and therefore,
\begin{align*}
P_1 &= AF(|\widehat{s} \cdot \widehat{n_1}|) \\
P_2 &= AF(|\widehat{s} \cdot \widehat{n_2}|)
\end{align*}
Hence,
\begin{align*}
\widehat{s} \cdot \widehat{n_1}
  &= \frac{1}{2}s_x + \frac{\sqrt{3}}{2}s_y
   = \frac{2.70 \times 10^{-10}}{AF} \\
\widehat{s} \cdot \widehat{n_2}
  &= -\frac{1}{2}s_x + \frac{\sqrt{3}}{2}s_y
   = \frac{4.70 \times 10^{-10}}{AF}
\end{align*}
Solving,
\begin{align*}
s_x &= -\frac{2.00 \times 10^{-10}}{AF}
  && \text{(0.5 marks)} \\
s_y &= \frac{7.40 \times 10^{-10}}{\sqrt{3}AF}
  && \text{(0.5 marks)}
\end{align*}
Hence, the angle the direction vector makes with respect to the positive
$y$-axis is
\[ \eta = \tan^{-1}\left(-\frac{s_x}{s_y}\right) \]
\hfill(1.0 marks)

Then, $\eta = 25.1^\circ = 0.438$\,rad. \hfill(0.5 marks)

\item[(T01.2)]
The time delay occurs because the satellites S$_1$ and S$_2$ are located at
different distances from the source. Light takes longer time to reach the
satellite that is farther (here, say S$_1$). Let $\Delta t = t_1 - t_2$.

\ioaafig{2025_TH_T01-s2.pdf}

If the direction of the source makes an angle $\theta$ (see the figure
above) with the line joining telescopes S$_1$ and S$_2$ we have the path
difference ($\Delta d$) given as,
\[ \Delta d = c\Delta t \]
\hfill(0.5 marks)\\
From the figure above the path difference can be written as,
\[ \Delta d = 2r\cos\theta \]
\hfill(1.0 marks)\\
Therefore,
\[ \cos\theta = \frac{c\Delta t}{2r} \]
Since the measurement of $\Delta t$ has some uncertainty, that is
$\Delta t \in [9.9, 10.1]$\,ms, we get a range $[\theta_1, \theta_2]$ for the
possible values of $\theta$.
\begin{align*}
\theta_1 &= \cos^{-1}\left(\frac{c(10.1\,\mathrm{ms})}{2r}\right)
  = 77.51^\circ
  && \text{(0.5 marks)} \\
\theta_2 &= \cos^{-1}\left(\frac{c(9.9\,\mathrm{ms})}{2r}\right)
  = 77.76^\circ
  && \text{(0.5 marks)}
\end{align*}

\ioaafig{2025_TH_T01-s3.pdf}

The source may lie anywhere in the annulus bounded by the red and blue
circles on the unit sphere above, whose area is given by
\[ 2\pi(\cos\theta_1 - \cos\theta_2) \]
\hfill(1.5 marks)\\
Hence the fraction of sky where the source might lie is
\[ f = \frac{1}{2}(\cos\theta_1 - \cos\theta_2) \]
\hfill(0.5 marks)
\[ f = 2.13 \times 10^{-3} \]
\hfill(0.5 marks)
\end{description}

\solhead{5.26}{Atmospheric Seeing}
% source id: 2025_TH_T09
\begin{description}
\item[(T09.1)]
Angular radius of the diffraction disk is given by
$\dfrac{1.22\lambda}{D}$. The radius of the image is therefore,
\[ d_{\mathrm{image}} = f\left(\frac{1.22\lambda}{D}\right) \]
% sic: the source says "The radius of the image is therefore" although
% d_image is the diameter asked for and the boxed value is the diameter.
\hfill(0.5 marks)
\[ d_{\mathrm{image}} = 8.95 \times 10^{-6}\,\mathrm{m} \]
\hfill(0.5 marks)

\item[(T09.2a)]
\ioaafig{2025_TH_T09-s1.pdf}
The ray-diagram at the boundary shows the incident and refracted rays with
the normal to the boundary: the angle of incidence is $\theta$, the
refracted ray in $n_1$ makes angle $(\theta - \alpha)$ with the normal, and
$\alpha$ is the angular shift of the refracted ray from the incident
direction. \hfill(2.0 marks)

\item[(T09.2b)]
Using Snell's Law,
\begin{gather*}
n_2\sin\theta = n_1\sin(\theta - \alpha) \\
\frac{n_2}{n_1}\sin\theta = \sin\theta\cos\alpha - \cos\theta\sin\alpha
\end{gather*}
\hfill(1.0 marks)\\
Using small angle approximation ($\sin\alpha \approx \alpha$ and
$\cos\alpha \approx 1$)
\[ \frac{n_2}{n_1}\sin\theta = \sin\theta - \alpha\cos\theta \]
\[ \alpha = \frac{\sin\theta}{\cos\theta}\left(1 - \frac{n_2}{n_1}\right) \]
\hfill(1.0 marks)

\item[(T09.2c)]
The shift in position of image due to change in $\theta$ is:
\begin{align*}
\Delta x_\theta &= f\alpha
  = f\Delta\tan\theta\left(1 - \frac{n_2}{n_1}\right)
  && \text{(1.0 marks)} \\
  &= f\left(1 - \frac{n_2}{n_1}\right)
     (\tan(\theta + \Delta\theta) - \tan\theta) \\
  &= f\left(1 - \frac{n_2}{n_1}\right)\sec^2\theta\,\Delta\theta
  && \text{(1.0 marks)} \\
  &= f \times (6.98 \times 10^{-8}\,\mathrm{rad})
\end{align*}
\[ \Delta x_\theta = 1.40 \times 10^{-7}\,\mathrm{m} \]
\hfill(1.0 marks)

\item[(T09.2d)]
Shift in the position of image due to the change in $n_2$ is:
\begin{align*}
\Delta x_n &= f\tan\theta
  \left(-\frac{n_2 + \Delta n_2}{n_1} + \frac{n_2}{n_1}\right) \\
  &= -f\tan\theta\left(\frac{0.0001}{100} \times \frac{n_2}{n_1}\right)
  && \text{(2.0 marks)} \\
  &= -f \times (5.77 \times 10^{-7}\,\mathrm{rad})
\end{align*}
\[ \Delta x_n = -1.15 \times 10^{-6}\,\mathrm{m} \]
\hfill(1.0 marks)

\item[(T09.3)]
Blue light has larger refractive index than the red light, hence it will
bend more. So incoming light will vary in colour from red to blue towards
right.

So, the correct answer will be F. \hfill(2.0 marks)

Note: convex lens inverts image, but does not affect refractive index
dependent bending.

\item[(T09.4a)]
\ioaafig{2025_TH_T09-s2.pdf}
The paths of the incident light rays are drawn below the zigzag boundary:
each planar segment produces a set of parallel refracted rays, and just
below each point where two planes make an upward bend the rays diverge,
leaving a region marked "X" on the plane of the telescope objective where no
light rays will reach. \hfill(4.0 marks)

\item[(T09.4b)]
Each set of planes, with specific tilt, will lead to one set of stripes of
parallel rays. From solution of part (T09.2b),
\[ \alpha = \tan\theta\,\frac{(n_1 - n_2)}{n_1} \]
Two sets of planes with relative angle of $2\theta$ will create angular
separation of $2\alpha$ between the stripes. \hfill(1.0 marks)\\
At a distance of $H$ this will create a horizontal separation of
$\approx 2\alpha H$. For $H = 1$\,km, $\theta = 10^\circ$, the horizontal
shift will be
\[ \Delta x \simeq 2\alpha H = 2\tan\theta\,\frac{(n_1 - n_2)}{n_1}H
   = 2 \times 1.76 \times 10^{-6} \times 1 \times 10^{5}\,\mathrm{cm}
   = 0.35\,\mathrm{cm}. \]
\hfill(1.0 marks)\\
Just below the point where two planes make upward bend, a patch of this
width 0.35\,cm will be centred where no light from the star will reach.
Hence, $W_{\mathrm{X}} = 0.35$\,cm. \hfill(1.0 marks)

\item[(T09.4c)]
From the figure one can see that two adjacent planes (with downward bend),
will lead to overlapping stripes of rays producing two images.

Therefore, the region in which only single image is formed will be 10.0\,cm
(including the region X). Hence, $D_{\mathrm{max}} = 10.0$\,cm.
\hfill(4.0 marks)

\item[(T09.5)]
For $D \gg d$, several stripes will fall on the lens, leading to multiple
images. Here, four images will form, two each from the set of planes with
zigzag shapes along $x$-axis and $y$-axis respectively.

First consider the zigzag shape of the boundary along the $x$ axis. From
solution of (T09.4), angle between these two sets of rays will be
$3.52 \times 10^{-6}$\,rad. With $f = 2$\,m, the separation between the two
images will be $= 7.0\,\mu\mathrm{m}$.

For $D = 100$\,cm, the diffraction peak size $= 1.34\,\mu\mathrm{m}$. So
the images formed by the two sets of planes (with tilts of $\pm 10^\circ$
from horizontal for each set) will form two spots, clearly separated along
the $x$-axis. Similarly for the $y$-axis.

The resulting speckles will be as shown below. \hfill(6.0 marks)
\ioaafig{2025_TH_T09-s3.pdf}

\item[(T09.6)]
A single image is formed when only one stripe of rays is allowed to fall on
the objective. This is only possible when the objective is placed suitably
below the zigzag pattern.

The diffraction-limited image size for a $D = 8$\,cm lens is
$8.4 \times 10^{-6}$\,m. \hfill(1.0 marks)

The angle of the atmospheric layer changes from $+10^\circ$ to $-10^\circ$\\
$\implies$ Change in angle of rays is $3.52 \times 10^{-6}$\,rad\\
$\implies$ Shift in image position $= 3.52\,\mu\mathrm{m}$.
\hfill(1.0 marks)

Now, the shift in the image position allowed is 1\,\% of the image size.\\
$\therefore$ Allowed shift $= 0.084\,\mu\mathrm{m}$ \hfill(1.0 marks)

Longest exposure time $=$ Time taken by the image to shift by
$0.084\,\mu\mathrm{m}$

Image shifts by $3.52\,\mu\mathrm{m}$ in 1\,s $\implies$
$0.084\,\mu\mathrm{m}$ shift occurs in 0.024\,s.
\[ t_{\mathrm{max}} = 0.024\,\mathrm{s} \]
\hfill(2.0 marks)
\end{description}

\solhead{5.27}{CCD Image Processing}
% source id: 2009_DA_D1
\begin{description}
\item[(a)] To measure instrumental magnitude we should choose an aperture.
Careful investigation of the image shows that a $5 \times 5$ pixel aperture
is enough to measure $m_I$ for all stars. $m_I$ can be calculated using:
\[ m_I = -2.5\log\!\left(\frac{\sum_{i=1}^{N} I_{i(\mathrm{star})} - N\bar{I}_{\mathrm{Sky}}}{\mathrm{Exp}}\right) \]
where $I_{i(\mathrm{star})}$ is the pixel value for each pixel inside the
aperture, $N$ is number of pixels inside the aperture,
$\bar{I}_{\mathrm{Sky}}$ is the average of sky value per pixel taken from
dark part of image and Exp is the exposure time. Table (1.4) lists values
for $m_I$ and $Zmag$ calculated for all three identified stars.
\[ \bar{I}_{\mathrm{Sky}} = 4.42 \]
\[ N = 25 \]
\[ Exp = 450 \]

\begin{center}
Table (1.4)\\[2pt]
\begin{tabular}{|c|c|c|c|}
\hline
Star & $m_I$ & $m_t$ & $Zmag$ \\
\hline
1 & $-3.02$ & 9.03 & 12.38 \\
\hline
3 & $-5.85$ & 6.22 & 12.40 \\
\hline
4 & $-4.04$ & 8.02 & 12.39 \\
\hline
\end{tabular}
\end{center}

\item[(b)] Average $Zmag = 12.4$.

\item[(c)] Following part (a) for stars 2 and 5, we can calculate true
magnitudes ($m_t$) for these stars.

\begin{center}
Table (1.5)\\[2pt]
\begin{tabular}{|c|c|c|}
\hline
Star & $m_I$ & $m_t$ \\
\hline
2 & $-2.13$ & 9.93 \\
\hline
5 & $-0.66$ & 11.4 \\
\hline
\end{tabular}
\end{center}

\item[(d)] Pixel scale for this CCD is calculated as
\[ p = \frac{\text{pixel size}}{\text{focal length}} \times \frac{180 \times 3600}{\pi} = 4.30'' \]

\item[(e)] Average sky brightness:
\[ m_{sky} = -2.5\log\frac{\bar{I}_{\mathrm{Sky}}}{(Exp)(p)^2} + Zmag = 20.6 \]

\item[(f)] To estimate astronomical seeing, first we plot pixel values in
$x$ or $y$ direction for one of the bright stars in the image. As plot (1)
shows, the FWHM of pixel values which is plotted for star 3, is 1 pixel,
hence astronomical seeing is equal to
\[ seeing \cong 4'' \]

\ioaafig{2009_DA_D1-s1.pdf}
\end{description}

\solhead{5.28}{CCD image}
% source id: 2010_DA_D1
\begin{enumerate}
\item \hfill(0.5 point for one star, totally 2 points)
\item The field of view of the camera is marked on the map. \hfill(2 p)

\ioaafig{2010_DA_D1-s1.pdf}

\item According to the pic of A2, it's easy to find the field of view of the
telescope. It's about $26'$, and the declination of the center of the CCD
image is $7.3^\circ$. Thus the side length of the field of view is:
\[ 26' \times \cos 7.3^\circ = 1550'' \]
Image scale is $d = f\theta$, so, $s = f/206.265$\,mm/arcsec
$= 0.0154$\,mm/arcsec.\\
The chip size is $1550 \times 0.0154 = 24$\,mm. \hfill(4p)

\item The star is 10 pixels across, so the FWHM is $10/3.5 = 2.9$ pixels.
\hfill(4p)\\
Seeing is $S = 2.9\,\mathrm{pixels} \times 1.5''/\mathrm{pixel}$ (from Q3
and 1024 pixels) $= 4.4''$.

\item Theoretical (Airy) diffraction disc is $2.44\lambda/D$ radians in
diameter:
\[ A = 2.44 \times 550 \times 10^{-9}/0.61\ \mathrm{rad} = 0.45'' \sim 0.3\ \text{pixels} \]
$A \ll S$ (seeing). \emph{(Accept all reasonable wavelengths:
450--650\,nm)} \hfill(4p)

\item Seeing $=$ FWHM $\times\ 1.5''/\mathrm{pixel}$ (from Q3) $= 1''$. So,
FWHM $= 1/1.5$ pixel $= 1$ pixel.\\
Printed image of star would then be
$s_2 = 3.5 \times \mathrm{FWHM} = 3.5$ pixels. \hfill(3p)\\
Note: if use: $s_2 = 1'' \times 10\,\mathrm{pix}/4.4'' = 2.3$\,pix, 2 points.

\item For object 1, the trail of the object is about $107''$ (measured from
pic 1, 300\,s exposure). Its angular velocity is:
\[ \omega_1 = 107''/300\,\mathrm{s} = 0.36\,''/\mathrm{s} \]
Note: accept to $v \pm 10\%$. \hfill(3p)

For object 2, it moves about 8 pixels between pic 2A and 2B. 8 pixels
$\sim 12''$, and the time between exposures is $17^{\mathrm{m}}27^{\mathrm{s}}$.
Its angular velocity is: \hfill(3p)
\[ \omega_2 = 12''/1047\,\mathrm{s} = 0.012\,''/\mathrm{s}\ (\text{accept} \pm 10\%). \]
\begin{description}
\item[a)] \textbf{wrong}: different masses of the objects, \hfill($+2$/$-1$p)
\item[b)] \textbf{right}: different distances of the objects from Earth,
  \hfill($+3$/$-1$p)
\item[c)] \textbf{right}: different orbital velocities of the objects,
  \hfill($+3$/$-1$p)
\item[d)] \textbf{wrong}: different projections of the objects' velocities,
  \hfill($+2$/$-1$p)
\item[e)] \textbf{rejected}: Object 1 orbits the Earth while Object 2 orbits
  the Sun. \hfill(0p)
\end{description}
\end{enumerate}

\solhead{5.29}{}
% source id: 2013_DA_D1
Calculate the angular distance, $\varphi$, (in degrees or minutes of arc),
of two bright stars, whose coordinates are given in the \emph{List of Bright
stars}. Preferably these stars should be chosen to be far apart (e.g.
$\alpha$ UMa [$a_1 = 11^{\mathrm{h}}03^{\mathrm{m}}$,
$\delta_1 = +61^\circ 45'$] and $\eta$ UMa
[$a_2 = 13^{\mathrm{h}}48^{\mathrm{m}}$, $\delta_2 = +49^\circ 19'$]).

\ioaafig{2013_DA_D1-s1.pdf}

\begin{enumerate}
\item In order to calculate the angular distance of the two stars, the
coordinates should be converted to decimal degrees (e.g. $\alpha$ UMa
[$a_1 = 165^\circ.75$, $\delta_1 = +61.75^\circ$] and $\eta$ UMa
[$a_2 = 207.0^\circ$, $\delta_2 = +49^\circ.32$]) \hfill(2.5 Point)

\item Use the \emph{cosine law} to calculate the angular distance between
the two stars:
\[ \varphi = \arccos(\sin\delta_1 \times \sin\delta_2
   + \cos\delta_1 \times \cos\delta_2 \times \cos(a_1 - a_2))
   \quad\rightarrow\quad \varphi = 25^\circ.8448 \]
\hfill(10 Points)

\item Measure the distance, $d_0$, of the two stars in mm in Figure 1.
$d_0 = 138$\,mm. \hfill(2.5 Points)

\item The photograph in Figure 1 does not have the same dimensions as the
original photograph. Measure the length, $\ell$, of the photograph in mm.
$\ell = 140$\,mm. This length corresponds to 22\,mm. Convert $d_0$ to the
distance, $d$, of the original photograph,
$d = 138 \times \dfrac{22}{140}$\,mm. $d = 21.6857$\,mm \hfill(1.5 Points)

\item % sic: the source numbers this step "4." again, duplicating the
% previous item's number
The \emph{image scale} $\left(\dfrac{\varphi}{d}\right)$ of the original
photograph is
$\dfrac{\varphi}{d} = \dfrac{25^\circ.8448}{21.6857}
 = 1.1918\,\dfrac{\mathrm{deg}}{\mathrm{mm}}$ \hfill(2.5 Points)

\item The focal length is given from equation (see Figure 2):
\[ \tan\frac{\varphi}{2} = \frac{d}{2f}
   \quad\text{or}\quad
   f = \frac{d}{2\tan\dfrac{\varphi}{2}} \]
from which the focal length is calculated:
\[ f = \frac{21.6857}{2 \times 0.2294}
   \quad\text{or}\quad
   \mathbf{f = 47.3\ \mathrm{mm}} \]
\hfill(1.5 Points)\\
(If errors are included \hfill(4 Points)
% sic: the source sentence "(If errors are included" is left unfinished
\end{enumerate}

\solhead{5.30}{Photometric comparison of surveys}
% source id: 2024_DA_D1
\begin{description}
\item[(a)] DES is Table 1 and SDSS Table 2, since faint stars will have
higher uncertainties. SDSS uses a smaller telescope (2.5\,m) than DES
(4\,m) hence, it is typically shallower and has larger errors. Just looking
at the error distribution should be enough to tell who has larger errors.

\item[(b)] The linear regression of each curve:

\ioaafig{2024_DA_D1-s1.pdf}

The reported parameters are:

\textbf{SDSS:}
\[ \log_{10}(err\ m_g) = A \cdot m_g + B \]
$A = 0.2780$ (acceptable from 0.2502 to 0.3058 which is within 10\% range
of the fit)\\
$B = -7.3108$ (acceptable from $-8.0419$ to $-6.5797$ which is within 10\%
range of the fit)

\textbf{DES:}
\[ \log_{10}(err\ m_g) = A \cdot m_g + B \]
$A = 0.4063$ (acceptable from 0.3657 to 0.4470 which is within 10\% range
of the fit)\\
$B = -10.7597$ (acceptable from $-10.2217$ to $-11.2977$ which is within
10\% range of the fit)

\item[(c)] Using that $S/N \sim 1/err\ m_g$ and above fits we can arrive at
the following:\\
\textbf{SDSS:} $\log_{10}(err\ m_g) = 0.2780 \cdot 21.5 - 7.3108
\rightarrow err\ m_g = 0.0464$\\
\textbf{DES:} $\log_{10}(err\ m_g) = 0.4063 \cdot 21.5 - 10.7597
\rightarrow err\ m_g = 0.0095$\\
Accepted answers follow below.

\begin{center}
\begin{tabular}{|l|c|l|}
\hline
 & err & $S/N$ \\
\hline
SDSS err at 21.5 mag & 0.05 & \textbf{22} (acceptable from 19.8 to 24.2
$\rightarrow$ within 10\% range of the result) \\
\hline
DES err at 21.5 mag & 0.01 & \textbf{106} (acceptable from 95.4 to 116.6
$\rightarrow$ within 10\% range of the result) \\
\hline
\end{tabular}
\end{center}

\item[(d)] One tip for this question is that the student should realize
that the SDSS RA coordinate is sorted, so it can be used as a reference for
the scanning of DES coordinates to perform the match. These are the stars
that can be matched between catalogs:

\begin{center}
\begin{tabular}{ccccc|ccccc|c}
$ID_1$ & RA & Dec & $m_g$ & err $m_g$ &
$ID_2$ & RA & Dec & $m_g$ & err $m_g$ & Sep \\
 & (deg) & (deg) & (mag) & (mag) &
 & (deg) & (deg) & (mag) & (mag) & (arcsec) \\
\hline
3 & 0.064911 & 0.026395 & 21.3931 & 0.0084 & 8 & 0.064910 & 0.026419 & 21.6428 & 0.0488 & 0.08475 \\
9 & 0.057422 & 0.076006 & 21.8897 & 0.0127 & 6 & 0.057414 & 0.075999 & 21.8884 & 0.0578 & 0.03724 \\
4 & 0.098343 & 0.054871 & 21.3934 & 0.0088 & 13 & 0.098343 & 0.054854 & 21.6542 & 0.0499 & 0.06186 \\
6 & 0.006188 & 0.066928 & 21.5490 & 0.0088 & 1 & 0.006167 & 0.066874 & 21.9020 & 0.0576 & 0.2076 \\
12 & 0.071089 & 0.091053 & 21.4390 & 0.0098 & 9 & 0.071102 & 0.091058 & 21.9259 & 0.0751 & 0.05009 \\
2 & 0.064741 & 0.021568 & 21.1111 & 0.0067 & 7 & 0.064745 & 0.021583 & 21.3634 & 0.0422 & 0.05655 \\
8 & 0.076715 & 0.079496 & 21.4808 & 0.0089 & 11 & 0.076709 & 0.079747 & 21.5303 & 0.0476 & 0.08276 \\
% sic: the solution table gives Dec 0.079747 for ID2 = 11, whereas Table 2
% of the problem lists 0.079474
\end{tabular}
\end{center}

\item[(e)] The figure of g magnitude of DES vs SDSS should look like this
with the best linear fit:

\ioaafig{2024_DA_D1-s2.pdf}
\ioaafig{2024_DA_D1-s3.pdf}

The student should identify the stars that are closest to the one-to-one
line ($x = y$) between surveys. This could be with a zero offset, or a
non-zero offset. The \textbf{two} stars that would be suitable for
photometric calibration with zero offset are the ones with ID's \textbf{8}
and \textbf{9} of Table 1 (DES). The \textbf{four} stars that would be
suitable for photometric calibration with non-zero offset are the ones with
ID's \textbf{2, 3, 4 and 6} of Table 1 (DES).
\end{description}

\solhead{5.31}{Binoculars: Hyades limiting magnitude, M31 shape}
% source id: 2007_OBS-SKY_2
\textbf{(2.1)} The correct value of brightness of $\gamma$-Tauri $= 3.6$.

If the answer is 3.5, 3.6 or 3.7 \hfill(5 points)\\
If the answer is 3.3, 3.4, 3.8 or 3.9 \hfill(3 points)\\
If the answer is none of the above \hfill(0 point)

\textbf{(2.2)} Correct shape and size (between 0.5 to 1 degree) --- 3
points. Correct orientation (within 23 degrees of the picture) --- 2
points.

\ioaafig{2007_OBS-SKY_2-s3.pdf}

\solhead{5.32}{CCD imaging and identification of a cluster field}
% source id: 2008_OBS-TEL_1
Model answer pages are provided for each of the five recommended regions:
the sky chart image, an example CCD image with the identified stars
numbered, and the identification table.

\ioaafig{2008_OBS-TEL_1-s1.pdf}
\ioaafig{2008_OBS-TEL_1-s2.pdf}
\ioaafig{2008_OBS-TEL_1-s3.pdf}
\ioaafig{2008_OBS-TEL_1-s4.pdf}
\ioaafig{2008_OBS-TEL_1-s5.pdf}
\ioaafig{2008_OBS-TEL_1-s6.pdf}
\ioaafig{2008_OBS-TEL_1-s7.pdf}
\ioaafig{2008_OBS-TEL_1-s8.pdf}
\ioaafig{2008_OBS-TEL_1-s9.pdf}
\ioaafig{2008_OBS-TEL_1-s10.pdf}

\textbf{M7} ($\alpha = 17^{\mathrm{h}}53.^{\mathrm{m}}3$,
$\delta = -34^\circ 46.'2$). Central equatorial coordinates:
$\alpha = 17^{\mathrm{h}}54^{\mathrm{m}}22.3^{\mathrm{s}}$,
$\delta = -34^\circ 47' 15''$. Exposure time: 5\,sec.

\begin{center}
\begin{tabular}{clllc}
Number & Names or Catalogue number & R.A. & Declination & Magnitude \\
\hline
1 & SAO 209436 & $17^{\mathrm{h}}\,54^{\mathrm{m}}\,14.380^{\mathrm{s}}$ & $-34^\circ 43' 39.535''$ & 6.88 \\
2 & GSC 7386:494 & $17^{\mathrm{h}}\,54^{\mathrm{m}}\,09^{\mathrm{s}}$ & $-34^\circ 44' 31''$ & 11.6 \\
3 & GSC 7386:601 & $17^{\mathrm{h}}\,54^{\mathrm{m}}\,06^{\mathrm{s}}$ & $-34^\circ 44' 44''$ & 12.2 \\
4 & GSC 7386:561 & $17^{\mathrm{h}}\,54^{\mathrm{m}}\,04^{\mathrm{s}}$ & $-34^\circ 44' 06''$ & 8.5 \\
5 & GSC 7386:15 & $17^{\mathrm{h}}\,53^{\mathrm{m}}\,55^{\mathrm{s}}$ & $-34^\circ 48' 20''$ & 11.9 \\
6 & GSC 7386:40 & $17^{\mathrm{h}}\,54^{\mathrm{m}}\,13^{\mathrm{s}}$ & $-34^\circ 48' 18''$ & 11.8 \\
\end{tabular}
\end{center}

\textbf{M8} ($\alpha = 18^{\mathrm{h}}04.^{\mathrm{m}}2$,
$\delta = -24^\circ 22.'1$). Central equatorial coordinates:
$\alpha = 18^{\mathrm{h}}04^{\mathrm{m}}51^{\mathrm{s}}.2$,
$\delta = -24^\circ 23' 04''$. Exposure time: 10\,sec.

\begin{center}
\begin{tabular}{clllc}
Number & Names or Catalogue number & R.A. & Declination & Magnitude \\
\hline
1 & GSC 6842:1751 & $18^{\mathrm{h}}\,04^{\mathrm{m}}\,10^{\mathrm{s}}$ & $-24^\circ 20' 53''$ & 12.5 \\
2 & GSC 6846:166 & $18^{\mathrm{h}}\,04^{\mathrm{m}}\,05^{\mathrm{s}}$ & $-24^\circ 23' 01''$ & 12.1 \\
3 & GSC 6846:211 & $18^{\mathrm{h}}\,03^{\mathrm{m}}\,58^{\mathrm{s}}$ & $-24^\circ 25' 35''$ & 12.0 \\
4 & GSC 6846:718 & $18^{\mathrm{h}}\,04^{\mathrm{m}}\,27^{\mathrm{s}}$ & $-24^\circ 23' 52''$ & 12.4 \\
5 & GSC 6842:1578 & $18^{\mathrm{h}}\,04^{\mathrm{m}}\,23^{\mathrm{s}}$ & $-24^\circ 21' 25''$ & 13.4 \\
\end{tabular}
\end{center}

\textbf{M20} ($\alpha = 18^{\mathrm{h}}02.^{\mathrm{m}}4$,
$\delta = -22^\circ 58.'7$). Central equatorial coordinates:
$\alpha = 18^{\mathrm{h}}02^{\mathrm{m}}56^{\mathrm{s}}.7$,
$\delta = -23^\circ 01' 56''$. Exposure time: 10\,sec.

\begin{center}
\begin{tabular}{clllc}
Number & Names or Catalogue number & R.A. & Declination & Magnitude \\
\hline
1 & GSC 6842:631 & $18^{\mathrm{h}}\,02^{\mathrm{m}}\,18^{\mathrm{s}}$ & $-23^\circ 00' 25''$ & 12.4 \\
2 & GSC 6842:240 & $18^{\mathrm{h}}\,02^{\mathrm{m}}\,19^{\mathrm{s}}$ & $-23^\circ 00' 52''$ & 12.8 \\
\end{tabular}
\end{center}

\textbf{M21} ($\alpha = 18^{\mathrm{h}}04.^{\mathrm{m}}2$,
$\delta = -22^\circ 29.'5$). Central equatorial coordinates:
$\alpha = 18^{\mathrm{h}}04^{\mathrm{m}}46^{\mathrm{s}}.3$,
$\delta = -22^\circ 29' 27''$. Exposure time: 10\,sec.

\begin{center}
\begin{tabular}{clllc}
Number & Names or Catalogue number & R.A. & Declination & Magnitude \\
\hline
1 & GSC 6842:985 & $18^{\mathrm{h}}\,04^{\mathrm{m}}\,09^{\mathrm{s}}$ & $-22^\circ 30' 28''$ & 12.9 \\
2 & GSC 6263:306 & $18^{\mathrm{h}}\,04^{\mathrm{m}}\,01^{\mathrm{s}}$ & $-22^\circ 29' 36''$ & 13.1 \\
\end{tabular}
\end{center}

\textbf{M23} ($\alpha = 17^{\mathrm{h}}57.^{\mathrm{m}}0$,
$\delta = -18^\circ 59.'4$). Central equatorial coordinates:
$\alpha = 17^{\mathrm{h}}57^{\mathrm{m}}32^{\mathrm{s}}.1$,
$\delta = -19^\circ 01' 08''$. Exposure time: 10\,sec.

\begin{center}
\begin{tabular}{clllc}
Number & Names or Catalogue number & R.A. & Declination & Magnitude \\
\hline
1 & GSC 6258:809 & $17^{\mathrm{h}}\,56^{\mathrm{m}}\,58^{\mathrm{s}}$ & $-18^\circ 58' 21''$ & 13.2 \\
2 & GSC 6258:825 & $17^{\mathrm{h}}\,56^{\mathrm{m}}\,59^{\mathrm{s}}$ & $-18^\circ 58' 43''$ & 12.5 \\
\end{tabular}
\end{center}

Marking: student who has the highest number of stars identified ($N_m$)
gets 60\%; other students who get $N$ stars: $N \times 0.60/N_m$.
Estimation of the limiting magnitude: 20\% ($\Delta < 1$ mag: full mark,
$1 < \Delta < 2$ mag: half mark, $\Delta > 3$ mag: 0). Good image quality
(unsaturated, not under-exposed): 20\%.

\solhead{5.33}{Telescope pointing and field identification}
% source id: 2009_OBS-TEL_3
\begin{center}
\begin{tabular}{|l|c|c|c|}
\hline
\multicolumn{4}{|c|}{Coordinate of Selected Star} \\
\hline
Star & Hour Angle & R.A. & Dec. \\
\hline
Deneb (Alpha Cygni) & ST$-$RA & $20^{\mathrm{h}}\,41^{\mathrm{m}}$ & $+45^\circ$ \\
\hline
Alfirk (Beta Cephei) & ST$-$RA & $21^{\mathrm{h}}\,28^{\mathrm{m}}$ & $+71^\circ$ \\
\hline
Algol (Beta Persei) & RA$-$ST & $03^{\mathrm{h}}\,08^{\mathrm{m}}$ & $+41^\circ$ \\
\hline
Capella (Alpha Aurigae) & RA$-$ST & $05^{\mathrm{h}}\,18^{\mathrm{m}}$ & $+46^\circ$ \\
\hline
\end{tabular}
\end{center}

(The answer key leaves the ``Local Time:'' and ``Sidereal Time:'' fields
blank, to be filled at the moment of observation; the hour angle of each
star is expressed as ST$-$RA, or RA$-$ST when the star is east of the
meridian.)

\solhead{5.34}{Field of view of the telescope (NGC 6231)}
% source id: 2012_OBS-TEL_1
Answer: 24.0 arc minutes ($0.4^\circ$). We will tolerate up to 10\% error
as correct answer. Error between 10 to 20\% will be considered 70\%
correct.

\solhead{5.35}{Night observation round}
% source id: 2014_OBS-TEL_1
\textit{No official solution for the 2014 outdoor telescope round was published; the surviving 2014 keys cover only the planetarium and sky-chart rounds.}

\solhead{5.36}{Main telescope: point to M7 and observe}
% source id: 2015_OBS-TEL_1
\begin{enumerate}
\item Checklist (Telescope Assistant): Object Accuracy (Right Object $=$
Full Score).

\item Cardinal points: as marked on the solution chart (E at top, W at
bottom, N at left, S at right). Each correct cardinal point: quarter score.

\item Three missing stars: marked on the solution chart. Region 1 (Full
Score), 2 (Half Score), 3 (Quarter Score). Note: the regions are marked
with the circles (region will be chosen during the IBM).

\ioaafig{2015_OBS-TEL_1-s2.pdf}

\item Missing stars:
\begin{description}
\item[1.] 7.4
\item[2.] 6.4
\item[3.] 7.0
\end{description}
Score range: $\pm 0.2$ = Full Score; $\pm 0.5$ = Half Score;
$\pm 0.8$ = Quarter Score.

\item Using the 15\,mm (eyepiece), the field of view of the instrument is:
$0.75^\circ$
\begin{itemize}
\item Using stopwatch to measure the time for a star to travel from the
  center of the FOV to the edge of the FOV. \hfill(Half Score)
\item Using stopwatch to measure the time for a star to travel from the
  edge of the FOV to the other edge of the FOV. \hfill(Quarter Score)
\item Calculating the FOV,
\[ FOV(\mathrm{deg}) = \frac{\text{Travel time (sec)}}{23^{\mathrm{h}}\,56^{\mathrm{m}}}
   \times \cos(\delta) \times 360\,\mathrm{deg} \]
\hfill(Half Score)
\end{itemize}
\end{enumerate}

\solhead{5.37}{Telescope already set: identify the field}
% source id: 2016_OBS-TEL_OT1
\textit{No official solution was published for the 2016/2017 observation rounds.}

\solhead{5.38}{Telescope observation task}
% source id: 2016_OBS-TEL_OT2
\textit{No official solution was published for the 2016/2017 observation rounds.}

\solhead{5.39}{Telescope observation task}
% source id: 2016_OBS-TEL_OT3
\textit{No official solution was published for the 2016/2017 observation rounds.}

\solhead{5.40}{Night observation with the real sky}
% source id: 2017_OBS-TEL_1
\textit{No official solution was published for the 2016/2017 observation rounds.}

\solhead{5.41}{Eyepiece field of view}
% source id: 2018_OBS-TEL_O5
\[ F \tan(\mathrm{FOV}/2) = f \tan(45^\circ/2) \]
$F = 700$\,mm, $f = 26$\,mm
\[ \mathrm{FOV}/2 = \arctan(0.414 \cdot f/F) = 0.881^\circ \]
\[ \mathrm{FOV} = 2 \times 0.881^\circ = 1.76^\circ \approx 1.8^\circ \]
10 points for $1.6^\circ$--$1.9^\circ$, 5 points for $1.3^\circ$--$1.5^\circ$
and $2.0^\circ$--$2.2^\circ$.

\solhead{5.42}{Text observed through the telescope}
% source id: 2018_OBS-TEL_O6
2 planets: Saturn, Jupiter.\\
2 Messier objects: M15, M74.\\
1 equatorial coordinate: $+30^\circ$ $21^{\mathrm{h}}$.

8 points for each text.

\solhead{5.43}{Telescope task 1}
% source id: 2019_OBS-TEL_1
Evaluation of the alignment and marking is done by the telescope
assistants. \hfill Max.: 5 points

\solhead{5.44}{Telescope task 2}
% source id: 2019_OBS-TEL_2
\textbf{Saturn and Titan.} The official answer depends on the night of
observation; reference values are given for both nights. Titan's position is
checked with a transparent template foil: if Titan is within the template
circle, 2 points; if outside, 0 points.

\ioaafig{2019_OBS-TEL_2-s1.pdf}
\ioaafig{2019_OBS-TEL_2-s2.pdf}
\ioaafig{2019_OBS-TEL_2-s3.pdf}
\ioaafig{2019_OBS-TEL_2-s4.pdf}

\medskip
For 2019.08.04 22:00 UT:
\[ d_{\mathrm{Titan}} = 164'' \]
\begin{center}
\begin{tabular}{ll}
$d$: $148$--$180''$ & 3 points \\
$d$: $139$--$189''$ & 2 points \\
$d$: $123$--$205''$ & 1 point \\
\end{tabular}
\end{center}
\[ \mathrm{PA}_{\mathrm{Titan}} = 109^\circ \]
\begin{center}
\begin{tabular}{ll}
PA: $99$--$119^\circ$ & 4 points \\
PA: $89$--$129^\circ$ & 2 points \\
\end{tabular}
\end{center}

For 2019.08.05 22:00 UT:
\[ d_{\mathrm{Titan}} = 129'' \]
\begin{center}
\begin{tabular}{ll}
$d$: $116$--$131''$ & 3 points \\
$d$: $110$--$148''$ & 2 points \\
$d$: $97$--$161''$ & 1 point \\
\end{tabular}
\end{center}
\[ \mathrm{PA}_{\mathrm{Titan}} = 124^\circ \]
\begin{center}
\begin{tabular}{ll}
PA: $114$--$134^\circ$ & 4 points \\
PA: $104$--$144^\circ$ & 2 points \\
\end{tabular}
\end{center}

\medskip
Scoring:

\begin{center}
\begin{tabular}{lll}
Titan position + name: & & 2 points \\
Titan is outside the circle: & & 0 point \\[4pt]
Ring continuity in front of the disk is proper & & 2 points \\
Ring tilt towards Earth is proper & & 2 points \\
\multicolumn{3}{l}{\quad ring does not exceed northern pole, ring is not too thin} \\
Ring tilt towards sky E--W direction is proper & & 2 points \\[4pt]
PA measurement of Titan: & $\pm 10^\circ$ & 4 points \\
 & $\pm 20^\circ$ & 2 points \\[4pt]
Distance calculation of Titan & $\pm 10$\% & 3 points \\
 & $\pm 15$\% & 2 points \\
 & $\pm 25$\% & 1 point \\
\end{tabular}
\end{center}

The Saturn drawing evaluation scheme uses the sub-Earth latitude
$\Phi = 24.8^\circ$ (ring tilt towards Earth, max 2 points: drawings
equivalent to $22.8$--$26.8^\circ$ get 2p; $20.8^\circ$ gets 1p, S pole
visible; $28.8^\circ$ gets 1p, very thick; thinner or thicker gets 0p) and
the equatorial position angle $\mathrm{PA} = 6.4^\circ$ (ring tilt towards
the sky E--W direction, max 2 points: $6.4^\circ$ or $9.4^\circ$ get 2p;
$3.4^\circ$ or $12.4^\circ$ get 1p; $0^\circ$ or $15.4^\circ$ get 0p).
Inverse, improper, too wide or too narrow rings get 0 points for ring
continuity.

\ioaafig{2019_OBS-TEL_2-s5.pdf}
\ioaafig{2019_OBS-TEL_2-s6.pdf}

\medskip
\textbf{Alternative task: observing Epsilon Lyr.} \hfill TOTAL: 15 points

Checking the object in the FOV, and pointing is evaluated by the telescope
assistant. \hfill(0--3 points)

FOV with 10 mm eyepiece: the components of the close pairs are not resolved.

\ioaafig{2019_OBS-TEL_2-s7.pdf}

FOV drawing: star field and directions:
\begin{center}
\begin{tabular}{ll}
Correct direction \& labelling of North and East relative to the star field: & 2 points \\
Correct drawing of at least 3 stars in the FOV: & 1 point \\
Correct drawing of at least 6 stars in the FOV: & 2 points \\
\end{tabular}
\end{center}

Wide pair distance: $208'' = 3.47'$ \hfill 2 points

\begin{center}
\begin{tabular}{lll}
Distance estimation of wide pair: & $d_{\epsilon 1 - \epsilon 2} = 3.2$--$3.8'$ & 2 points \\
 & $d_{\epsilon 1 - \epsilon 2} = 2.8$--$4.2'$ & 1 point \\
\end{tabular}
\end{center}

Wide pair PA: $172^\circ$ \hfill 3 points

\begin{center}
\begin{tabular}{lll}
Position angle of the wide pair: & $\mathrm{PA}_{\epsilon 2} = 170$--$175^\circ$ & 3 points \\
 & $\mathrm{PA}_{\epsilon 2} = 167$--$177^\circ$ & 2 points \\
 & $\mathrm{PA}_{\epsilon 2} = 162$--$182^\circ$ & 1 point \\
\end{tabular}
\end{center}

Correct estimation of the relative angle of close pairs: \hfill 3 points

Relative angle between the direction of lines fitted onto the two close
pairs: $92^\circ$ (also the $88^\circ$ complementary angle is acceptable ---
referring to this, the evaluation bands are centered to $90^\circ$):
\begin{center}
\begin{tabular}{ll}
$85$--$95^\circ$ & 3 points \\
$80$--$100^\circ$ & 2 points \\
$75$--$105^\circ$ & 1 point \\
\end{tabular}
\end{center}

\solhead{5.45}{Telescope task 3}
% source id: 2019_OBS-TEL_3
Checking the object in the FOV, and pointing is evaluated by the telescope
assistant. \hfill(0--4 points)

The field of view (in the 25 mm eyepiece) drawn on the star chart is checked
against a transparent template foil showing the correct FOV circle around
M57, with dashed circles marking the $\pm 10$\% and $\pm 20$\% accuracy
bands.

\ioaafig{2019_OBS-TEL_3-s1.pdf}
\ioaafig{2019_OBS-TEL_3-s2.pdf}

\begin{center}
\begin{tabular}{ll}
Accuracy within $\pm 10$\%: & 6 points \\
Accuracy within $\pm 20$\%: & 3 points \\
\end{tabular}
\end{center}

\solhead{5.46}{Telescope task 4}
% source id: 2019_OBS-TEL_4
Checking the object in the FOV, and pointing is evaluated by the telescope
assistant. \hfill(0--8 points)

Brightness estimation:
\begin{center}
\begin{tabular}{ll}
Magnitude $\pm 0.2$ mag: & 6 points \\
Magnitude $\pm 0.3$ mag: & 4 points \\
Magnitude $\pm 0.4$ mag: & 2 points \\
\end{tabular}
\end{center}

Time in UTC (if corrected for Summer Time and Time Zone differences):
\hfill 1 point

Reference estimation: according to 3 independent visual amateur observations
during the night.

\solhead{5.47}{Telescope task 5}
% source id: 2019_OBS-TEL_5
Brightness estimation of STAR 2 ($\zeta$ UMi): \textbf{4.3} mag

Brightness estimation of STAR 1 ($\gamma$ UMi): \textbf{3.0} mag

\begin{center}
\begin{tabular}{ll}
Values within $\pm 0.2$ mag: & 2--2 points \\
Values within $\pm 0.3$ mag: & 1--1 point \\
\end{tabular}
\end{center}

Angular distance estimation between $\gamma$ UMi and Polaris:
$\mathbf{19^\circ}$

Values within $\pm 3^\circ$: \hfill 1 point

\solhead{5.48}{Night observation O1}
% source id: 2022_OBS-TEL_O1
Each correct answer is worth 3 pts.

\begin{center}
\begin{tabular}{lc|lc}
Object & Number & Object & Number \\
\hline
Jupiter & --- & $\beta$ Her / Antilicus & --- \\
% sic: the official solution table spells the name "Antilicus"; the question
% paper has "Antilicus" as well (the star is usually called Kornephoros)
Saturn & --- & $\alpha$ Oph / Rasalhague & --- \\
Mars & --- & $\epsilon$ Peg / Enif & 3 \\
$\alpha$ And / Alpheratz & --- & $\alpha$ Per / Mirfak & 1 \\
$\alpha$ Aql / Altair & --- & $\alpha$ Sco / Antares & 2 \\
$\alpha$ Boo / Arcturus & 4 & $\alpha$ Ser / Unukalhai & --- \\
$\alpha$ CrB / Alphecca & --- & $\epsilon$ Sgr / Kaus Australis & --- \\
$\beta$ Dra / Rastaban & --- & $\beta$ UMi / Kochab & 5 \\
\end{tabular}
\end{center}

\solhead{5.49}{Night observation O2}
% source id: 2022_OBS-TEL_O2
One method is to aim the telescope at an object, with easily recognizable
markings (such as a building) and scroll through several ($n$) fields of
view using only a single axis, while recording the difference between angles
($\Delta\alpha$) marked on the axis.

The final formula is:
\[
FOV = \frac{\Delta\alpha}{n}
\]

Full markings will be given for the range values of
$FOV = 0.6-1.2^\circ$, but only if accompanied with a working method.

\solhead{5.50}{Night observation O3}
% source id: 2022_OBS-TEL_O3
\[
F = f \cdot \frac{AFOV}{FOV}
\]
Full markings will be given for the range $f = \textbf{940--1875 mm}$, but
only if calculated using a correct formula and a value of $FOV$ from the
previous question.
% sic: the official solution writes "f" for the telescope focal length in
% the quoted range although the formula calls it F

\solhead{5.51}{Night observation O4}
% source id: 2022_OBS-TEL_O4
One method is to follow the line with a telescope to the end of the building
and remember the point. After that students could use rulers to measure $x$
and $y$ coordinate differences at an arm's length and calculate the angle.

Full markings will be given for the range $\textbf{39.5--47.5}^\circ$.

\solhead{5.52}{Observation round: the alien saucer screens}
% source id: 2023_OBS-TEL_1
\textbf{Observation 1: `Asteroid occultation'.}

Beginning of occultation: 23:03:31\\
End of occultation: 23:03:48

\begin{center}
\begin{tabular}{c|c|c|c}
Mid-time of occultation & $\pm$ error & Duration of occultation & $\pm$ error \\
\hline
23:03:39.5 & 0.4 s & 17 s & 0.2 s \\
\end{tabular}
\end{center}

Marking:
\begin{enumerate}
\item Mid-time $t$ of occultation:
  \begin{itemize}
  \item $23{:}03{:}39.3 \le t < 23{:}03{:}39.7$ $\Longrightarrow$ 4 points
  \item $23{:}03{:}39.0 \le t < 23{:}03{:}39.3$ or
    $23{:}03{:}39.7 < t \le 23{:}03{:}40.0$ $\Longrightarrow$ 3 points
  \item $23{:}03{:}38.0 \le t < 23{:}03{:}39.0$ or
    $23{:}03{:}40.0 < t \le 23{:}03{:}41.0$ $\Longrightarrow$ 1 point
  \item outside this range $\Longrightarrow$ 0 points
  \end{itemize}
\item Error in mid-time $\Delta t$:
  \begin{itemize}
  \item $0.1 < \Delta x \le 0.5$ $\Longrightarrow$ 3 points
  \item $0.5 < \Delta x \le 0.7$ $\Longrightarrow$ 2 points
  \item $0.7 < \Delta x \le 1.0$ $\Longrightarrow$ 1 point
  \item outside this range or missing $\Longrightarrow$ 0 points
  \end{itemize}
  % sic: the official solution uses "Delta x" in the bands for the error in
  % mid-time although the heading defines it as Delta t
\item Duration $x$ of occultation (should be to 1 s.f.):
  \begin{itemize}
  \item $16.8~\mathrm{s} \le x \le 17.2~\mathrm{s}$ $\Longrightarrow$ 4 points
  \item $16~\mathrm{s} \le x < 16.8~\mathrm{s}$ or
    $17.2~\mathrm{s} < x \le 18~\mathrm{s}$ $\Longrightarrow$ 3 points
  \item $15~\mathrm{s} \le x < 16~\mathrm{s}$ or
    $18~\mathrm{s} < x \le 19~\mathrm{s}$ $\Longrightarrow$ 1 point
  \item outside this range $\Longrightarrow$ 0 points
  \end{itemize}
\item Error in duration $\Delta x$:
  \begin{itemize}
  \item $0.05 < \Delta x \le 0.2$ $\Longrightarrow$ 3 points
  \item $0.2 < \Delta x \le 0.4$ $\Longrightarrow$ 2 points
  \item $0.4 < \Delta x \le 1.0$ $\Longrightarrow$ 1 point
  \item outside this range or missing $\Longrightarrow$ 0 points
  \end{itemize}
\item Error of duration lower than and different from error in mid time
  ($\Delta x < \Delta t$) $\Longrightarrow$ 1 point
\end{enumerate}

\medskip
\textbf{Observation 2: `Starlink'.}

The satellites are at an altitude $h = 20^\circ$ and their height above the
surface of the Earth is $H = 400$ km.

1) In the (simplified) solution neglecting the Earth's curvature, the
student will assume the distance to the satellites is equal to:
\[
d_{flat} = \frac{H}{\sin h} = 1170~\mathrm{km}
\]

\ioaafig{2023_OBS-TEL_1-s4.pdf}

2) In the solution taking into account the Earth's curvature, let us draw
the $OSC$ (observer--satellite--center of the Earth) triangle. The angles
are: $\angle COS = 90^\circ + h$, $\angle OSC$ that we shall denote as
$\eta$, and $\angle SCO = 180^\circ - (90^\circ + h) - \eta = 70^\circ - \eta$.

\ioaafig{2023_OBS-TEL_1-s5.pdf}

The law of sines can be applied:
\[
\frac{\sin(90^\circ + h)}{R_\oplus + H} = \frac{\sin\eta}{R_\oplus}
= \frac{\sin(70^\circ - \eta)}{d_{curve}}
\]
We calculate the distance to the satellites:
\[
\eta = \arcsin\left(\frac{R_\oplus}{R_\oplus + H}\sin(90^\circ + h)\right)
= 62^\circ.2
\]
\[
d_{curve} = R_\oplus \frac{\sin(70^\circ - \eta)}{\sin\eta} = 980~\mathrm{km}
\]

Knowing the radius of the orbit is $R_\oplus + H$, we calculate
$v_{orbit} = \sqrt{GM_\oplus/(R_\oplus + H)} = 7.6$ km/s. The observer can
only measure the tangential component of motion $v_t$.

The student will estimate the angular velocity to be
$v_t/d = v_{orbit}\sin h/d = 2.2\cdot 10^{-3}\,[1/\mathrm{s}] = 7.6\,['/\mathrm{s}]$
(for the solution neglecting the Earth's curvature) or
$2.6\cdot 10^{-3}\,[1/\mathrm{s}] = 9.0\,['/\mathrm{s}]$ (for the solution
taking into account the Earth's curvature).

\begin{flushright}
\textbf{5 points for predicting the angular velocity:}\\
5 points if within $8.55 - 9.45\,['/\mathrm{s}]$;\\
4 points if within $8.1 - 9.9\,['/\mathrm{s}]$;\\
3 points if within $7.5 - 10.5\,['/\mathrm{s}]$;\\
0 points otherwise
\end{flushright}

The simulated angular velocity is $8.5\,['/\mathrm{s}]$; the student should
conclude this value is similar to their prediction (or optionally comment it
is slightly lower and might imply e.g.\ the real orbit radius is minimally
larger than predicted, but it is not necessary).

\begin{flushright}
\textbf{4 points for measuring the angular velocity:}\\
4 points if agrees within 5\% ($8.1$ - $8.9$ $['/\mathrm{s}]$);\\
2 points if agrees within 10\% ($7.65$ - $9.35$ $['/\mathrm{s}]$);\\
0 points otherwise\\[4pt]
\textbf{1 point for correct comparison}
\end{flushright}

The student should then measure the time interval between the passes to be
$t = 2$ s.

\begin{flushright}
\textbf{2 points for measuring the time interval:}\\
2 points if agrees within 5\% ($1.9$ - $2.1$ s); 0 points otherwise
\end{flushright}

In such a short time, the path along the orbit can be approximated with a
straight line; the distance passed is $S = v_{orbit}t = 15$ km, which is the
distance between the satellites. (Alternatively, if the student chooses to
use their own $v_t$ measurement, $S = v_t t/\sin h = 14$ km).

\begin{flushright}
\textbf{2 points for calculating the distance}\\
2 points if within 13.8-15.2 km\\
1 point if within 12.5-16.5 km\\
0 points otherwise
\end{flushright}

Finally, the student should mark the observed satellite path on the sky map.

\begin{flushright}
\textbf{1 point for correctly marking the satellite trail}
\end{flushright}

\medskip
\textbf{Observation 3: `Planetary Moons'.}

The displayed planet is Saturn. On the reference solution image the moons
are marked and labelled with their names and magnitudes: Japetus (11.0),
Rhea (9.9), Mimas (13.0), Tethys (10.4), Enceladus (11.8), Dione (10.6) and
Titan (8.5) around the disk of Saturn.

\ioaafig{2023_OBS-TEL_1-s6.pdf}

2 points for each moon: position (1 pt) and name (1 pt).

\medskip
\textbf{Observation 4: `Supernova'.}

\begin{center}
\begin{tabular}{c|c|c}
Right ascension & Declination & est.\ magnitude \\
\hline
$9^h\,55^m\,54^s$ & $+69^\circ\,09'\,11''$ & 10.2 mag \\
\end{tabular}
\end{center}

\begin{enumerate}
\item Declination:
  \begin{itemize}
  \item[(a)] $< \pm 1.5'$ $\Longrightarrow$ 4 points
  \item[(b)] $\le \pm 3'$ $\Longrightarrow$ 2 points
  \item[(c)] $> \pm 3'$ $\Longrightarrow$ 0 points
  \end{itemize}
\item Right ascension:
  \begin{itemize}
  \item[(a)] $< \pm\,22.5$ s $\Longrightarrow$ 4 points
  \item[(b)] $\le \pm\,45$ s $\Longrightarrow$ 2 points
  \item[(c)] $> \pm\,45$ s $\Longrightarrow$ 0 points
  \end{itemize}
\item Magnitude:
  \begin{itemize}
  \item[(a)] $< \pm\,0.4$ mag $\Longrightarrow$ 2 points
  \item[(b)] $\le \pm\,0.8$ mag $\Longrightarrow$ 1 point
  \item[(c)] $> \pm\,0.8$ mag $\Longrightarrow$ 0 points
  \end{itemize}
\end{enumerate}

\solhead{5.53}{Discovery of 'Z'}
% source id: 2025_OBS-TEL_OT01
\textbf{(OT01.1)} The four stars are circled and labelled on the sky map:
S3 ($\beta$ Peg, $23^h3^m47^s$, $28^\circ4'58''$) at top, S1 ($\alpha$ And,
$0^h8^m24^s$, $29^\circ5'16''$) to its lower left on the map, S2
($\alpha$ Peg, $23^h4^m46^s$, $15^\circ12'17''$) at the centre and S4
($\gamma$ Peg, $0^h13^m14^s$, $15^\circ10'59''$) below.

\ioaafig{2025_OBS-TEL_OT01-s2.pdf}

\begin{center}
\begin{tabular}{l r}
\hline
Circle around each correct star & 0.5 \\
Correct identification as S1--S4 of each star & 1.0 \\
\hline
\end{tabular}
\end{center}

\textbf{(OT01.2)}
\begin{itemize}
\item Realize that two stars have almost the same RA and two of them have
  almost same declination.
\item We draw a line through them indicating constant RA and constant
  Declination. (On the official overlay the construction lines are labelled
  $\delta \approx 15^\circ$, $\alpha \approx 23^h$, and their parallels
  through Z are labelled $\alpha \approx 21^h36^m$ and
  $\delta \approx -26^\circ$.)
\item Given the scale of the map and the separation between the 4 stars and
  the new object we can assume that the celestial grid will be almost
  linear.
\item We can thus mark `Z' at its given coordinates
  ($\alpha = 21^h36^m$; $\delta = -26^\circ10'$).
\end{itemize}

\begin{center}
\begin{tabular}{l r}
\hline
Object marked within inner rectangle & 7.0 \\
\quad Object marked between outer and inner rectangles & 4.0 \\
\hline
\end{tabular}
\end{center}

\textbf{(OT01.3)}
\begin{itemize}
\item After marking the new object on the given sky map, we carefully look
  at its position and the star pattern surrounding it.
\item Using the 25 mm focal length eyepiece we try to locate it.
\item By doing some fine adjustments, we try to bring the same pattern in
  the 10 mm focal length eye piece (which has a crosshair).
\end{itemize}

\begin{center}
\begin{tabular}{l r}
\hline
Pointing towards the screen & 1.0 \\
\hline
\multicolumn{2}{l}{If 10 mm objective is used:} \\
% sic: the official marking table says "objective" for the eyepiece
\quad Correct choice of eyepiece & 2.0 \\
\quad Object is anywhere in FOV & 3.0 \\
\quad Any part of the object is touching the origin of crosshair & 6.0 \\
\hline
\multicolumn{2}{l}{If 25 mm objective is used:} \\
\quad Object is anywhere in FOV & 3.0 \\
\hline
\end{tabular}
\end{center}

\solhead{5.54}{Lensing Time Delay}
% source id: 2025_OBS-TEL_OT02
The solution is $t_{\mathrm{DA}} = 1694$ d, $t_{\mathrm{DB}} = 1651$ d and
$t_{\mathrm{DC}} = 2485$ d.

\begin{center}
\begin{tabular}{l r}
\hline
Pointing towards the screen & 1.0 \\
Choice of correct eyepiece & 1.0 \\
\hline
\multicolumn{2}{l}{For each of images A, B, C, D:} \\
\quad Timings of image & 2.0 \\
\qquad At least one reading & 1.0 \\
\qquad Extra reading (0.5 each, maximum 1.0) & 1.0 \\
\hline
\multicolumn{2}{l}{Calculation of time differences:} \\
\quad Calculation each of $t_{\mathrm{DA}}$, $t_{\mathrm{DB}}$ and
  $t_{\mathrm{DC}}$ & 3.0 \\
\qquad full credit for value $-25$ d to value $+25$ d; & \\
\qquad half credit down to value $-50$ d or up to value $+50$ d & \\
\quad Scaling from seconds to days for each interval & 1.0 \\
\hline
Correct order ($t_{\mathrm{DC}} > t_{\mathrm{DA}} > t_{\mathrm{DB}}$) & 3.0 \\
\hline
\end{tabular}
\end{center}
